% 自 Notizen/W02_0921/ 各表由 Skripte/gen_tex.py 產生；改字改表後重產，不手改本檔。
\chapter{Tremblay 論違憲審查的一般正當性（2026-09-21）}

2026 年 9 月 21 日上午第 3、4 節課堂錄音（約 122 分鐘）之逐字稿。左邊界為段首時間碼；未標說話人者為湯德宗老師。同週的課堂整理見課程問答集第 2 部。


\section{篇名的逐字解讀}

\tc{0:00:00}
% cue 1 0.930 15.500
\hypertarget{w02:cue:1}{}所以以前也不認識加拿大人\footnote{Luc B. Tremblay，加拿大法學者，本文刊出時任教於蒙特婁大學（Université de Montréal）法學院。另著有 The Rule of Law, Justice, and Interpretation（Montreal: McGill-Queen's University Press, 1997），本文第 528 頁註 7 引該書第 29 至 34 頁說明法治概念。}。可以寫這麼難的文章。好，那他寫的這個文章的深度。我不知道你們可能明天或是哪天要去參加 Bruce Ackerman\footnote{美國憲法學者，1943 年生，耶魯大學 Sterling 法學與政治學講座教授，著有 We the People 三卷（1991、1998、2014）。2026 年 6 月 18 日獲唐獎法治獎；同年 9 月 23 日出席在國立臺灣大學舉行的第七屆唐獎大師論壇閉幕場。\url{https://law.yale.edu/bruce-ackerman}。} 的演講的話，
% cue 2 15.500 26.840
\hypertarget{w02:cue:2}{}那就比那個 Bruce Ackerman 講的像小說一樣的東西，那是難多的嘛。你們聽聽看這個人。他的標題的 General Legitimacy\footnote{Luc B. Tremblay, General Legitimacy of Judicial Review and the Fundamental Basis of Constitutional Law，刊於 Oxford Journal of Legal Studies 第 23 卷第 4 期，頁 525–562（2003）。課綱第 2 週（9 月 21 日）指定文獻，同週另指定 Tamas Gyorfi 的 The Debate about the Justification of Constitutional Review: Six Theses。}，就是一般正當性。

\tc{0:00:26}
% cue 3 26.840 58.270
\hypertarget{w02:cue:3}{}你們就覺得很奇怪。正當性就正當性，還要一般正當性存在心理。他等一下會講，跟一般正當性相對比的。它叫做 specific legitimacy。我把它翻成叫特定正當性。但沒關係啦。一般正當性，什麼東西的一般正當性？of judicial review。所以這是主標題。換句話就是說，違憲審查的一般正當性。下面 and the fundamental basis of constitutional law.這句有點難了。憲法就憲法。憲法還有
% cue 4 58.270 94.230
\hypertarget{w02:cue:4}{}沒有一個根本基礎。那這話中就有話了。所以他把這兩件事情擺在一起。你們看到這個標題，你就要去想。我們講違憲審查的正當性。為什麼會扯到憲法的根本基礎。他講基礎，還要加個 fundamental。所以它不只是 basis，它是 fundamental basis。基礎的基礎。那這兩件事情是什麼關係？簡單的講，他認為要討論正當性的問題，不能不討論憲法的根本基礎。簡單講是這樣。

\section{摘要所列的四個問題與本文目的}

\tc{0:01:34}
% cue 5 94.230 127.750
\hypertarget{w02:cue:5}{}好，我們看一下 abstract,4 questions dominate normative contemporary constitutional theory，我們現代的所謂的憲法理論，這個字圈起來，constitutional theory，這個以後只要你看美式英文，現在加拿大也是，他們一定會用這個字，所謂的憲法理論，憲法理論比我們想像中複雜，他現在又加了一個字，normative，規範性的，規範性的當代憲法理論，換句話說，憲法理論有些可能不是規範，他可能是 emma。

\tc{0:02:07}
% cue 6 127.750 159.600
\hypertarget{w02:cue:6}{}有些可能是別的，好，那裡面呢，主要是四個問題，四個問題主宰了這個，那也就是說這個包含了四個問題，討論主要是四個問題，What is the purpose of the constitution？憲法的目的是什麼？這個事情很複雜啊，那你覺得憲法的目的會影響到憲法理論嗎？What makes the constitution legitimate？什麼東西使得憲法變成正當？這個問題在我們這裡幾乎不討論，憲法會不正當嗎？
% cue 7 159.600 164.220
\hypertarget{w02:cue:7}{}你還問說什麼東西使憲法變成正當，那你意思是有不正當的憲法嗎？

\tc{0:02:44}
% cue 8 164.220 170.630
\hypertarget{w02:cue:8}{}What kind of argument are legitimate within the process of constitutional interpretation?
% cue 9 170.630 191.430
\hypertarget{w02:cue:9}{}這個跟我們上去一起講的比較有關係，在我們解釋憲法的過程中，那一些個論辯是正當的，換句話說，有一些憲法上的解釋，你可能聽到一些不正當的論辯，它沒有辦法幫助你這個解釋變成正當。
% cue 10 191.430 207.210
\hypertarget{w02:cue:10}{}然後第四個問題，what can make judicial review of the legislation legitimate in principle？什麼東西是可以使得對於立法進行違憲審查在原則上變得正當？
% cue 11 207.440 209.230
\hypertarget{w02:cue:11}{}這就是我們本文要講的重點。

\tc{0:03:29}
% cue 12 209.230 216.230
\hypertarget{w02:cue:12}{}The main purpose of this text is to provide one general answer to the last question.
% cue 13 216.230 222.590
\hypertarget{w02:cue:13}{}本文主要是針對第四個問題提供一個一般性的解答，不是特定的一般。
% cue 14 222.920 229.160
\hypertarget{w02:cue:14}{}換句話說，是普遍適用的，我講正當就是哪裡都正當，不是只有在特定情況下正當的。

\tc{0:03:49}
% cue 15 229.160 247.460
\hypertarget{w02:cue:15}{}The secondary purpose，這是本文主要目的的 main purpose，但是有 secondary purpose，次要的目的，是 to show，來顯示，how this answer，我提出來的這個違憲審查，一般正當性理論，may bear upon，會跟什麼東西發生關聯，
% cue 16 247.460 258.510
\hypertarget{w02:cue:16}{}會跟我們對於憲法的根本基礎的理解產生關聯。
% cue 17 258.510 265.080
\hypertarget{w02:cue:17}{}These two purposes should suggest particular answers to the first three questions.
% cue 18 265.080 277.580
\hypertarget{w02:cue:18}{}這兩個目的可以幫助我們提供對於前三個問題的答案，是以第四個問題的答案會提供前三個問題，也能提供一定的答案。

\section{結論的一般正當性命題與憲法命題}

\tc{0:04:37}
% cue 19 277.580 297.580
\hypertarget{w02:cue:19}{}好，這是他的 extract，那我們先看一下，我怕等一下時間來不及，我們先看一下他的結論，因為他的結論很短，在 561 頁，你看現在你們如果沒有紙本，你看你是不是跟著我累死了，你拉到最底下告回來，他的結論很短。
% cue 20 297.580 314.730
\hypertarget{w02:cue:20}{}第四，conclusion,the main purpose of this text was to provide a general answer to one of the questions that have dominated contemporary and normative constitutional theory,what can make a judicial review,morally legitimate in principle?

\tc{0:05:14}
% cue 21 314.730 320.190
\hypertarget{w02:cue:21}{}好，在這邊要注意的是，你看他的文句，那邊重複了多份我就不講了。
% cue 22 320.190 331.990
\hypertarget{w02:cue:22}{}本文主要的目的就是對於主宰現代所謂的規範性的憲法理論裡面的四個問題提供一個答案。
% cue 23 331.990 341.960
\hypertarget{w02:cue:23}{}那什麼呢？就是什麼東西可以使得違憲審查，在原則上具有道德的正當性？
% cue 24 341.960 347.070
\hypertarget{w02:cue:24}{}morally，剛才我們前面只講了 legitimate，他現在把 morally 加進去。
% cue 25 347.070 353.860
\hypertarget{w02:cue:25}{}所以他的意思是，當我討論正不正當，我主要是從道德的觀點來看，正當與否。
% cue 26 353.860 360.390
\hypertarget{w02:cue:26}{}我們中文要講這麼多句，你用直翻的話，就是道德上的正當性。
% cue 27 360.390 370.050
\hypertarget{w02:cue:27}{}但是這樣子的句子在中文裡面是不好的，嚴格來講是從道德的觀點來看是否正當，原則上是正當的。
% cue 28 370.050 374.330
\hypertarget{w02:cue:28}{}為什麼違憲審查從道德的觀點來看，原則上是正當的？
% cue 29 374.330 377.100
\hypertarget{w02:cue:29}{}這個文章主要在講這件事情。

\tc{0:06:17}
% cue 30 377.100 387.720
\hypertarget{w02:cue:30}{}我的答案是，我的答案是，法律的違憲審查是正當的。
% cue 31 387.720 393.530
\hypertarget{w02:cue:31}{}由一種普通的政治性的理性，
% cue 32 393.530 406.560
\hypertarget{w02:cue:32}{}這句話真的太難了。
% cue 33 406.560 413.980
\hypertarget{w02:cue:33}{}他說，為什麼違憲審查就是對立法進行是否合憲的這個審查是正當的？
% cue 34 413.980 429.790
\hypertarget{w02:cue:34}{}這是由於有一個政治，這個政治道德的原則是說，所有的政治權利在一個國家之內所有政治權利。
% cue 35 429.790 436.000
\hypertarget{w02:cue:35}{}因為我們白話講就是，憲法所規定的所有憲法機關行使公權力的時候，
% cue 36 436.000 444.690
\hypertarget{w02:cue:36}{}那包含司法行使公權力，違憲審查的時候，
% cue 37 444.690 452.970
\hypertarget{w02:cue:37}{}我們的以一種合法的方式來行使，那這個就是一個基本原則。
% cue 38 452.970 460.480
\hypertarget{w02:cue:38}{}按照這個基本原則來看，我認為，違憲審查原則上是正當的。
% cue 39 460.480 463.590
\hypertarget{w02:cue:39}{}這是他的答案。
% cue 40 463.590 470.540
\hypertarget{w02:cue:40}{}我把這個答案呢，就叫做一般正當性命題，
% cue 41 470.540 472.720
\hypertarget{w02:cue:41}{}等一下它會有講，一般正當性命題，
% cue 42 472.720 474.680
\hypertarget{w02:cue:42}{}我覺得是五個，六個，複雜的槍字題。

\tc{0:07:54}
% cue 43 474.680 479.930
\hypertarget{w02:cue:43}{}第二個目的，這是主要目的，是主要目的，
% cue 44 479.930 486.080
\hypertarget{w02:cue:44}{}是主要目的，是主要目的，是要從這個命題點生出來一些後果。
% cue 45 486.080 494.180
\hypertarget{w02:cue:45}{}主要的是，牽涉到我們對於憲法的根本基礎的理解。
% cue 46 494.180 508.230
\hypertarget{w02:cue:46}{}我認為，這個就是他剛剛講的所謂的 fundamental basis。
% cue 47 508.230 516.800
\hypertarget{w02:cue:47}{}我主張，有關於憲法的基本的規則和原則。
% cue 48 516.800 523.740
\hypertarget{w02:cue:48}{}規則就比方說，過馬路闖紅燈則死，那你有沒有闖紅燈，有，那就是死，
% cue 49 523.740 530.660
\hypertarget{w02:cue:49}{}規則，一番兩燈，有，原則，比方說，有了自由，英語，保障，
% cue 50 530.660 534.220
\hypertarget{w02:cue:50}{}那你會討論有時是原則，但有些例外，等等。
% cue 51 534.220 537.140
\hypertarget{w02:cue:51}{}所以它是有彈性的，不管啦，反正就是，
% cue 52 537.140 541.580
\hypertarget{w02:cue:52}{}憲法上的規則或原則，dealing with the validity，
% cue 53 541.580 548.750
\hypertarget{w02:cue:53}{}牽涉到所謂的有效性，就有關於有效性的憲法，
% cue 54 548.750 554.910
\hypertarget{w02:cue:54}{}原則和法則，and interpretation of legislation，
% cue 55 554.910 557.140
\hypertarget{w02:cue:55}{}以及立法的解釋。

\tc{0:09:17}
% cue 56 557.140 563.550
\hypertarget{w02:cue:56}{}Must be conceived as deriving from a set of antecedent norms
% cue 57 563.550 573.210
\hypertarget{w02:cue:57}{}that provide conditions,legislation must satisfy,in order to be accepted as morally legitimate.
% cue 58 573.210 580.380
\hypertarget{w02:cue:58}{}為了要使得立法能夠被認為，被接受為道德上正當，
% cue 59 580.380 589.830
\hypertarget{w02:cue:59}{}所以它的正當是道德上正當，那它必須要符合某一些先驗的，更前，
% cue 60 589.830 596.880
\hypertarget{w02:cue:60}{}更先行的規範，這樣子你才能夠說這個法律是正當的，and consequently
% cue 61 596.880 603.830
\hypertarget{w02:cue:61}{}as morally entitled to judicial enforcement，
% cue 62 603.830 613.200
\hypertarget{w02:cue:62}{}所以說，那麼法律就有權利獲得司法的執行，那就是我們法院要 honor 這個規定。

\tc{0:10:13}
% cue 63 613.200 622.190
\hypertarget{w02:cue:63}{}如果你不符合到我們法院就會說，哦，你違憲，那所以我就不執行，這個跟今天我們通告的某個問題很像，
% cue 64 622.190 635.560
\hypertarget{w02:cue:64}{}卓榮泰\footnote{卓榮泰於 2024 年 5 月 20 日就任行政院院長。2025 年 12 月 15 日宣布不副署立法院同年 11 月三讀通過的財政收支劃分法修正案；憲法第 37 條規定，總統依法公布法律，須經行政院院長之副署，或行政院院長及有關部會首長之副署。\url{https://www.ey.gov.tw/Page/3C927BC7FA83C969}。}說他不執行法律，那我們現在違憲審查是不是司法機關不執行法律，有個法律通過了，但是法律說我不執行，我不執行的理由，因為他不正當，他為什麼不正當。

\tc{0:10:35}
% cue 65 635.560 637.560
\hypertarget{w02:cue:65}{}因為他違反了那個法律必須要成為正當所需要符合的若干先行規範，所以那個 legitimacy 跟那個先行規範是有關係的，幾乎劃上等號，如果你違反了法律應該要具備正當性的那些先行規範，那你違反了那些先行規範，那你就不具有正當性，你不具有正當性，那我就有正當理由不去執行，司法不執行法律，I have court argument that it is, I have courted argument that it is, that it is not a law。

\tc{0:10:37}
% cue 66 637.560 669.340
\hypertarget{w02:cue:66}{}所以，你違反了，你違反了，那個法律必須要成為正當所需要符合的若干先行規範，所以那個 legitimacy 跟那個先行規範是有關係的，幾乎劃上等號，如果你違反了法律應該要具備正當性的那些先行規範，那你違反了那些先行規範，那你就不具有正當性，你不具有正當性，那我就有正當理由不去執行，司法不執行法律，I have called the argument that supports this claim。

\tc{0:11:09}
% cue 67 669.340 688.100
\hypertarget{w02:cue:67}{}the Constitutional Law Thesis，我把它叫做憲法的正多複雜，一般性核線，一般核線性命題，又有憲法命題，等一下你還有看到特殊正當性的命題，再往下一段，
% cue 68 688.100 702.310
\hypertarget{w02:cue:68}{}the general legitimacy thesis and the Constitutional Law Thesis therefore have some bearing on the three other questions that are in the core of contemporary normative constitutional theory，
% cue 69 702.310 716.170
\hypertarget{w02:cue:69}{}當代的規範性的憲法理論的核心四個問題，這個剛才講的兩個命題，一個是一般正當性命題，一個是憲法命題，跟他們都有關係，怎麼有關係呢？

\tc{0:11:56}
% cue 70 716.170 722.680
\hypertarget{w02:cue:70}{}first, they entail the main purpose or function of a constitution，
% cue 71 722.680 729.750
\hypertarget{w02:cue:71}{}憲法的目的，或者是憲法的功能，insofar as it deals with legislation，
% cue 72 729.750 733.750
\hypertarget{w02:cue:72}{}當它跟立法有關的時候，is to make legislation legitimate，
% cue 73 733.750 742.820
\hypertarget{w02:cue:73}{}憲法的功能就是使得法律變得正當，for the purpose of judicial enforcement of law，
% cue 74 742.820 749.340
\hypertarget{w02:cue:74}{}為了便於司法來執行那個法律，這個話有點複雜對不對？
% cue 75 749.340 756.290
\hypertarget{w02:cue:75}{}簡單在講究說，他認為第一點，憲法的功能就是要幫立法變得正當，
% cue 76 756.290 764.560
\hypertarget{w02:cue:76}{}以便於司法的執行，這是他剛剛講的第一個問題，憲法的目的，
% cue 77 764.560 769.880
\hypertarget{w02:cue:77}{}憲法的目的，從他的觀點看，就是要幫助立法變得正當。

\tc{0:12:49}
% cue 78 769.880 776.360
\hypertarget{w02:cue:78}{}A constitution must be understood，憲法應該被理解為，必須被理解為，
% cue 79 776.360 784.160
\hypertarget{w02:cue:79}{}Expressing or embodying a set of rules or principles，是包含，或者它是表述，
% cue 80 784.160 789.780
\hypertarget{w02:cue:80}{}一系列的規則和原則，providing the criteria of validity，
% cue 81 789.780 796.570
\hypertarget{w02:cue:81}{}對於什麼叫做有效性，它提供了一個標準，有效性標準，
% cue 82 796.570 802.410
\hypertarget{w02:cue:82}{}提供有效性標準的這一系列的規則和原則，and，
% cue 83 802.410 812.670
\hypertarget{w02:cue:83}{}以及 of proper interpretation of legislation，以及有關於法律適當解釋的這一系列的規則和原則，
% cue 84 812.670 824.500
\hypertarget{w02:cue:84}{}that make it morally entitled to judicial enforcement，使得它可以獲得有權利，獲得司法的執行，也就是它是和諧的。

\tc{0:13:44}
% cue 85 824.500 833.480
\hypertarget{w02:cue:85}{}The basic purpose and function of a constitutional rules and principles is to enable the making of a legitimate law，
% cue 86 833.480 844.290
\hypertarget{w02:cue:86}{}所以它才講，我們從立法的觀點來看，我們之所以要憲法，憲法的功能，或者憲法的目的，就是要使者提供一些標準，
% cue 87 844.290 851.110
\hypertarget{w02:cue:87}{}如何使得立法能夠變成正當的上次第一點。

\tc{0:14:11}
% cue 88 851.110 855.650
\hypertarget{w02:cue:88}{}第二點，the legitimacy of a constitution，
% cue 89 855.650 857.740
\hypertarget{w02:cue:89}{}憲法的正當性，
% cue 90 857.740 867.820
\hypertarget{w02:cue:90}{}憲法的正當性有部分是來自於下列的事實，它是有助於創造維持促進一個政府體系，
% cue 91 867.820 869.820
\hypertarget{w02:cue:91}{}那個政府的體系可以產生正當的立法，這個句子太複雜，而且有點頭跟尾巴，
% cue 92 869.820 874.580
\hypertarget{w02:cue:92}{}而且像頭搖，頭搖，尾巴要頭。
% cue 93 874.580 877.060
\hypertarget{w02:cue:93}{}the basic ground of the normative force of the constitutional rules and the principles is not source-based，
% cue 94 877.060 879.060
\hypertarget{w02:cue:94}{}它們不是所謂的看什麼樣的基礎，看什麼樣的來源，不是以法源來論斷，
% cue 95 879.060 881.570
\hypertarget{w02:cue:95}{}不是以法源來論斷。
% cue 96 881.570 883.980
\hypertarget{w02:cue:96}{}it is not a matter of will or consent，
% cue 97 883.980 885.970
\hypertarget{w02:cue:97}{}也不牽涉到人民的意志或論斷，
% cue 98 885.970 889.150
\hypertarget{w02:cue:98}{}也不牽涉到人民的意志或論斷。
% cue 99 889.150 891.490
\hypertarget{w02:cue:99}{}而是以法源來論斷，
% cue 100 891.490 893.830 刪：與 cue 99 同文（連續 4 個重複解碼）
% cue 101 893.830 896.110 刪：與 cue 99 同文（連續 4 個重複解碼）
% cue 102 896.110 898.070 刪：與 cue 99 同文（連續 4 個重複解碼）
% cue 103 902.190 911.330 刪：與 cue 94 同文（連續 4 個重複解碼）

\tc{0:15:11}
% cue 104 911.330 914.140
\hypertarget{w02:cue:104}{}It is not a matter of will or consent，
% cue 105 914.140 917.620
\hypertarget{w02:cue:105}{}也不牽涉到人民的意志或擁有同意的問題。
% cue 106 918.250 920.250
\hypertarget{w02:cue:106}{}at least not in an exclusionary sense，
% cue 107 920.250 922.250
\hypertarget{w02:cue:107}{}至少不是絕對如此。

\tc{0:15:22}
% cue 108 922.250 924.310
\hypertarget{w02:cue:108}{}那不是這樣是什麼呢？
% cue 109 924.310 928.500
\hypertarget{w02:cue:109}{}這些所謂憲法的規則和法則的，
% cue 110 928.500 933.000
\hypertarget{w02:cue:110}{}規範性的基本效力是怎麼來的呢？
% cue 111 933.000 936.980
\hypertarget{w02:cue:111}{}It is primarily a matter of content，
% cue 112 936.980 939.120
\hypertarget{w02:cue:112}{}它主要是牽涉到內容，
% cue 113 939.120 943.170
\hypertarget{w02:cue:113}{}It is a matter of being in accordance with a set of moral
% cue 114 943.170 945.170
\hypertarget{w02:cue:114}{}antecedent norms
% cue 115 945.170 953.170
\hypertarget{w02:cue:115}{}the substance of which determines the conditions of the legitimate legislation。

\tc{0:15:53}
% cue 116 953.170 957.090
\hypertarget{w02:cue:116}{}這句話就是剛才講的那個，
% cue 117 957.090 959.810
\hypertarget{w02:cue:117}{}它說這是牽涉到內容，
% cue 118 959.810 966.980
\hypertarget{w02:cue:118}{}這個內容是因為它根據一套道德的先行的規範，
% cue 119 966.980 970.160
\hypertarget{w02:cue:119}{}來決定了這個法律，
% cue 120 970.160 972.450
\hypertarget{w02:cue:120}{}這個立法，
% cue 121 972.450 974.560
\hypertarget{w02:cue:121}{}它是不是正當的那些條件，
% cue 122 974.560 977.230
\hypertarget{w02:cue:122}{}是有一套先行的規範，
% cue 123 977.230 979.510
\hypertarget{w02:cue:123}{}來決定法律正當，
% cue 124 979.510 983.450
\hypertarget{w02:cue:124}{}因為它要正當，
% cue 125 983.450 985.440
\hypertarget{w02:cue:125}{}它要符合某些條件，
% cue 126 985.440 987.760
\hypertarget{w02:cue:126}{}所以它主要牽涉到的是內容。
% cue 127 987.760 989.880
\hypertarget{w02:cue:127}{}它不是說你是根據憲法嗎？
% cue 128 989.880 991.850
\hypertarget{w02:cue:128}{}來源，法源，
% cue 129 991.850 995.880
\hypertarget{w02:cue:129}{}不是，也不跟有沒有得到人民同意有關，
% cue 130 995.880 996.880
\hypertarget{w02:cue:130}{}至少它不決然，
% cue 131 996.880 998.320
\hypertarget{w02:cue:131}{}絕對如此。
% cue 132 998.320 999.440
\hypertarget{w02:cue:132}{}這是第二點。

\tc{0:16:39}
% cue 133 999.440 1000.440
\hypertarget{w02:cue:133}{}第三點，
% cue 134 1000.440 1003.530
\hypertarget{w02:cue:134}{}Constitutional interpretation is made legitimate，
% cue 135 1003.530 1006.660
\hypertarget{w02:cue:135}{}憲法的解釋會變得正當，
% cue 136 1006.660 1008.690
\hypertarget{w02:cue:136}{}By virtue of interpretive argument，
% cue 137 1008.690 1013.750
\hypertarget{w02:cue:137}{}經由下列這些所謂的解釋性的論證，
% cue 138 1013.750 1016.100
\hypertarget{w02:cue:138}{}那個解釋性的論證，
% cue 139 1016.100 1021.340
\hypertarget{w02:cue:139}{}that contribute to make rules and principles of a constitutional law legitimate，
% cue 140 1021.340 1025.650
\hypertarget{w02:cue:140}{}使得憲法的原則，
% cue 141 1025.650 1028.190
\hypertarget{w02:cue:141}{}法則變得有助於，
% cue 142 1028.190 1032.380
\hypertarget{w02:cue:142}{}有助於使這些憲法的原則、法則變得正當。
% cue 143 1032.380 1034.690
\hypertarget{w02:cue:143}{}that is，
% cue 144 1034.690 1035.690
\hypertarget{w02:cue:144}{}也就是說，
% cue 145 1035.690 1039.740
\hypertarget{w02:cue:145}{}consistent with the antecedent norms
% cue 146 1039.740 1042.770
\hypertarget{w02:cue:146}{}with respect to the legitimacy of the legislation。

\tc{0:17:22}
% cue 147 1042.770 1044.060
\hypertarget{w02:cue:147}{}也就是說，
% cue 148 1044.060 1048.080
\hypertarget{w02:cue:148}{}這一些個解釋性的論辯，
% cue 149 1048.080 1054.870
\hypertarget{w02:cue:149}{}是符合那個使得法律法變得正當，
% cue 150 1054.870 1059.130
\hypertarget{w02:cue:150}{}因為它符合那某一套預先設立的那些規範，
% cue 151 1059.130 1061.510
\hypertarget{w02:cue:151}{}也就是你符合那些預設的規範，
% cue 152 1061.510 1064.600
\hypertarget{w02:cue:152}{}這個法律立法才會變得正當。

\tc{0:17:44}
% cue 153 1064.600 1066.760
\hypertarget{w02:cue:153}{}那我們來 argue 說，
% cue 154 1066.760 1069.760
\hypertarget{w02:cue:154}{}這個法律怎麼來解釋它，
% cue 155 1069.760 1073.760
\hypertarget{w02:cue:155}{}就是透過各種要的解釋性的論辯，
% cue 156 1073.760 1075.670
\hypertarget{w02:cue:156}{}來有助於促成，
% cue 157 1075.670 1079.860
\hypertarget{w02:cue:157}{}這些憲法最基本的這些原則、法則，
% cue 158 1079.860 1082.880
\hypertarget{w02:cue:158}{}能夠維持，使它變得正當。
% cue 159 1082.880 1086.200
\hypertarget{w02:cue:159}{}這樣的句子是粘在一起的。

\tc{0:18:06}
% cue 160 1086.200 1088.930
\hypertarget{w02:cue:160}{}legitimate interpretive argument，
% cue 161 1088.930 1092.710
\hypertarget{w02:cue:161}{}正當的解釋論證，
% cue 162 1092.710 1096.770
\hypertarget{w02:cue:162}{}must support constitutional interpretation，
% cue 163 1096.770 1098.830
\hypertarget{w02:cue:163}{}能夠支持憲法的解釋，
% cue 164 1098.830 1100.050
\hypertarget{w02:cue:164}{}那個解釋是，
% cue 165 1100.050 1105.050
\hypertarget{w02:cue:165}{}make legislation morally entitled to judicial enforcement。
% cue 166 1105.050 1108.070
\hypertarget{w02:cue:166}{}這句話真太複雜了，
% cue 167 1108.070 1110.090
\hypertarget{w02:cue:167}{}它的意思是說，
% cue 168 1110.090 1111.430
\hypertarget{w02:cue:168}{}好像是說，
% cue 169 1111.430 1112.680
\hypertarget{w02:cue:169}{}我的理解，
% cue 170 1112.680 1114.930
\hypertarget{w02:cue:170}{}如果你現在做一個解釋，
% cue 171 1114.930 1124.360
\hypertarget{w02:cue:171}{}而那個解釋無助於使得這個規定，
% cue 172 1124.360 1127.320
\hypertarget{w02:cue:172}{}能夠獲得法院的支持去執行，
% cue 173 1127.320 1130.400
\hypertarget{w02:cue:173}{}那這樣的解釋是不正當的，
% cue 174 1130.400 1131.360
\hypertarget{w02:cue:174}{}沒有用的。

\tc{0:18:51}
% cue 175 1131.360 1136.770
\hypertarget{w02:cue:175}{}this applies to written as well as to unwritten constitutional rules and principles，
% cue 176 1136.770 1139.800
\hypertarget{w02:cue:176}{}我們剛才講的憲法的原則，
% cue 177 1139.800 1142.800
\hypertarget{w02:cue:177}{}基礎的原則或者是法則，
% cue 178 1142.800 1144.820
\hypertarget{w02:cue:178}{}有可能是成文的，
% cue 179 1144.820 1145.820
\hypertarget{w02:cue:179}{}有可能是不成文的，
% cue 180 1145.820 1149.620
\hypertarget{w02:cue:180}{}散布在各種慣例習慣，
% cue 181 1149.620 1152.620
\hypertarget{w02:cue:181}{}還有法官講的法律法裡，
% cue 182 1152.620 1154.360
\hypertarget{w02:cue:182}{}都一樣。
% cue 183 1154.360 1158.420
\hypertarget{w02:cue:183}{}你都必須要是有一套先行的這個原則，
% cue 184 1158.420 1159.480
\hypertarget{w02:cue:184}{}跟法則，
% cue 185 1159.480 1161.540
\hypertarget{w02:cue:185}{}你必須要符合它，
% cue 186 1161.540 1164.540
\hypertarget{w02:cue:186}{}它是來決定法律是不是正當的，
% cue 187 1164.540 1165.540
\hypertarget{w02:cue:187}{}那法律是正當的，
% cue 188 1165.540 1168.540
\hypertarget{w02:cue:188}{}這樣子你才有辦法讓司法去執行它，
% cue 189 1168.540 1170.760
\hypertarget{w02:cue:189}{}司法執行也才是正當的。

\tc{0:19:30}
% cue 190 1170.760 1172.550
\hypertarget{w02:cue:190}{}therefore
% cue 191 1172.550 1174.550
\hypertarget{w02:cue:191}{}insofar as political legitimacy
% cue 192 1174.550 1179.570
\hypertarget{w02:cue:192}{}is conceived as coming within some version of democratic theory
% cue 193 1179.570 1185.140
\hypertarget{w02:cue:193}{}as it is generally postulated within contemporary Western society
% cue 194 1185.140 1188.330
\hypertarget{w02:cue:194}{}the purpose and function of the constitutional law
% cue 195 1188.330 1196.990
\hypertarget{w02:cue:195}{}dealing with legislation should be conceived as enabling the making of democratic legislation
% cue 196 1196.990 1199.270
\hypertarget{w02:cue:196}{}whatever that means。

\tc{0:19:59}
% cue 197 1199.270 1201.940
\hypertarget{w02:cue:197}{}那憲法的目的，
% cue 198 1201.940 1206.150
\hypertarget{w02:cue:198}{}從對付立法的觀點來說，
% cue 199 1206.150 1208.140
\hypertarget{w02:cue:199}{}應該被理解為，
% cue 200 1208.140 1213.980
\hypertarget{w02:cue:200}{}使得民主立法的過程，
% cue 201 1213.980 1217.980
\hypertarget{w02:cue:201}{}enabling the making，
% cue 202 1217.980 1224.670
\hypertarget{w02:cue:202}{}就是我們怎麼樣能夠創造出一個民主的立法，
% cue 203 1224.670 1228.180
\hypertarget{w02:cue:203}{}能夠協助創造出民主立法，
% cue 204 1228.180 1231.370
\hypertarget{w02:cue:204}{}這就是我們憲法的目的。
% cue 205 1231.370 1235.970
\hypertarget{w02:cue:205}{}如果憲法沒有辦法協助創造出民主的立法，
% cue 206 1235.970 1239.110
\hypertarget{w02:cue:206}{}管理你對這個民主立法是怎麼理解，
% cue 207 1239.110 1241.190
\hypertarget{w02:cue:207}{}如果它無助於形成民主立法，
% cue 208 1241.190 1243.400
\hypertarget{w02:cue:208}{}那這不是憲法的目的。

\tc{0:20:43}
% cue 209 1243.400 1244.840
\hypertarget{w02:cue:209}{}Furthermore
% cue 210 1244.840 1251.120
\hypertarget{w02:cue:210}{}the legitimate authority of the constitutional rules and principles dealing with the legislation
% cue 211 1251.120 1258.780
\hypertarget{w02:cue:211}{}should depend on whether their enforcement contribute to create and enhance democratic legislation，
% cue 212 1258.780 1261.130
\hypertarget{w02:cue:212}{}這又念到一起了。
% cue 213 1261.130 1264.060
\hypertarget{w02:cue:213}{}這有點雞生蛋蛋生雞了。
% cue 214 1264.060 1269.710
\hypertarget{w02:cue:214}{}而一個所謂有關於立法的，
% cue 215 1269.710 1275.100
\hypertarget{w02:cue:215}{}這個憲法法則和原則的正當性權利，
% cue 216 1275.100 1279.480
\hypertarget{w02:cue:216}{}是取決於他們的執行，
% cue 217 1279.480 1286.220
\hypertarget{w02:cue:217}{}是否有助於創造、維護、促進這個民主的立法，
% cue 218 1286.220 1289.250
\hypertarget{w02:cue:218}{}如果無助於民主的立法，
% cue 219 1289.250 1292.960
\hypertarget{w02:cue:219}{}那我們憲法的這套法則就沒有用了，
% cue 220 1292.960 1294.500
\hypertarget{w02:cue:220}{}它就不算是正當的。

\tc{0:21:34}
% cue 221 1294.500 1300.900
\hypertarget{w02:cue:221}{}同理類似的 legitimate constitutional interpretation，
% cue 222 1300.900 1306.500
\hypertarget{w02:cue:222}{}合法的正當的憲法解釋，
% cue 223 1306.500 1308.760
\hypertarget{w02:cue:223}{}should depend on，
% cue 224 1308.760 1309.820
\hypertarget{w02:cue:224}{}應該取決於，
% cue 225 1309.820 1311.880
\hypertarget{w02:cue:225}{}whether the norms it entails，
% cue 226 1311.880 1314.040
\hypertarget{w02:cue:226}{}它所產生，
% cue 227 1314.040 1318.040
\hypertarget{w02:cue:227}{}因過經過憲法解釋所產生出來的這個規範，
% cue 228 1318.040 1321.600
\hypertarget{w02:cue:228}{}enable the making of democratic legislation，
% cue 229 1321.600 1328.060
\hypertarget{w02:cue:229}{}是不是能夠有助於達成的所謂民主立法，
% cue 230 1328.060 1331.020
\hypertarget{w02:cue:230}{}有助於達成，
% cue 231 1331.020 1332.020
\hypertarget{w02:cue:231}{}有助於完成，
% cue 232 1332.020 1334.010
\hypertarget{w02:cue:232}{}有助於創造民主立法。
% cue 233 1334.010 1335.140
\hypertarget{w02:cue:233}{}那我們就說，
% cue 234 1335.140 1337.390
\hypertarget{w02:cue:234}{}這個憲法解釋是正當的，
% cue 235 1337.390 1339.810
\hypertarget{w02:cue:235}{}一個正當的憲法解釋，
% cue 236 1339.810 1342.060
\hypertarget{w02:cue:236}{}它產生出來的規範，
% cue 237 1342.060 1345.880
\hypertarget{w02:cue:237}{}應有助於達成、創造民主立法。

\tc{0:22:25}
% cue 238 1345.880 1349.890
\hypertarget{w02:cue:238}{}where these conditions are satisfied，
% cue 239 1349.890 1353.470
\hypertarget{w02:cue:239}{}當這些條件都被滿足的話，
% cue 240 1353.470 1356.850
\hypertarget{w02:cue:240}{}the constitution can all legitimately be characterized，
% cue 241 1356.850 1360.950
\hypertarget{w02:cue:241}{}這個憲法就可以被正當的定性為，
% cue 242 1360.950 1363.200
\hypertarget{w02:cue:242}{}democratic 民主的憲法，
% cue 243 1363.200 1366.940
\hypertarget{w02:cue:243}{}independently of specific character
% cue 244 1366.940 1368.940
\hypertarget{w02:cue:244}{}the formal source，
% cue 245 1368.940 1373.190
\hypertarget{w02:cue:245}{}跟它正式的依據什麼東西，
% cue 246 1373.190 1375.570
\hypertarget{w02:cue:246}{}來界定它的特殊性格，
% cue 247 1375.570 1376.540
\hypertarget{w02:cue:247}{}是無關的。
% cue 248 1376.540 1377.950
\hypertarget{w02:cue:248}{}我現在講的是一般性，
% cue 249 1377.950 1380.040
\hypertarget{w02:cue:249}{}它現在提供的是一個，
% cue 250 1380.040 1382.350
\hypertarget{w02:cue:250}{}違憲審查一般性的論證，
% cue 251 1382.350 1384.830
\hypertarget{w02:cue:251}{}它現在講的是一般性，
% cue 252 1384.830 1386.790
\hypertarget{w02:cue:252}{}跟它特定的哪個來源無關，
% cue 253 1386.790 1388.820
\hypertarget{w02:cue:253}{}那是它等一下要講的來源。

\tc{0:23:08}
% cue 254 1388.820 1391.230
\hypertarget{w02:cue:254}{}好了，我現在講完這個事情，
% cue 255 1391.230 1393.640
\hypertarget{w02:cue:255}{}我說，今天我碰到一個鐵板，
% cue 256 1393.640 1395.640
\hypertarget{w02:cue:256}{}踢到一個鐵板，很難嗎？
% cue 257 1395.640 1399.010
\hypertarget{w02:cue:257}{}因為應該可以同意的嗎？
% cue 258 1399.010 1401.070
\hypertarget{w02:cue:258}{}就光光看那一串名詞，
% cue 259 1401.070 1403.130
\hypertarget{w02:cue:259}{}就已經被它搞得多昏倒障了，
% cue 260 1403.130 1404.310
\hypertarget{w02:cue:260}{}對吧？
% cue 261 1404.310 1405.310
\hypertarget{w02:cue:261}{}那我們等一下，
% cue 262 1405.310 1406.310
\hypertarget{w02:cue:262}{}像這個就是說，
% cue 263 1406.310 1408.310
\hypertarget{w02:cue:263}{}屬於你們要做的一些事情，
% cue 264 1408.310 1409.560
\hypertarget{w02:cue:264}{}剛開始，
% cue 265 1409.560 1410.880
\hypertarget{w02:cue:265}{}一回生兩回熟，
% cue 266 1410.880 1413.940
\hypertarget{w02:cue:266}{}你就慢慢把它那些關鍵語句，
% cue 267 1413.940 1415.910
\hypertarget{w02:cue:267}{}你先想辦法把它變成中文，
% cue 268 1415.910 1416.910
\hypertarget{w02:cue:268}{}你放在旁邊，
% cue 269 1416.910 1418.070
\hypertarget{w02:cue:269}{}然後你讀到的時候，
% cue 270 1418.070 1420.070
\hypertarget{w02:cue:270}{}你慢慢去推想，
% cue 271 1420.070 1421.070
\hypertarget{w02:cue:271}{}推想揣摩，
% cue 272 1421.070 1422.100
\hypertarget{w02:cue:272}{}它是什麼意思，
% cue 273 1422.100 1424.610
\hypertarget{w02:cue:273}{}因為這個英文確實已經很難。

\section{全文的節次與頁次}

\tc{0:23:44}
% cue 274 1424.610 1425.610
\hypertarget{w02:cue:274}{}好，我們現在來看一下這個結構，
% cue 275 1425.610 1426.630
\hypertarget{w02:cue:275}{}這個結構呢，
% cue 276 1426.630 1428.690
\hypertarget{w02:cue:276}{}在剛才我們看到的是 4，
% cue 277 1428.690 1430.820
\hypertarget{w02:cue:277}{}在 4 字，
% cue 278 1430.820 1432.820
\hypertarget{w02:cue:278}{}1,2,3,4,4 字結論，
% cue 279 1432.820 1434.780
\hypertarget{w02:cue:279}{}好，那 4 字前，
% cue 280 1434.780 1435.710
\hypertarget{w02:cue:280}{}我們先看 1，
% cue 281 1435.710 1436.800
\hypertarget{w02:cue:281}{}1 的前面呢，
% cue 282 1436.800 1437.800
\hypertarget{w02:cue:282}{}最早先是 abstract，
% cue 283 1437.800 1438.800
\hypertarget{w02:cue:283}{}剛才講過了。
% cue 284 1438.800 1440.800
\hypertarget{w02:cue:284}{}abstract 後面有一段文字，
% cue 285 1440.800 1441.770
\hypertarget{w02:cue:285}{}算是它沒有標題，
% cue 286 1441.770 1443.770
\hypertarget{w02:cue:286}{}那就算是它的引言吧，
% cue 287 1443.770 1445.790
\hypertarget{w02:cue:287}{}然後呢，
% cue 288 1445.790 1446.720
\hypertarget{w02:cue:288}{}1 標題，
% cue 289 1446.720 1448.100
\hypertarget{w02:cue:289}{}阿拉伯數字 1 在哪裡，
% cue 290 1448.100 1449.100
\hypertarget{w02:cue:290}{}在 527 頁，
% cue 291 1449.100 1450.770
\hypertarget{w02:cue:291}{}它的標題叫做
% cue 292 1450.770 1452.770
\hypertarget{w02:cue:292}{}Specific Legitimacy Thesis，
% cue 293 1452.770 1454.990
\hypertarget{w02:cue:293}{}我們等一下會從這裡談起。

\tc{0:24:14}
% cue 294 1454.990 1456.110
\hypertarget{w02:cue:294}{}因為，
% cue 295 1456.110 1458.110
\hypertarget{w02:cue:295}{}我們大部分通告的是這種情況，
% cue 296 1458.110 1461.310
\hypertarget{w02:cue:296}{}這是一種特定的正當性命題，
% cue 297 1461.310 1463.750
\hypertarget{w02:cue:297}{}我們現在講的是違憲審查，
% cue 298 1463.750 1466.190
\hypertarget{w02:cue:298}{}所以是講特定違憲審查，
% cue 299 1466.190 1468.190
\hypertarget{w02:cue:299}{}特定的違憲審查正當性命題，
% cue 300 1468.190 1469.760
\hypertarget{w02:cue:300}{}這是 1。
% cue 301 1469.760 1471.170
\hypertarget{w02:cue:301}{}那本文作者，
% cue 302 1471.170 1473.190
\hypertarget{w02:cue:302}{}它要講的不是特定的微正當性，
% cue 303 1473.190 1475.190
\hypertarget{w02:cue:303}{}它還講的是一般正當性，
% cue 304 1475.190 1476.190
\hypertarget{w02:cue:304}{}所以你會發現，
% cue 305 1476.190 1479.310
\hypertarget{w02:cue:305}{}它一定還會有各不同的招集，
% cue 306 1479.310 1480.810
\hypertarget{w02:cue:306}{}所以我們現在要看 2。

\tc{0:24:40}
% cue 307 1480.810 1482.010
\hypertarget{w02:cue:307}{}2 在哪裡呢，
% cue 308 1482.010 1484.370
\hypertarget{w02:cue:308}{}看你們沒有紙本，
% cue 309 1484.370 1485.370
\hypertarget{w02:cue:309}{}看怎麼怎麼找，
% cue 310 1485.370 1489.900
\hypertarget{w02:cue:310}{}2 在第 538 頁，
% cue 311 1489.900 1494.480
\hypertarget{w02:cue:311}{}538 頁。
% cue 312 1494.480 1496.480
\hypertarget{w02:cue:312}{}它要論證本文重點，
% cue 313 1496.480 1499.850
\hypertarget{w02:cue:313}{}一般性違憲審查，
% cue 314 1499.850 1502.330
\hypertarget{w02:cue:314}{}一般性正當性，
% cue 315 1502.330 1504.330
\hypertarget{w02:cue:315}{}一般正當性命題，
% cue 316 1504.330 1505.750
\hypertarget{w02:cue:316}{}它叫做 the general，
% cue 317 1505.750 1507.740
\hypertarget{w02:cue:317}{}剛才是 the specific，
% cue 318 1507.740 1509.730
\hypertarget{w02:cue:318}{}現在是 the general legitimacy。
% cue 319 1509.730 1511.980
\hypertarget{w02:cue:319}{}這個命題都有好多點，
% cue 320 1511.980 1513.200
\hypertarget{w02:cue:320}{}像你們慢慢看，
% cue 321 1513.200 1514.190
\hypertarget{w02:cue:321}{}這有 6 點，
% cue 322 1514.190 1515.190
\hypertarget{w02:cue:322}{}我沒有記住，
% cue 323 1515.190 1516.860
\hypertarget{w02:cue:323}{}6 點複雜。
% cue 324 1516.860 1518.630
\hypertarget{w02:cue:324}{}一個人寫文章要寫到這樣，
% cue 325 1518.630 1519.630
\hypertarget{w02:cue:325}{}實際是會發，
% cue 326 1519.630 1521.200
\hypertarget{w02:cue:326}{}你就發現外國人頭腦，
% cue 327 1521.200 1523.750
\hypertarget{w02:cue:327}{}跟我們頭腦真的長得不，
% cue 328 1523.750 1524.830
\hypertarget{w02:cue:328}{}我這個文道如果這樣寫出來，
% cue 329 1524.830 1525.830
\hypertarget{w02:cue:329}{}在臺灣一定登不出來。

\tc{0:25:25}
% cue 330 1525.830 1526.780
\hypertarget{w02:cue:330}{}申告一，
% cue 331 1526.780 1529.770
\hypertarget{w02:cue:331}{}第三點，
% cue 332 1529.770 1530.770
\hypertarget{w02:cue:332}{}the constitutional thesis，
% cue 333 1530.770 1532.930
\hypertarget{w02:cue:333}{}這叫憲法命題。
% cue 334 1532.930 1535.290
\hypertarget{w02:cue:334}{}我們就出現了三個命題，
% cue 335 1535.290 1536.290
\hypertarget{w02:cue:335}{}什麼都沒講究其三個命題，
% cue 336 1536.290 1539.560
\hypertarget{w02:cue:336}{}三個命題裡面還有一大堆命題。
% cue 337 1539.560 1542.230
\hypertarget{w02:cue:337}{}那它為什麼要講憲法命題，
% cue 338 1542.230 1544.350
\hypertarget{w02:cue:338}{}我們等一下看看它怎麼解釋，
% cue 339 1544.350 1547.270
\hypertarget{w02:cue:339}{}憲法命題是在 544 頁，
% cue 340 1547.270 1548.270
\hypertarget{w02:cue:340}{}阿拉伯數字三，
% cue 341 1548.270 1550.240
\hypertarget{w02:cue:341}{}the Constitutional Law Thesis，
% cue 342 1550.240 1552.490
\hypertarget{w02:cue:342}{}那這裡面有分 a，
% cue 343 1552.490 1553.550
\hypertarget{w02:cue:343}{}546，
% cue 344 1553.550 1557.700
\hypertarget{w02:cue:344}{}the hypothesis of unwritten constitution，
% cue 345 1557.700 1562.220
\hypertarget{w02:cue:345}{}是有關於不成文憲法的假設。
% cue 346 1562.220 1566.490
\hypertarget{w02:cue:346}{}有不成文那就有成文，
% cue 347 1566.490 1568.480
\hypertarget{w02:cue:347}{}那它先把不成文講在前面，
% cue 348 1568.480 1570.430
\hypertarget{w02:cue:348}{}因為這是一個加拿大人，
% cue 349 1570.430 1571.440
\hypertarget{w02:cue:349}{}它原來是英國的傳統，
% cue 350 1571.440 1573.530
\hypertarget{w02:cue:350}{}所以它會先講加拿大熟悉，
% cue 351 1573.530 1575.530
\hypertarget{w02:cue:351}{}但是加拿大後來也有成文憲法，
% cue 352 1575.530 1576.530
\hypertarget{w02:cue:352}{}1982\footnote{加拿大的 Constitution Act, 1867 為英國國會制定的成文憲法性法律。1982 年英國國會通過 Canada Act 1982，依其第 1 條制定附表 B 所載的 Constitution Act, 1982，該法第一篇為 Canadian Charter of Rights and Freedoms（加拿大權利與自由憲章）；Canada Act 1982 第 2 條規定，此後英國國會通過的法律不再成為加拿大法律的一部分。\url{https://laws-lois.justice.gc.ca/eng/Const/FullText.html}。} 以後。

\tc{0:26:16}
% cue 353 1576.530 1577.680
\hypertarget{w02:cue:353}{}好，
% cue 354 1577.680 1579.680
\hypertarget{w02:cue:354}{}那我們現在看 551 頁，
% cue 355 1579.680 1581.460
\hypertarget{w02:cue:355}{}b
% cue 356 1581.460 1583.460
\hypertarget{w02:cue:356}{}the hypothesis of written constitution，
% cue 357 1583.460 1585.710
\hypertarget{w02:cue:357}{}有關於成文憲法的假設。
% cue 358 1585.710 1587.510
\hypertarget{w02:cue:358}{}好，
% cue 359 1587.510 1588.830
\hypertarget{w02:cue:359}{}所以它剛剛講，
% cue 360 1588.830 1589.820
\hypertarget{w02:cue:360}{}它這一套，
% cue 361 1589.820 1592.860
\hypertarget{w02:cue:361}{}它這一套所謂違憲審查，
% cue 362 1592.860 1594.860
\hypertarget{w02:cue:362}{}一般正當性的論述，
% cue 363 1594.860 1597.750
\hypertarget{w02:cue:363}{}既適用在不成文憲法，
% cue 364 1597.750 1599.970
\hypertarget{w02:cue:364}{}也適用在成文憲法，
% cue 365 1599.970 1601.060
\hypertarget{w02:cue:365}{}所以它就是簡單講，
% cue 366 1601.060 1603.030
\hypertarget{w02:cue:366}{}我分別適用給你看。

\tc{0:26:43}
% cue 367 1603.030 1605.530
\hypertarget{w02:cue:367}{}那我再下來，
% cue 368 1605.530 1607.450
\hypertarget{w02:cue:368}{}才是 561 頁，
% cue 369 1607.450 1609.670
\hypertarget{w02:cue:369}{}就是 conclusion，
% cue 370 1609.670 1610.670
\hypertarget{w02:cue:370}{}剛才我念了 conclusion，
% cue 371 1610.670 1612.520
\hypertarget{w02:cue:371}{}猜想大概很難，
% cue 372 1612.520 1615.870
\hypertarget{w02:cue:372}{}看到一個文章，
% cue 373 1615.870 1617.240
\hypertarget{w02:cue:373}{}conclusion 都聽不懂。
% cue 374 1617.240 1621.060
\hypertarget{w02:cue:374}{}就是我剛才念的那個，
% cue 375 1621.060 1622.150
\hypertarget{w02:cue:375}{}我自己覺得，
% cue 376 1622.150 1623.150
\hypertarget{w02:cue:376}{}我用翻的，
% cue 377 1623.150 1624.150
\hypertarget{w02:cue:377}{}我都翻的很辛苦，
% cue 378 1624.150 1625.530
\hypertarget{w02:cue:378}{}我猜想你們也是聽，
% cue 379 1625.530 1628.000
\hypertarget{w02:cue:379}{}應該是聽得霧煞煞。
% cue 380 1628.000 1629.440
\hypertarget{w02:cue:380}{}那我就說，
% cue 381 1629.440 1630.440
\hypertarget{w02:cue:381}{}我們今天看到這個就是，
% cue 382 1630.440 1632.590
\hypertarget{w02:cue:382}{}英零快甜，
% cue 383 1632.590 1634.160
\hypertarget{w02:cue:383}{}那我覺得，
% cue 384 1634.160 1636.410
\hypertarget{w02:cue:384}{}我能講多少算多少。

\section{引言與反多數決困難}

\tc{0:27:16}
% cue 385 1636.410 1637.860
\hypertarget{w02:cue:385}{}來我們看一下，
% cue 386 1637.860 1638.920
\hypertarget{w02:cue:386}{}從頭開始，
% cue 387 1638.920 1640.010
\hypertarget{w02:cue:387}{}那個引言，
% cue 388 1640.010 1644.400
\hypertarget{w02:cue:388}{}引言相對是最好的。
% cue 389 1644.400 1646.680
\hypertarget{w02:cue:389}{}我為了避免讓你們完全聽不懂，
% cue 390 1646.680 1650.500
\hypertarget{w02:cue:390}{}引言一定要 traditionally，
% cue 391 1650.500 1651.500
\hypertarget{w02:cue:391}{}傳統來說，
% cue 392 1651.500 1652.500
\hypertarget{w02:cue:392}{}the answer to the question
% cue 393 1652.500 1654.660
\hypertarget{w02:cue:393}{}of the legitimacy of judicial review
% cue 394 1654.660 1655.820
\hypertarget{w02:cue:394}{}of the legislation，
% cue 395 1655.820 1658.560
\hypertarget{w02:cue:395}{}只把審查立法的，
% cue 396 1658.560 1660.560
\hypertarget{w02:cue:396}{}違憲審查正當性，
% cue 397 1660.560 1663.530
\hypertarget{w02:cue:397}{}has been conceived as the solutions
% cue 398 1663.530 1664.530
\hypertarget{w02:cue:398}{}to what has been called
% cue 399 1664.530 1666.560
\hypertarget{w02:cue:399}{}the countermajoritarian difficulty\footnote{語出 Alexander M. Bickel, The Least Dangerous Branch（New Haven: Yale University Press, 1962），本文第 525 頁註 1 引該書第 16 至 23 頁。Bickel 該節首句為「The root difficulty is that judicial review is a countermajoritarian force in our system.」（根本的難題在於，違憲審查在我們的體制裡是一股反多數的力量）。課綱譯作「反抗多數（決）的難題」，第 4 週指定該書第 16 至 26 頁。\cite{bickel2004}。}。

\tc{0:27:46}
% cue 400 1666.560 1668.100
\hypertarget{w02:cue:400}{}我們為什麼要討論，
% cue 401 1668.100 1669.100
\hypertarget{w02:cue:401}{}違憲審查正當性，
% cue 402 1669.100 1671.100
\hypertarget{w02:cue:402}{}我們是把它拿來對治，
% cue 403 1671.100 1673.220
\hypertarget{w02:cue:403}{}一個所謂的，
% cue 404 1673.220 1675.220
\hypertarget{w02:cue:404}{}反多數決困難，
% cue 405 1675.220 1677.220
\hypertarget{w02:cue:405}{}這個題目的答案，
% cue 406 1677.220 1678.620
\hypertarget{w02:cue:406}{}因為，
% cue 407 1678.620 1679.760
\hypertarget{w02:cue:407}{}我們碰到了，
% cue 408 1679.760 1680.760
\hypertarget{w02:cue:408}{}反多數決困難，
% cue 409 1680.760 1681.790
\hypertarget{w02:cue:409}{}怎麼辦？
% cue 410 1681.790 1683.460
\hypertarget{w02:cue:410}{}所以我們提出，
% cue 411 1683.460 1684.460
\hypertarget{w02:cue:411}{}違憲審查的正當性，
% cue 412 1684.460 1685.650
\hypertarget{w02:cue:412}{}違憲審查正當性，
% cue 413 1685.650 1687.650
\hypertarget{w02:cue:413}{}就是解答，
% cue 414 1687.650 1689.120
\hypertarget{w02:cue:414}{}違憲審查，
% cue 415 1689.120 1691.090
\hypertarget{w02:cue:415}{}反多數決困難的答案，
% cue 416 1691.090 1692.880
\hypertarget{w02:cue:416}{}這是第一句話。
% cue 417 1692.880 1694.670
\hypertarget{w02:cue:417}{}有聽懂嗎？
% cue 418 1694.670 1696.410
\hypertarget{w02:cue:418}{}應該可以。

\tc{0:28:16}
% cue 419 1696.410 1697.410
\hypertarget{w02:cue:419}{}on this assumption，
% cue 420 1697.410 1698.410
\hypertarget{w02:cue:420}{}在這個假設之下，
% cue 421 1698.410 1703.230
\hypertarget{w02:cue:421}{}no law is acceptable as legitimate，
% cue 422 1703.230 1705.480
\hypertarget{w02:cue:422}{}沒有任何一個法律，
% cue 423 1705.480 1706.480
\hypertarget{w02:cue:423}{}可以被接納為正當的，
% cue 424 1706.480 1709.180
\hypertarget{w02:cue:424}{}沒有一個法律是正當的。
% cue 425 1709.180 1710.180
\hypertarget{w02:cue:425}{}注意這句。

\tc{0:28:30}
% cue 426 1710.180 1711.180
\hypertarget{w02:cue:426}{}unless change line，
% cue 427 1711.180 1712.240
\hypertarget{w02:cue:427}{}除非，
% cue 428 1712.240 1713.300
\hypertarget{w02:cue:428}{}換句話說，
% cue 429 1713.300 1715.520
\hypertarget{w02:cue:429}{}推定為不正當，
% cue 430 1715.520 1717.640
\hypertarget{w02:cue:430}{}除非你能證明它會正當。
% cue 431 1717.640 1723.070
\hypertarget{w02:cue:431}{}unless it expresses the will of representative of the people
% cue 432 1723.070 1726.350
\hypertarget{w02:cue:432}{}as determined in accordance with the majority rule，
% cue 433 1726.350 1728.670
\hypertarget{w02:cue:433}{}除非是依照多數決的原則，
% cue 434 1728.670 1734.410
\hypertarget{w02:cue:434}{}由人民的代表所表達的意思。
% cue 435 1734.410 1735.950
\hypertarget{w02:cue:435}{}換句話說，
% cue 436 1735.950 1737.010
\hypertarget{w02:cue:436}{}是經過多數決，
% cue 437 1737.010 1740.010
\hypertarget{w02:cue:437}{}除非是經過代表人民的代表，
% cue 438 1740.010 1742.170
\hypertarget{w02:cue:438}{}民意代表多數決所通過的，
% cue 439 1742.170 1744.230
\hypertarget{w02:cue:439}{}代表了人民的意思。
% cue 440 1744.230 1749.090
\hypertarget{w02:cue:440}{}不然這個法律是沒有辦法認為，
% cue 441 1749.090 1752.200
\hypertarget{w02:cue:441}{}所以他一直開始跟你講的這個，
% cue 442 1752.200 1754.250
\hypertarget{w02:cue:442}{}所以他一直開始跟你講的這個，
% cue 443 1754.250 1757.730
\hypertarget{w02:cue:443}{}一定要能顯示明明，
% cue 444 1757.730 1760.040
\hypertarget{w02:cue:444}{}要合乎多數決。
% cue 445 1760.040 1761.160
\hypertarget{w02:cue:445}{}ok?

\tc{0:29:21}
% cue 446 1761.160 1763.250
\hypertarget{w02:cue:446}{}not been elected
% cue 447 1763.250 1765.670
\hypertarget{w02:cue:447}{}and not otherwise
% cue 448 1765.670 1767.680
\hypertarget{w02:cue:448}{}since，
% cue 449 1767.680 1768.600
\hypertarget{w02:cue:449}{}那是另外一句，
% cue 450 1768.600 1769.840
\hypertarget{w02:cue:450}{}我剛剛講到這裡，
% cue 451 1769.840 1770.840
\hypertarget{w02:cue:451}{}應該可以聽得懂，
% cue 452 1770.840 1772.820
\hypertarget{w02:cue:452}{}對不對？

\tc{0:29:32}
% cue 453 1772.820 1773.950
\hypertarget{w02:cue:453}{}我們講一般，
% cue 454 1773.950 1775.980
\hypertarget{w02:cue:454}{}講違憲審查正當性，
% cue 455 1775.980 1779.170
\hypertarget{w02:cue:455}{}是因為我們通過違憲審查，
% cue 456 1779.170 1780.260
\hypertarget{w02:cue:456}{}多數決的困難，
% cue 457 1780.260 1783.060
\hypertarget{w02:cue:457}{}那我們對這個問題要提一個答案。
% cue 458 1783.060 1784.380
\hypertarget{w02:cue:458}{}那答案呢，
% cue 459 1784.380 1785.440
\hypertarget{w02:cue:459}{}就是我們提出來，
% cue 460 1785.440 1787.440
\hypertarget{w02:cue:460}{}違憲審查正當性作為答案，
% cue 461 1787.440 1789.430
\hypertarget{w02:cue:461}{}我們現在先不講它內容是什麼，
% cue 462 1789.430 1791.030
\hypertarget{w02:cue:462}{}為什麼會這樣子。

\tc{0:29:51}
% cue 463 1791.030 1792.090
\hypertarget{w02:cue:463}{}since，
% cue 464 1792.090 1793.340
\hypertarget{w02:cue:464}{}因為，
% cue 465 1793.340 1794.340
\hypertarget{w02:cue:465}{}judges do not represent people，
% cue 466 1794.340 1795.340
\hypertarget{w02:cue:466}{}法官不代表人民，
% cue 467 1795.340 1798.100
\hypertarget{w02:cue:467}{}先不管代表是什麼意思。
% cue 468 1798.100 1800.380
\hypertarget{w02:cue:468}{}not been elected，
% cue 469 1800.380 1801.380
\hypertarget{w02:cue:469}{}也沒有被人民選取，
% cue 470 1801.380 1802.660
\hypertarget{w02:cue:470}{}and not otherwise
% cue 471 1802.660 1805.260
\hypertarget{w02:cue:471}{}personifying the citizenry，
% cue 472 1805.260 1808.480
\hypertarget{w02:cue:472}{}法官也不能說他象徵人民，
% cue 473 1808.480 1810.480
\hypertarget{w02:cue:473}{}法官有這資格講他象徵人民嗎？
% cue 474 1810.480 1813.860
\hypertarget{w02:cue:474}{}and since they are not accountable to the people，
% cue 475 1813.860 1814.950
\hypertarget{w02:cue:475}{}也不對人民負責，
% cue 476 1814.950 1816.200
\hypertarget{w02:cue:476}{}or to citizen，
% cue 477 1816.200 1817.320
\hypertarget{w02:cue:477}{}也不用對人民負責，
% cue 478 1817.320 1819.260
\hypertarget{w02:cue:478}{}也不對公民負責。
% cue 479 1819.260 1825.090
\hypertarget{w02:cue:479}{}judicial review of legislation constitutes a countermajoritarian force
% cue 480 1825.090 1828.400
\hypertarget{w02:cue:480}{}in a system of governance committed to democracy。

\tc{0:30:28}
% cue 481 1828.400 1833.320
\hypertarget{w02:cue:481}{}這就是為什麼講反多數決的困境，
% cue 482 1833.320 1839.320
\hypertarget{w02:cue:482}{}你把一個聲稱要實現民主的政府體制，
% cue 483 1839.320 1845.200
\hypertarget{w02:cue:483}{}那你卻是用違反民主不代表民意的違憲審查，
% cue 484 1845.200 1848.000
\hypertarget{w02:cue:484}{}來推翻多數決人所通過的，
% cue 485 1848.000 1851.480
\hypertarget{w02:cue:485}{}它變成一個反多數決的力量，
% cue 486 1851.480 1854.660
\hypertarget{w02:cue:486}{}認為這樣是合理。
% cue 487 1854.660 1856.200
\hypertarget{w02:cue:487}{}所以他才會說，
% cue 488 1856.200 1859.480
\hypertarget{w02:cue:488}{}有 counter-majoritarian difficulty。

\tc{0:30:59}
% cue 489 1859.480 1862.700
\hypertarget{w02:cue:489}{}所以這段話是在解釋，
% cue 490 1862.700 1864.760
\hypertarget{w02:cue:490}{}什麼叫做 counter-majoritarian difficulty，
% cue 491 1864.760 1866.790
\hypertarget{w02:cue:491}{}你也會碰到，
% cue 492 1866.790 1868.790
\hypertarget{w02:cue:492}{}這叫做 anti-democratic character，
% cue 493 1868.790 1870.940
\hypertarget{w02:cue:493}{}反民主的性格，
% cue 494 1870.940 1873.130
\hypertarget{w02:cue:494}{}違憲審查的反民主性格。
% cue 495 1873.130 1874.420
\hypertarget{w02:cue:495}{}好。

\tc{0:31:14}
% cue 496 1874.420 1875.600
\hypertarget{w02:cue:496}{}it follows that，
% cue 497 1875.600 1877.330
\hypertarget{w02:cue:497}{}接下來，
% cue 498 1877.330 1881.740
\hypertarget{w02:cue:498}{}those who claim that this institution is legitimate
% cue 499 1881.740 1884.130
\hypertarget{w02:cue:499}{}have the burden of showing，
% cue 500 1884.130 1888.850
\hypertarget{w02:cue:500}{}在這種以民主為職志的，
% cue 501 1888.850 1890.260
\hypertarget{w02:cue:501}{}體制之下，
% cue 502 1890.260 1893.670
\hypertarget{w02:cue:502}{}居然出現了這種反民主，
% cue 503 1893.670 1894.670
\hypertarget{w02:cue:503}{}反多數決的力量。
% cue 504 1894.670 1895.580
\hypertarget{w02:cue:504}{}沒有違憲審查，
% cue 505 1895.580 1900.780
\hypertarget{w02:cue:505}{}那麼你主張違憲審查這個制度是正當的人，
% cue 506 1900.780 1902.880
\hypertarget{w02:cue:506}{}你就負有舉證責任，
% cue 507 1902.880 1904.120
\hypertarget{w02:cue:507}{}you have the burden of showing
% cue 508 1904.120 1908.400
\hypertarget{w02:cue:508}{}that it can be reconciled with the underlying assumption
% cue 509 1908.400 1910.400
\hypertarget{w02:cue:509}{}of the democratic governance。
% cue 510 1910.400 1911.490
\hypertarget{w02:cue:510}{}你說不怪不怪，
% cue 511 1911.490 1912.490
\hypertarget{w02:cue:511}{}它為什麼是正當的？
% cue 512 1912.490 1917.550
\hypertarget{w02:cue:512}{}因為我有辦法調和它背後所依據的這個民主統治，
% cue 513 1917.550 1919.990
\hypertarget{w02:cue:513}{}這個基礎的假設，
% cue 514 1919.990 1923.010
\hypertarget{w02:cue:514}{}那你一定要說它怎麼不違反民主，
% cue 515 1923.010 1924.260
\hypertarget{w02:cue:515}{}它怎麼可以調和，
% cue 516 1924.260 1927.200
\hypertarget{w02:cue:516}{}不然它就是沒有辦法過這個難關呢？

\tc{0:32:07}
% cue 517 1927.200 1928.820
\hypertarget{w02:cue:517}{}the various answer，
% cue 518 1928.820 1929.820
\hypertarget{w02:cue:518}{}好，這裡要注意了，
% cue 519 1929.820 1930.820
\hypertarget{w02:cue:519}{}我們有很多種答案，
% cue 520 1930.820 1934.110
\hypertarget{w02:cue:520}{}to the question of the legitimacy of the judicial review
% cue 521 1934.110 1936.200
\hypertarget{w02:cue:521}{}have generally postulated，
% cue 522 1936.200 1938.510
\hypertarget{w02:cue:522}{}postulated，
% cue 523 1938.510 1939.990
\hypertarget{w02:cue:523}{}通常我們是假設，
% cue 524 1939.990 1945.180
\hypertarget{w02:cue:524}{}that no form of judicial review is acceptable as a legitimate
% cue 525 1945.180 1955.450
\hypertarget{w02:cue:525}{}unless it can be shown that it is specifically authorized by some positive democratic legitimating fact。

\tc{0:32:35}
% cue 526 1955.450 1958.430
\hypertarget{w02:cue:526}{}這句話就比較難解了。
% cue 527 1958.430 1961.160
\hypertarget{w02:cue:527}{}我們這樣講的時候，
% cue 528 1961.160 1963.190
\hypertarget{w02:cue:528}{}我們通常是假設，
% cue 529 1963.190 1965.540
\hypertarget{w02:cue:529}{}假設什麼呢？
% cue 530 1965.540 1966.730
\hypertarget{w02:cue:530}{}沒有任何一種，
% cue 531 1966.730 1970.240
\hypertarget{w02:cue:531}{}沒有任何一種形式的違憲審查，
% cue 532 1970.240 1972.870
\hypertarget{w02:cue:532}{}是可以被認定為正當的，
% cue 533 1972.870 1975.210
\hypertarget{w02:cue:533}{}除非你能夠證明，
% cue 534 1975.210 1978.710
\hypertarget{w02:cue:534}{}它特定的，
% cue 535 1978.710 1980.710
\hypertarget{w02:cue:535}{}特別的被授權，
% cue 536 1980.710 1981.660
\hypertarget{w02:cue:536}{}具體的被授權，
% cue 537 1981.660 1982.740
\hypertarget{w02:cue:537}{}被什麼東西授權，
% cue 538 1982.740 1984.870
\hypertarget{w02:cue:538}{}某一種積極的，
% cue 539 1984.870 1986.780
\hypertarget{w02:cue:539}{}民主的，
% cue 540 1986.780 1988.000
\hypertarget{w02:cue:540}{}正當性的事實。
% cue 541 1988.000 1989.580
\hypertarget{w02:cue:541}{}這句話真難翻，
% cue 542 1989.580 1992.560
\hypertarget{w02:cue:542}{}聽得懂嗎？
% cue 543 1992.560 1995.860
\hypertarget{w02:cue:543}{}一定要有一個 positive，
% cue 544 1995.860 1997.020
\hypertarget{w02:cue:544}{}不是 negative，
% cue 545 1997.020 1998.970
\hypertarget{w02:cue:545}{}一個積極的，
% cue 546 1998.970 2002.190
\hypertarget{w02:cue:546}{}一個正向的 legitimating fact，
% cue 547 2002.190 2004.730
\hypertarget{w02:cue:547}{}使它變得正當的事實。
% cue 548 2004.730 2007.010
\hypertarget{w02:cue:548}{}而且這個使它變成正當，
% cue 549 2007.010 2008.010
\hypertarget{w02:cue:549}{}因為我們是民主國，
% cue 550 2008.010 2010.010
\hypertarget{w02:cue:550}{}所以它一定是民主的正當事，
% cue 551 2010.010 2014.860
\hypertarget{w02:cue:551}{}使它民主又正當的這樣的一個授權，
% cue 552 2014.860 2017.150
\hypertarget{w02:cue:552}{}這樣你才能夠說，
% cue 553 2017.150 2018.240
\hypertarget{w02:cue:553}{}違憲審查是正當的。
% cue 554 2018.240 2020.410
\hypertarget{w02:cue:554}{}那這個正當的授權呢，
% cue 555 2020.410 2021.590
\hypertarget{w02:cue:555}{}一定是特定的，
% cue 556 2021.590 2022.590
\hypertarget{w02:cue:556}{}所以是 specifically。
% cue 557 2022.590 2025.250
\hypertarget{w02:cue:557}{}那這是等一下它會講的。

\tc{0:33:45}
% cue 558 2025.250 2026.530
\hypertarget{w02:cue:558}{}下面一句，
% cue 559 2026.530 2029.070
\hypertarget{w02:cue:559}{}They ought to establish what I call
% cue 560 2029.070 2032.380
\hypertarget{w02:cue:560}{}the specific legitimacy of judiciary.
% cue 561 2032.380 2035.940
\hypertarget{w02:cue:561}{}也就是一般形形，
% cue 562 2035.940 2038.060
\hypertarget{w02:cue:562}{}你必須要能夠證明，
% cue 563 2038.060 2040.560
\hypertarget{w02:cue:563}{}你具有特定的，
% cue 564 2040.560 2041.560
\hypertarget{w02:cue:564}{}違憲審查正當性，
% cue 565 2041.560 2042.910
\hypertarget{w02:cue:565}{}你特定的被授權，
% cue 566 2042.910 2045.100
\hypertarget{w02:cue:566}{}你特定的有一個，
% cue 567 2045.100 2046.030
\hypertarget{w02:cue:567}{}符合民主的，
% cue 568 2046.030 2049.060
\hypertarget{w02:cue:568}{}使它成為正當的授權事實，
% cue 569 2049.060 2051.760
\hypertarget{w02:cue:569}{}這樣有個特定的授權事實。

\tc{0:34:11}
% cue 570 2051.760 2053.370
\hypertarget{w02:cue:570}{}In this text
% cue 571 2053.370 2054.370
\hypertarget{w02:cue:571}{}I should argue，
% cue 572 2054.370 2055.490
\hypertarget{w02:cue:572}{}在本文裡，
% cue 573 2055.490 2056.490
\hypertarget{w02:cue:573}{}我主張，
% cue 574 2056.490 2060.770
\hypertarget{w02:cue:574}{}the institution of the judicial review of legislation does not need，
% cue 575 2060.770 2062.250
\hypertarget{w02:cue:575}{}我不認為，
% cue 576 2062.250 2067.760
\hypertarget{w02:cue:576}{}需要 specifically authorized in order to be accepted as legitimate.
% cue 577 2067.760 2070.290
\hypertarget{w02:cue:577}{}討論民主，
% cue 578 2070.290 2071.290
\hypertarget{w02:cue:578}{}討論違憲審查正當性的授權，
% cue 579 2071.290 2072.290
\hypertarget{w02:cue:579}{}一般都是強，
% cue 580 2072.290 2073.310
\hypertarget{w02:cue:580}{}特定授權。
% cue 581 2073.310 2074.460
\hypertarget{w02:cue:581}{}我不認為，
% cue 582 2074.460 2076.800
\hypertarget{w02:cue:582}{}我不認為需要特定授權。
% cue 583 2076.800 2078.000
\hypertarget{w02:cue:583}{}那是什麼意思呢？

\tc{0:34:38}
% cue 584 2078.000 2079.060
\hypertarget{w02:cue:584}{}再看下一句，
% cue 585 2079.060 2080.380
\hypertarget{w02:cue:585}{}the mere fact，
% cue 586 2080.380 2082.370
\hypertarget{w02:cue:586}{}下面這個單純的事實，
% cue 587 2082.370 2086.460
\hypertarget{w02:cue:587}{}the judiciary already exercises effective coercive power，
% cue 588 2086.460 2092.570
\hypertarget{w02:cue:588}{}司法已經能夠行使有效的強制力。
% cue 589 2092.570 2093.730
\hypertarget{w02:cue:589}{}你可以不遵守嗎？
% cue 590 2093.730 2094.840
\hypertarget{w02:cue:590}{}當然，
% cue 591 2094.840 2096.170
\hypertarget{w02:cue:591}{}你要說，
% cue 592 2096.170 2097.170
\hypertarget{w02:cue:592}{}我們大法官做了解釋，
% cue 593 2097.170 2098.170
\hypertarget{w02:cue:593}{}努力法者不遵從，
% cue 594 2098.170 2100.200
\hypertarget{w02:cue:594}{}那他就沒有強制力。
% cue 595 2100.200 2102.490
\hypertarget{w02:cue:595}{}那我們現在是講他有強制力，
% cue 596 2102.490 2103.730
\hypertarget{w02:cue:596}{}and political authority，
% cue 597 2103.730 2105.730
\hypertarget{w02:cue:597}{}而且有政治權威，
% cue 598 2105.730 2107.730
\hypertarget{w02:cue:598}{}within society，
% cue 599 2107.730 2108.880
\hypertarget{w02:cue:599}{}這是一個事實。
% cue 600 2108.880 2110.880
\hypertarget{w02:cue:600}{}這是一個這樣的單純事實。
% cue 601 2110.880 2111.810
\hypertarget{w02:cue:601}{}constitute，
% cue 602 2111.810 2114.260
\hypertarget{w02:cue:602}{}前面是 fact，
% cue 603 2114.260 2115.230
\hypertarget{w02:cue:603}{}主詞，
% cue 604 2115.230 2118.310
\hypertarget{w02:cue:604}{}這個事實就構成一個 sufficient reason，
% cue 605 2118.310 2120.750
\hypertarget{w02:cue:605}{}足充分的理由，
% cue 606 2120.750 2123.710
\hypertarget{w02:cue:606}{}to recognize what I shall call the general legitimacy of the judicial review。

\tc{0:35:23}
% cue 607 2123.710 2126.700
\hypertarget{w02:cue:607}{}我認為不需要特定的，
% cue 608 2126.700 2129.050
\hypertarget{w02:cue:608}{}不需要再有特定的授權，
% cue 609 2129.050 2130.110
\hypertarget{w02:cue:609}{}為什麼？
% cue 610 2130.110 2133.480
\hypertarget{w02:cue:610}{}因為我們都知道司法審查，
% cue 611 2133.480 2135.480
\hypertarget{w02:cue:611}{}這是一個客觀存在的事實，
% cue 612 2135.480 2136.570
\hypertarget{w02:cue:612}{}它有拘束力，
% cue 613 2136.570 2137.700
\hypertarget{w02:cue:613}{}這樣的東西就夠了，
% cue 614 2137.700 2142.800
\hypertarget{w02:cue:614}{}那我們就可以承認它有一般的正當性。

\tc{0:35:42}
% cue 615 2142.800 2146.070
\hypertarget{w02:cue:615}{}my argument does not deny，
% cue 616 2146.070 2147.480
\hypertarget{w02:cue:616}{}我並不否認，
% cue 617 2147.480 2152.700
\hypertarget{w02:cue:617}{}some form of specific authorization could support judicial review，
% cue 618 2152.700 2153.700
\hypertarget{w02:cue:618}{}我並不否認，
% cue 619 2153.700 2156.890
\hypertarget{w02:cue:619}{}也許有一種特定的授權也可以支撐違憲審查，
% cue 620 2156.890 2160.330
\hypertarget{w02:cue:620}{}and enhance its moral legitimacy，
% cue 621 2160.330 2161.260
\hypertarget{w02:cue:621}{}我並不否認，
% cue 622 2161.260 2162.460
\hypertarget{w02:cue:622}{}否定，
% cue 623 2162.460 2163.490
\hypertarget{w02:cue:623}{}我講的一般性，
% cue 624 2163.490 2164.610
\hypertarget{w02:cue:624}{}正一般正當性以外，
% cue 625 2164.610 2166.610
\hypertarget{w02:cue:625}{}可以有特定正當性，
% cue 626 2166.610 2167.610
\hypertarget{w02:cue:626}{}我不否認這件事情。
% cue 627 2167.610 2170.020
\hypertarget{w02:cue:627}{}但是這不是我的目的。

\tc{0:36:10}
% cue 628 2170.020 2171.210
\hypertarget{w02:cue:628}{}my claim，
% cue 629 2171.210 2172.210
\hypertarget{w02:cue:629}{}我的主張就是，
% cue 630 2172.210 2175.320
\hypertarget{w02:cue:630}{}that such support is not necessary，
% cue 631 2175.320 2177.940
\hypertarget{w02:cue:631}{}但我認為是不必要的。
% cue 632 2177.940 2181.960
\hypertarget{w02:cue:632}{}judicial review of legislation is legitimate and in principle，
% cue 633 2181.960 2182.960
\hypertarget{w02:cue:633}{}in principle 圈起來。
% cue 634 2182.960 2184.340
\hypertarget{w02:cue:634}{}我的主張就是，
% cue 635 2184.340 2187.310
\hypertarget{w02:cue:635}{}違憲審查原則上是正當的。
% cue 636 2187.310 2190.520
\hypertarget{w02:cue:636}{}你看看這兩個字正好相反，
% cue 637 2190.520 2193.810
\hypertarget{w02:cue:637}{}specific legitimacy is 原則上不正當的，
% cue 638 2193.810 2195.830
\hypertarget{w02:cue:638}{}除非你得到特定授權，
% cue 639 2195.830 2196.800
\hypertarget{w02:cue:639}{}你才能證明為正當。

\tc{0:36:36}
% cue 640 2196.800 2199.060
\hypertarget{w02:cue:640}{}我的說法正好相反。
% cue 641 2199.060 2200.060
\hypertarget{w02:cue:641}{}作者說，
% cue 642 2200.060 2202.320
\hypertarget{w02:cue:642}{}我認為有這個客觀的事實在那裡存在，
% cue 643 2202.320 2203.320
\hypertarget{w02:cue:643}{}它運行了那麼多年，
% cue 644 2203.320 2205.320
\hypertarget{w02:cue:644}{}它其實大家都接受了，
% cue 645 2205.320 2206.290
\hypertarget{w02:cue:645}{}所以呢，
% cue 646 2206.290 2208.290
\hypertarget{w02:cue:646}{}我們應該推定它為合法，
% cue 647 2208.290 2209.480
\hypertarget{w02:cue:647}{}推定它為正當，
% cue 648 2209.480 2211.730
\hypertarget{w02:cue:648}{}原則上是正當。

\tc{0:36:51}
% cue 649 2211.730 2213.040
\hypertarget{w02:cue:649}{}indeed
% cue 650 2213.040 2214.040
\hypertarget{w02:cue:650}{}I shall argue，
% cue 651 2214.040 2215.130
\hypertarget{w02:cue:651}{}我要主張的是，
% cue 652 2215.130 2216.120
\hypertarget{w02:cue:652}{}It is only，
% cue 653 2216.120 2219.210
\hypertarget{w02:cue:653}{}It is not only morally permissible，
% cue 654 2219.210 2220.560
\hypertarget{w02:cue:654}{}這個所謂的正當，
% cue 655 2220.560 2221.560
\hypertarget{w02:cue:655}{}不只是道德上，
% cue 656 2221.560 2224.030
\hypertarget{w02:cue:656}{}可允許有違憲審查。
% cue 657 2224.030 2226.120
\hypertarget{w02:cue:657}{}It might be required，
% cue 658 2226.120 2227.120
\hypertarget{w02:cue:658}{}在某些情況之下，
% cue 659 2227.120 2229.120
\hypertarget{w02:cue:659}{}它甚至於還可能是，
% cue 660 2229.120 2230.120
\hypertarget{w02:cue:660}{}也許會是必要的。
% cue 661 2230.120 2232.120
\hypertarget{w02:cue:661}{}It is a matter of political morality，
% cue 662 2232.120 2235.370
\hypertarget{w02:cue:662}{}基於政治道德而言，
% cue 663 2235.370 2237.710
\hypertarget{w02:cue:663}{}它可能還是必要的。
% cue 664 2237.710 2240.060
\hypertarget{w02:cue:664}{}不僅是允許有違憲審查，
% cue 665 2240.060 2243.380
\hypertarget{w02:cue:665}{}甚至還非得有違憲審查不可，
% cue 666 2243.380 2244.380
\hypertarget{w02:cue:666}{}所以才叫 required。

\tc{0:37:24}
% cue 667 2244.380 2247.630
\hypertarget{w02:cue:667}{}它這句話講的很，
% cue 668 2247.630 2249.290
\hypertarget{w02:cue:668}{}英文來講沒有，
% cue 669 2249.290 2250.460
\hypertarget{w02:cue:669}{}不能懂，
% cue 670 2250.460 2251.650
\hypertarget{w02:cue:670}{}中文很難懂。
% cue 671 2251.650 2253.850
\hypertarget{w02:cue:671}{}因為我們中國人的頭腦沒有那麼複雜。
% cue 672 2253.850 2255.460
\hypertarget{w02:cue:672}{}不知道我是不是中國人。
% cue 673 2255.460 2256.460
\hypertarget{w02:cue:673}{}我的頭腦，
% cue 674 2256.460 2257.460
\hypertarget{w02:cue:674}{}我假如代表中國人頭腦，
% cue 675 2257.460 2259.460
\hypertarget{w02:cue:675}{}我的頭腦看到這樣的句子，
% cue 676 2259.460 2260.460
\hypertarget{w02:cue:676}{}我是可以懂的。
% cue 677 2260.460 2261.460
\hypertarget{w02:cue:677}{}但是我認為，
% cue 678 2261.460 2263.390
\hypertarget{w02:cue:678}{}一般人中國人是聽不懂的。
% cue 679 2263.390 2264.620
\hypertarget{w02:cue:679}{}就它講的話，
% cue 680 2264.620 2265.620
\hypertarget{w02:cue:680}{}已經好幾個層次了。
% cue 681 2265.620 2266.780
\hypertarget{w02:cue:681}{}好啦。

\tc{0:37:46}
% cue 682 2266.780 2267.940
\hypertarget{w02:cue:682}{}One might think of this question
% cue 683 2267.940 2269.940
\hypertarget{w02:cue:683}{}become academic in countries
% cue 684 2269.940 2272.940
\hypertarget{w02:cue:684}{}where judicial review established as constitutional practice。
% cue 685 2272.940 2275.170
\hypertarget{w02:cue:685}{}也許有人會說，
% cue 686 2275.170 2278.340
\hypertarget{w02:cue:686}{}那你區別特定正當性跟一般正當性，
% cue 687 2278.340 2279.340
\hypertarget{w02:cue:687}{}這有意義嗎？
% cue 688 2279.340 2281.340
\hypertarget{w02:cue:688}{}那純粹是學者之見。
% cue 689 2281.340 2282.690
\hypertarget{w02:cue:689}{}他說呢，
% cue 690 2282.690 2283.690
\hypertarget{w02:cue:690}{}尤其是在違憲審查，
% cue 691 2283.690 2286.750
\hypertarget{w02:cue:691}{}在某些國家已經被認為是一個，
% cue 692 2286.750 2288.730
\hypertarget{w02:cue:692}{}既成的憲政事實，
% cue 693 2288.730 2290.040
\hypertarget{w02:cue:693}{}憲法運作的時候，
% cue 694 2290.040 2292.040
\hypertarget{w02:cue:694}{}這樣講還有意義嗎？
% cue 695 2292.040 2294.040
\hypertarget{w02:cue:695}{}他說典型的例子像 Michael Perry\footnote{Michael J. Perry，美國憲法學者，2003 年起任 Emory 大學 Robert W. Woodruff 講座，現為該校法學院榮譽教授，研究憲法權利、人權理論及法律與宗教。本文第 526 頁引其 What is 'the Constitution'? (and Other Fundamental Questions)，收於 L. Alexander 編 Constitutionalism: Philosophical Foundations（Cambridge University Press, 1998）第 120 頁：違憲審查幾乎自建國之初即為美國憲政秩序的基石，對美國人而言，仍待回答的問題（live questions）在於這項權力應如何行使。\url{https://law.emory.edu/faculty/faculty-emeritus/perry-profile.html}。}。
% cue 696 2294.040 2296.040
\hypertarget{w02:cue:696}{}Michael Perry 我們以前有讀過他的文章，
% cue 697 2296.040 2298.000
\hypertarget{w02:cue:697}{}那他就有這樣的論述，
% cue 698 2298.000 2299.000
\hypertarget{w02:cue:698}{}被他挑出來了。
% cue 699 2299.000 2300.000
\hypertarget{w02:cue:699}{}然後他說，
% cue 700 2300.000 2302.000
\hypertarget{w02:cue:700}{}加拿大還有兩個人\footnote{本文第 526 頁所引為加拿大憲法學者 Patrick Monahan 與 Andrew Petter 合著的 Developments in Constitutional Law: The 1985-86 Term，刊於 (1987) 9 Supreme Court Law Review 69，引文在第 74 頁。兩人認為要不要有違憲審查的「big question」已因 1982 年制定 Charter 而解決，未解決的是違憲審查在加拿大的內在界限。Monahan 自 1982 年起任教於 York 大學 Osgoode Hall 法學院，2003 年 7 月起任院長。}也是這樣的主張。

\tc{0:38:22}
% cue 701 2302.000 2304.000
\hypertarget{w02:cue:701}{}中間這你們自己幹，
% cue 702 2304.000 2304.990
\hypertarget{w02:cue:702}{}不重要。
% cue 703 2304.990 2306.190
\hypertarget{w02:cue:703}{}注意的是他不同意，
% cue 704 2306.190 2308.380
\hypertarget{w02:cue:704}{}這些位置是不錯的，
% cue 705 2308.380 2309.440
\hypertarget{w02:cue:705}{}有沒有？
% cue 706 2309.440 2310.490
\hypertarget{w02:cue:706}{}跳跳跳。

\tc{0:38:30}
% cue 707 2310.490 2311.560
\hypertarget{w02:cue:707}{}文字下面的大行，
% cue 708 2311.560 2314.700
\hypertarget{w02:cue:708}{}526 頁最後一大段。
% cue 709 2314.700 2316.140
\hypertarget{w02:cue:709}{}他說這個觀點，
% cue 710 2316.140 2317.310
\hypertarget{w02:cue:710}{}這個立場，
% cue 711 2317.310 2318.420
\hypertarget{w02:cue:711}{}是誤導的。
% cue 712 2318.420 2320.760
\hypertarget{w02:cue:712}{}這兩個問題是，
% cue 713 2320.760 2322.980
\hypertarget{w02:cue:713}{}這兩個問題是，
% cue 714 2322.980 2324.100
\hypertarget{w02:cue:714}{}違憲審查，
% cue 715 2324.100 2325.100
\hypertarget{w02:cue:715}{}應如何進行，
% cue 716 2325.100 2327.350
\hypertarget{w02:cue:716}{}what are its inherent limitations，
% cue 717 2327.350 2329.350
\hypertarget{w02:cue:717}{}它有沒有本質上的限制，
% cue 718 2329.350 2331.570
\hypertarget{w02:cue:718}{}depends on the answer to the big questions，
% cue 719 2331.570 2333.570
\hypertarget{w02:cue:719}{}這兩個問題，
% cue 720 2333.570 2335.630
\hypertarget{w02:cue:720}{}要取決一個更大的問題，
% cue 721 2335.630 2337.580
\hypertarget{w02:cue:721}{}更大的問題是什麼？
% cue 722 2337.580 2340.290
\hypertarget{w02:cue:722}{}whether judicial review can be made legitimate at all，
% cue 723 2340.290 2341.510
\hypertarget{w02:cue:723}{}根本的問題是，
% cue 724 2341.510 2345.380
\hypertarget{w02:cue:724}{}違憲審查能不能成為正當。

\tc{0:39:05}
% cue 725 2345.380 2348.600
\hypertarget{w02:cue:725}{}違憲審查如果不能成為正當，
% cue 726 2348.600 2352.600
\hypertarget{w02:cue:726}{}你討論違憲審查應如何行事，
% cue 727 2352.600 2354.570
\hypertarget{w02:cue:727}{}你違憲審查有沒有本質限制，
% cue 728 2354.570 2357.340
\hypertarget{w02:cue:728}{}就好像是太逾越了，
% cue 729 2357.340 2360.190
\hypertarget{w02:cue:729}{}太跳躍，太超過了，
% cue 730 2360.190 2361.190
\hypertarget{w02:cue:730}{}你一步一步來。
% cue 731 2361.190 2363.250
\hypertarget{w02:cue:731}{}如果現成它都不能正名為正當，
% cue 732 2363.250 2365.250
\hypertarget{w02:cue:732}{}你還跟我討論要如何行事，
% cue 733 2365.250 2367.530
\hypertarget{w02:cue:733}{}它有什麼本質限制，
% cue 734 2367.530 2369.660
\hypertarget{w02:cue:734}{}你講到最後一步不來。

\tc{0:39:29}
% cue 735 2369.660 2373.750
\hypertarget{w02:cue:735}{}whether it is a good idea，
% cue 736 2373.750 2378.720
\hypertarget{w02:cue:736}{}and whether it has been historically established by we the people\footnote{美國聯邦憲法前言首句：「We the People of the United States, in Order to form a more perfect Union, ... do ordain and establish this Constitution for the United States of America.」（我們合眾國人民，為建立更完善的聯邦……制定並確立這部美利堅合眾國憲法）。本文第 526 頁所問違憲審查是否「historically established by We the people」，即用此語；Ackerman 的三卷本 We the People 亦以此語立題。\url{https://www.archives.gov/founding-docs/constitution-transcript}。}，
% cue 737 2378.720 2380.690
\hypertarget{w02:cue:737}{}我們人民到底有沒有承認這個制度。
% cue 738 2380.690 2382.690
\hypertarget{w02:cue:738}{}we the people 是美國憲法的用語，
% cue 739 2382.690 2385.580
\hypertarget{w02:cue:739}{}人民立憲的時候有沒有這個想法。

\tc{0:39:45}
% cue 740 2385.580 2390.990
\hypertarget{w02:cue:740}{}where the institution of the judicial review is accepted as legitimate
% cue 741 2390.990 2395.580
\hypertarget{w02:cue:741}{}the determination of its proper limitations
% cue 742 2395.580 2398.640
\hypertarget{w02:cue:742}{}including the correct modes of interpretation and deference
% cue 743 2398.640 2406.020
\hypertarget{w02:cue:743}{}is constrained by the value and considerations that have justified it in the first place。

\tc{0:40:06}
% cue 744 2406.020 2409.110
\hypertarget{w02:cue:744}{}好，這幾個地方要稍微注意一下，
% cue 745 2409.110 2411.070
\hypertarget{w02:cue:745}{}等一下你會看到了。
% cue 746 2411.070 2413.430
\hypertarget{w02:cue:746}{}它講了一些根本，根本是什麼？
% cue 747 2413.430 2415.590
\hypertarget{w02:cue:747}{}就是價值還有考慮，
% cue 748 2415.590 2417.750
\hypertarget{w02:cue:748}{}那種價值跟考慮呢？
% cue 749 2417.750 2420.840
\hypertarget{w02:cue:749}{}是使得違憲審查變成正當的。
% cue 750 2420.840 2423.480
\hypertarget{w02:cue:750}{}好，所以真正重要的是，
% cue 751 2423.480 2425.600
\hypertarget{w02:cue:751}{}憲法裡面也好，
% cue 752 2425.600 2427.940
\hypertarget{w02:cue:752}{}不管是成文的，
% cue 753 2427.940 2428.910
\hypertarget{w02:cue:753}{}或者是不成文的，
% cue 754 2428.910 2429.910
\hypertarget{w02:cue:754}{}外面也好，
% cue 755 2429.910 2436.550
\hypertarget{w02:cue:755}{}到底有沒有使得違憲審查成為正當的這種價值和考慮？
% cue 756 2436.550 2438.220
\hypertarget{w02:cue:756}{}that is，
% cue 757 2438.220 2439.180
\hypertarget{w02:cue:757}{}也就是說，
% cue 758 2439.180 2443.490
\hypertarget{w02:cue:758}{}by the answer given to the more abstract
% cue 759 2443.490 2446.490
\hypertarget{w02:cue:759}{}the normative and speculative question of legitimacy。

\tc{0:40:46}
% cue 760 2446.490 2449.130
\hypertarget{w02:cue:760}{}好了，中間跳過。
% cue 761 2449.130 2451.290
\hypertarget{w02:cue:761}{}it is therefore important to determine，
% cue 762 2451.290 2453.730
\hypertarget{w02:cue:762}{}很重要的要決定，
% cue 763 2453.730 2463.400
\hypertarget{w02:cue:763}{}whether judicial review of democratically enacted legislation can be made morally permissible in principle。
% cue 764 2463.400 2465.720
\hypertarget{w02:cue:764}{}剛才講的是 legitimacy，
% cue 765 2465.720 2467.850
\hypertarget{w02:cue:765}{}legitimacy 它現在換一個字，
% cue 766 2467.850 2469.960
\hypertarget{w02:cue:766}{}叫做 morally permissible，
% cue 767 2469.960 2471.850
\hypertarget{w02:cue:767}{}he，
% cue 768 2471.850 2474.040
\hypertarget{w02:cue:768}{}道德上是可接受的，
% cue 769 2474.040 2475.070
\hypertarget{w02:cue:769}{}原則上的原則。

\section{全文三節的論證安排}

\tc{0:41:15}
% cue 770 2475.070 2480.000
\hypertarget{w02:cue:770}{}my argument is divided into three sections.
% cue 771 2480.000 2482.350
\hypertarget{w02:cue:771}{}它下面要講的就是這件事情，
% cue 772 2482.350 2486.180
\hypertarget{w02:cue:772}{}我要怎麼論證在一個民主國家，
% cue 773 2486.180 2490.290
\hypertarget{w02:cue:773}{}對一個民主程序所通過的法律，
% cue 774 2490.290 2496.600
\hypertarget{w02:cue:774}{}我們由不民主不民選的一群法官可以來進行審查，
% cue 775 2496.600 2502.140
\hypertarget{w02:cue:775}{}這件事情在原則上去道德上允許。
% cue 776 2502.140 2505.380
\hypertarget{w02:cue:776}{}這就是它講的 legitimacy，
% cue 777 2505.380 2507.380
\hypertarget{w02:cue:777}{}legitimacy 就是 morally permissible。
% cue 778 2507.380 2509.370
\hypertarget{w02:cue:778}{}你以後去注意，
% cue 779 2509.370 2510.370
\hypertarget{w02:cue:779}{}我現在講的，
% cue 780 2510.370 2511.370
\hypertarget{w02:cue:780}{}這就是你讀了以後，
% cue 781 2511.370 2513.370
\hypertarget{w02:cue:781}{}你要把它寫在旁邊，
% cue 782 2513.370 2516.730
\hypertarget{w02:cue:782}{}這兩個是同位常常會互用的，
% cue 783 2516.730 2518.550
\hypertarget{w02:cue:783}{}interchangeable。

\tc{0:41:58}
% cue 784 2518.550 2519.870
\hypertarget{w02:cue:784}{}好，那它下面呢，
% cue 785 2519.870 2521.870
\hypertarget{w02:cue:785}{}分成三個章節，
% cue 786 2521.870 2522.930
\hypertarget{w02:cue:786}{}section 1，
% cue 787 2522.930 2527.340
\hypertarget{w02:cue:787}{}I will examine the structure and the limits of the theories，
% cue 788 2527.340 2528.340
\hypertarget{w02:cue:788}{}好多個理論，
% cue 789 2528.340 2535.610
\hypertarget{w02:cue:789}{}seeking to establish specific legitimacy of judicial review，
% cue 790 2535.610 2537.310
\hypertarget{w02:cue:790}{}好幾個不同的理論，
% cue 791 2537.310 2541.530
\hypertarget{w02:cue:791}{}要來證明所謂違憲審查特定正當性，
% cue 792 2541.530 2543.620
\hypertarget{w02:cue:792}{}所以特定正當性是一節。
% cue 793 2543.620 2544.710
\hypertarget{w02:cue:793}{}剛開標題，
% cue 794 2544.710 2545.710
\hypertarget{w02:cue:794}{}是不是這樣子，
% cue 795 2545.710 2547.060
\hypertarget{w02:cue:795}{}好，section 2，
% cue 796 2547.060 2551.530
\hypertarget{w02:cue:796}{}I will establish the general legitimacy of judicial review，
% cue 797 2551.530 2554.620
\hypertarget{w02:cue:797}{}我認為我要講的是一般正當性，
% cue 798 2554.620 2555.810
\hypertarget{w02:cue:798}{}這是第二節。

\tc{0:42:35}
% cue 799 2555.810 2557.810
\hypertarget{w02:cue:799}{}好，那為什麼跑出來第三節，
% cue 800 2557.810 2559.810
\hypertarget{w02:cue:800}{}在這種考驗呢，
% cue 801 2559.810 2560.810
\hypertarget{w02:cue:801}{}這句話要看一下。
% cue 802 2560.810 2561.740
\hypertarget{w02:cue:802}{}第三節，
% cue 803 2561.740 2562.870
\hypertarget{w02:cue:803}{}我主要要論，
% cue 804 2562.870 2563.870
\hypertarget{w02:cue:804}{}我剛才提的，
% cue 805 2563.870 2565.880
\hypertarget{w02:cue:805}{}違憲審查，
% cue 806 2565.880 2566.880
\hypertarget{w02:cue:806}{}一般正當性的這個命題，
% cue 807 2566.880 2567.880
\hypertarget{w02:cue:807}{}或者這個理論，
% cue 808 2567.880 2570.170
\hypertarget{w02:cue:808}{}它可能對我們的代議制度，
% cue 809 2570.170 2572.480
\hypertarget{w02:cue:809}{}對我們憲法基礎的這個代議制度，
% cue 810 2572.480 2573.480
\hypertarget{w02:cue:810}{}可能產生什麼影響？
% cue 811 2573.480 2575.750
\hypertarget{w02:cue:811}{}你說違憲審查是正當的，
% cue 812 2575.750 2576.730
\hypertarget{w02:cue:812}{}而且原則上是道德允許的，
% cue 813 2576.730 2576.830
\hypertarget{w02:cue:813}{}好，那你這麼講，
% cue 814 2576.730 2578.330
\hypertarget{w02:cue:814}{}對於我們民主憲法，
% cue 815 2578.330 2579.550
\hypertarget{w02:cue:815}{}它的基本原則，
% cue 816 2579.550 2580.860
\hypertarget{w02:cue:816}{}會產生什麼影響呢？
% cue 817 2580.860 2583.110
\hypertarget{w02:cue:817}{}這就是我第三節要講的。
% cue 818 2583.110 2584.210
\hypertarget{w02:cue:818}{}所以他會提出來一個，
% cue 819 2584.210 2585.210
\hypertarget{w02:cue:819}{}Constitutional Law，
% cue 820 2585.210 2586.210
\hypertarget{w02:cue:820}{}就是要論，
% cue 821 2586.210 2587.210
\hypertarget{w02:cue:821}{}我剛才提的，
% cue 822 2587.210 2588.130
\hypertarget{w02:cue:822}{}我剛才提的，
% cue 823 2588.130 2589.620
\hypertarget{w02:cue:823}{}違憲審查，
% cue 824 2589.620 2590.800 刪：與 cue 806 同文（連續 12 個重複解碼）
% cue 825 2590.800 2591.800 刪：與 cue 807 同文（連續 12 個重複解碼）
% cue 826 2591.800 2593.800 刪：與 cue 808 同文（連續 12 個重複解碼）
% cue 827 2593.800 2595.800 刪：與 cue 809 同文（連續 12 個重複解碼）
% cue 828 2595.800 2600.640 刪：與 cue 810 同文（連續 12 個重複解碼）
% cue 829 2600.640 2602.960 刪：與 cue 811 同文（連續 12 個重複解碼）
% cue 830 2602.960 2605.630 刪：與 cue 812 同文（連續 12 個重複解碼）
% cue 831 2605.630 2608.350 刪：與 cue 813 同文（連續 12 個重複解碼）
% cue 832 2608.350 2610.730 刪：與 cue 814 同文（連續 12 個重複解碼）
% cue 833 2610.730 2612.680 刪：與 cue 815 同文（連續 12 個重複解碼）
% cue 834 2612.680 2614.060 刪：與 cue 816 同文（連續 12 個重複解碼）
% cue 835 2614.060 2616.410 刪：與 cue 817 同文（連續 12 個重複解碼）

\tc{0:43:36}
% cue 836 2616.410 2618.410
\hypertarget{w02:cue:836}{}你也可以把它想成是，
% cue 837 2618.410 2621.010
\hypertarget{w02:cue:837}{}Democratic Constitution Thesis，
% cue 838 2621.010 2623.010
\hypertarget{w02:cue:838}{}民主憲法的基礎，
% cue 839 2623.010 2623.970
\hypertarget{w02:cue:839}{}命題。
% cue 840 2623.970 2624.970
\hypertarget{w02:cue:840}{}OK。

\tc{0:43:44}
% cue 841 2624.970 2625.920
\hypertarget{w02:cue:841}{}In a short conclusion，
% cue 842 2625.920 2629.360
\hypertarget{w02:cue:842}{}最後一個很簡短的結論，
% cue 843 2629.360 2630.840
\hypertarget{w02:cue:843}{}我建議，
% cue 844 2630.840 2632.990
\hypertarget{w02:cue:844}{}一些答案到三個其他問題，
% cue 845 2632.990 2635.120
\hypertarget{w02:cue:845}{}更多的其他問題，
% cue 846 2635.120 2637.150
\hypertarget{w02:cue:846}{}那些主憲法的現實，
% cue 847 2637.150 2639.560
\hypertarget{w02:cue:847}{}對於主宰的規範性，
% cue 848 2639.560 2641.590
\hypertarget{w02:cue:848}{}當代的憲政理論，
% cue 849 2641.590 2643.840
\hypertarget{w02:cue:849}{}影響的另外三個問題，
% cue 850 2643.840 2645.840
\hypertarget{w02:cue:850}{}我想的也會提供，
% cue 851 2645.840 2647.840
\hypertarget{w02:cue:851}{}某些答案。
% cue 852 2647.840 2649.870
\hypertarget{w02:cue:852}{}What is the purpose of the Constitution?
% cue 853 2649.870 2651.960
\hypertarget{w02:cue:853}{}憲法的目的，
% cue 854 2651.960 2652.960
\hypertarget{w02:cue:854}{}有沒有剛才講的，
% cue 855 2652.960 2654.180
\hypertarget{w02:cue:855}{}我們結論第一點，
% cue 856 2654.180 2656.200
\hypertarget{w02:cue:856}{}就是 the purpose of the Constitution。
% cue 857 2656.200 2657.200
\hypertarget{w02:cue:857}{}第二點，
% cue 858 2657.200 2659.420
\hypertarget{w02:cue:858}{}what makes the Constitution legitimate?
% cue 859 2659.420 2661.420
\hypertarget{w02:cue:859}{}什麼使得憲法變成正當？
% cue 860 2661.420 2662.540
\hypertarget{w02:cue:860}{}第三點，
% cue 861 2662.540 2663.540
\hypertarget{w02:cue:861}{}what kind of argument are legitimate
% cue 862 2663.540 2665.860
\hypertarget{w02:cue:862}{}within the process of the Constitutional Interpretation?
% cue 863 2665.860 2667.860
\hypertarget{w02:cue:863}{}在憲法解釋的過程當中，
% cue 864 2667.860 2669.860
\hypertarget{w02:cue:864}{}哪一些論證使得它變得正當？

\tc{0:44:29}
% cue 865 2669.860 2670.860
\hypertarget{w02:cue:865}{}我們剛剛講 123，
% cue 866 2670.860 2671.860
\hypertarget{w02:cue:866}{}是不是正 123?
% cue 867 2671.860 2672.860
\hypertarget{w02:cue:867}{}所以你這樣看，
% cue 868 2672.860 2674.860
\hypertarget{w02:cue:868}{}你如果覺得剛才講的不清楚，
% cue 869 2674.860 2677.110
\hypertarget{w02:cue:869}{}結論我實際上是可能再重複了。
% cue 870 2677.110 2679.880
\hypertarget{w02:cue:870}{}我今天用 fade 也飛不到那裡。
% cue 871 2679.880 2681.880
\hypertarget{w02:cue:871}{}但是我剛已經講過了，
% cue 872 2681.880 2683.880
\hypertarget{w02:cue:872}{}你現在就知道它的結構。

\section{難讀文獻的讀法}

\tc{0:44:43}
% cue 873 2683.880 2685.940
\hypertarget{w02:cue:873}{}好，我現在回應一個問題，
% cue 874 2685.940 2687.040
\hypertarget{w02:cue:874}{}上次我學說，
% cue 875 2687.040 2689.250
\hypertarget{w02:cue:875}{}老師這麼難的東西要怎麼讀？
% cue 876 2689.250 2690.220
\hypertarget{w02:cue:876}{}我跟這個問題，
% cue 877 2690.220 2693.340
\hypertarget{w02:cue:877}{}我跟你們一樣，
% cue 878 2693.340 2694.340
\hypertarget{w02:cue:878}{}老師也不是神，
% cue 879 2694.340 2695.340
\hypertarget{w02:cue:879}{}我也是在摸索，
% cue 880 2695.340 2698.340
\hypertarget{w02:cue:880}{}我只比你摸索多 40 年而已。
% cue 881 2698.340 2701.460
\hypertarget{w02:cue:881}{}那我就是跟你們一樣讀。
% cue 882 2701.460 2702.460
\hypertarget{w02:cue:882}{}那你現在看到了，
% cue 883 2702.460 2704.460
\hypertarget{w02:cue:883}{}你讀讀頭，讀讀尾巴，
% cue 884 2704.460 2706.650
\hypertarget{w02:cue:884}{}你已經大概知道它的整體結構。

\tc{0:45:06}
% cue 885 2706.650 2707.620
\hypertarget{w02:cue:885}{}那現在呢，
% cue 886 2707.620 2708.780
\hypertarget{w02:cue:886}{}看你的需要，
% cue 887 2708.780 2712.040
\hypertarget{w02:cue:887}{}那我認為我們這個課快不得了，
% cue 888 2712.040 2713.040
\hypertarget{w02:cue:888}{}快不得了。
% cue 889 2713.040 2715.290
\hypertarget{w02:cue:889}{}有關於那個 Constitutional Thesis，
% cue 890 2715.290 2717.540
\hypertarget{w02:cue:890}{}Constitutional Law Thesis，
% cue 891 2717.540 2718.540
\hypertarget{w02:cue:891}{}今天一定講不到，
% cue 892 2718.540 2720.600
\hypertarget{w02:cue:892}{}但是你們回去看。
% cue 893 2720.600 2722.760
\hypertarget{w02:cue:893}{}然後它裡面又分成成文憲法，
% cue 894 2722.760 2723.760
\hypertarget{w02:cue:894}{}跟非成文憲法。
% cue 895 2723.760 2725.760
\hypertarget{w02:cue:895}{}那個你就算是進階版，
% cue 896 2725.760 2727.260
\hypertarget{w02:cue:896}{}就是說，
% cue 897 2727.260 2728.260
\hypertarget{w02:cue:897}{}你的學問要很好，
% cue 898 2728.260 2729.450
\hypertarget{w02:cue:898}{}你才需要去處理。

\tc{0:45:29}
% cue 899 2729.450 2730.450
\hypertarget{w02:cue:899}{}現在我們能處理一跟二，
% cue 900 2730.450 2732.630
\hypertarget{w02:cue:900}{}就已經 Orbit of War。
% cue 901 2732.630 2733.770
\hypertarget{w02:cue:901}{}你一跟二，
% cue 902 2733.770 2734.770
\hypertarget{w02:cue:902}{}你就看到一大堆人名字在裡，
% cue 903 2734.770 2736.240
\hypertarget{w02:cue:903}{}不管是 John E.P.
% cue 904 2736.240 2737.240
\hypertarget{w02:cue:904}{}還有包括這個，
% cue 905 2737.240 2739.240
\hypertarget{w02:cue:905}{}Bruce Ackerman 都在裡面。

\tc{0:45:39}
% cue 906 2739.240 2740.240
\hypertarget{w02:cue:906}{}所以我跟你說，
% cue 907 2740.240 2741.390
\hypertarget{w02:cue:907}{}這是一個很難的討論。
% cue 908 2741.390 2742.370
\hypertarget{w02:cue:908}{}他就把現在我們討論，
% cue 909 2742.370 2744.510
\hypertarget{w02:cue:909}{}這麼大堆複雜的事情，
% cue 910 2744.510 2746.600
\hypertarget{w02:cue:910}{}它統統包在這篇文章。
% cue 911 2746.600 2747.820
\hypertarget{w02:cue:911}{}那這篇文章，
% cue 912 2747.820 2748.820
\hypertarget{w02:cue:912}{}當然很安心。
% cue 913 2748.820 2749.820
\hypertarget{w02:cue:913}{}我是很佩服這個人。
% cue 914 2749.820 2750.820
\hypertarget{w02:cue:914}{}要是我寫了，
% cue 915 2750.820 2751.820
\hypertarget{w02:cue:915}{}我坦誠，
% cue 916 2751.820 2752.820
\hypertarget{w02:cue:916}{}我寫不出來。
% cue 917 2752.820 2754.110
\hypertarget{w02:cue:917}{}好。

\section{特定正當性命題與政治正當性}

\tc{0:45:54}
% cue 918 2754.110 2755.110
\hypertarget{w02:cue:918}{}來，先看一。
% cue 919 2755.110 2758.030
\hypertarget{w02:cue:919}{}Specific legitimacy thesis，
% cue 920 2758.030 2760.370
\hypertarget{w02:cue:920}{}特定正當性的命題。
% cue 921 2760.370 2762.370
\hypertarget{w02:cue:921}{}The question of political legitimacy，
% cue 922 2762.370 2766.060
\hypertarget{w02:cue:922}{}先講什麼叫做政治正當性。
% cue 923 2766.060 2768.400
\hypertarget{w02:cue:923}{}我們前面都講 legitimacy，
% cue 924 2768.400 2770.500
\hypertarget{w02:cue:924}{}從來沒有加 political。

\tc{0:46:10}
% cue 925 2770.500 2771.790
\hypertarget{w02:cue:925}{}所以你現在看到這個作者，
% cue 926 2771.790 2772.700
\hypertarget{w02:cue:926}{}它實際上是，
% cue 927 2772.700 2774.820
\hypertarget{w02:cue:927}{}你可以說它很 tricky，
% cue 928 2774.820 2777.130
\hypertarget{w02:cue:928}{}你每次偷渡一點，
% cue 929 2777.130 2778.130
\hypertarget{w02:cue:929}{}每次偷渡一點，
% cue 930 2778.130 2781.450
\hypertarget{w02:cue:930}{}這次偷渡進來的是 political legitimacy。
% cue 931 2781.450 2782.640
\hypertarget{w02:cue:931}{}因為他認為，
% cue 932 2782.640 2785.020
\hypertarget{w02:cue:932}{}政府正當是政治問題。
% cue 933 2785.020 2787.200
\hypertarget{w02:cue:933}{}有關於政治正當性的問題，
% cue 934 2787.200 2789.650
\hypertarget{w02:cue:934}{}這是政治正當性的問題，
% cue 935 2789.650 2791.250
\hypertarget{w02:cue:935}{}這是人生的問題。
% cue 936 2791.250 2792.250
\hypertarget{w02:cue:936}{}當一個人或一個人，
% cue 937 2792.250 2793.210
\hypertarget{w02:cue:937}{}當一個人或一個人，
% cue 938 2793.210 2794.280
\hypertarget{w02:cue:938}{}擔任的公司，
% cue 939 2794.280 2795.280
\hypertarget{w02:cue:939}{}擔任的公司，
% cue 940 2795.280 2796.630
\hypertarget{w02:cue:940}{}要有權利益。

\tc{0:46:36}
% cue 941 2796.630 2797.780
\hypertarget{w02:cue:941}{}注意這句話，
% cue 942 2797.780 2798.780
\hypertarget{w02:cue:942}{}當一個人或一個機構，
% cue 943 2798.780 2800.780
\hypertarget{w02:cue:943}{}對個人主張權威的時候，
% cue 944 2800.780 2804.830
\hypertarget{w02:cue:944}{}它就會產生政治正當性的問題。
% cue 945 2804.830 2808.140
\hypertarget{w02:cue:945}{}後面那個 or，
% cue 946 2808.140 2809.140
\hypertarget{w02:cue:946}{}可以看，
% cue 947 2809.140 2811.580
\hypertarget{w02:cue:947}{}你可以不看。
% cue 948 2811.580 2812.580
\hypertarget{w02:cue:948}{}就是你到這裡就可以懂。
% cue 949 2812.580 2814.580
\hypertarget{w02:cue:949}{}後面這是換一種方法來表述。

\tc{0:46:54}
% cue 950 2814.580 2815.930
\hypertarget{w02:cue:950}{}再講一遍，
% cue 951 2815.930 2816.990
\hypertarget{w02:cue:951}{}政治正當性，
% cue 952 2816.990 2819.020
\hypertarget{w02:cue:952}{}就是當一個人或一個機構，
% cue 953 2819.020 2821.020
\hypertarget{w02:cue:953}{}對個人主張權威的時候，
% cue 954 2821.020 2823.180
\hypertarget{w02:cue:954}{}你得聽我的，
% cue 955 2823.180 2824.210
\hypertarget{w02:cue:955}{}不然會處罰你。
% cue 956 2824.210 2825.340
\hypertarget{w02:cue:956}{}你憑什麼？
% cue 957 2825.340 2826.780
\hypertarget{w02:cue:957}{}我有正當性。
% cue 958 2826.780 2827.780
\hypertarget{w02:cue:958}{}你有什麼正當性？
% cue 959 2827.780 2829.420
\hypertarget{w02:cue:959}{}我有政治正當性。
% cue 960 2829.420 2830.420
\hypertarget{w02:cue:960}{}所謂的正當性，
% cue 961 2830.420 2831.420
\hypertarget{w02:cue:961}{}這是政治正當性。

\tc{0:47:11}
% cue 962 2831.420 2832.830
\hypertarget{w02:cue:962}{}好了。
% cue 963 2832.830 2833.860
\hypertarget{w02:cue:963}{}再講一下，
% cue 964 2833.860 2834.920
\hypertarget{w02:cue:964}{}這個問題，
% cue 965 2834.920 2835.820
\hypertarget{w02:cue:965}{}這個問題，
% cue 966 2835.820 2836.910
\hypertarget{w02:cue:966}{}seeks to determine，
% cue 967 2836.910 2837.940
\hypertarget{w02:cue:967}{}尋求決定，
% cue 968 2837.940 2839.220
\hypertarget{w02:cue:968}{}尋求斷定，
% cue 969 2839.220 2841.280
\hypertarget{w02:cue:969}{}the moral considerations，
% cue 970 2841.280 2843.980
\hypertarget{w02:cue:970}{}所以你看呢，
% cue 971 2843.980 2844.980
\hypertarget{w02:cue:971}{}他講的這些事情，
% cue 972 2844.980 2845.950
\hypertarget{w02:cue:972}{}都是道德，
% cue 973 2845.950 2848.210
\hypertarget{w02:cue:973}{}道德上的考慮。
% cue 974 2848.210 2850.430
\hypertarget{w02:cue:974}{}什麼樣的道德考慮？
% cue 975 2850.430 2851.490
\hypertarget{w02:cue:975}{}by virtue of which，
% cue 976 2851.490 2853.490
\hypertarget{w02:cue:976}{}經由那個道德考慮，
% cue 977 2853.490 2854.450
\hypertarget{w02:cue:977}{}this person
% cue 978 2854.450 2855.450
\hypertarget{w02:cue:978}{}or this institution
% cue 979 2855.450 2857.450
\hypertarget{w02:cue:979}{}is entitled to command。

\tc{0:47:37}
% cue 980 2857.450 2858.450
\hypertarget{w02:cue:980}{}我是個個人，
% cue 981 2858.450 2859.830
\hypertarget{w02:cue:981}{}我是個機構，
% cue 982 2859.830 2861.820
\hypertarget{w02:cue:982}{}我對於其他的個人，
% cue 983 2861.820 2862.820
\hypertarget{w02:cue:983}{}主張權威。
% cue 984 2862.820 2864.620
\hypertarget{w02:cue:984}{}那我說，
% cue 985 2864.620 2866.620
\hypertarget{w02:cue:985}{}這裡牽涉到政治正當性。
% cue 986 2866.620 2867.780
\hypertarget{w02:cue:986}{}那換個方法講，
% cue 987 2867.780 2868.740
\hypertarget{w02:cue:987}{}政治正當性呢，
% cue 988 2868.740 2872.120
\hypertarget{w02:cue:988}{}就是來詢問，
% cue 989 2872.120 2874.080
\hypertarget{w02:cue:989}{}你政治正當性裡面，
% cue 990 2874.080 2876.190
\hypertarget{w02:cue:990}{}到底有什麼，
% cue 991 2876.190 2878.290
\hypertarget{w02:cue:991}{}moral considerations，
% cue 992 2878.290 2880.510
\hypertarget{w02:cue:992}{}有什麼樣的道德考慮，
% cue 993 2880.510 2882.570
\hypertarget{w02:cue:993}{}使得你有權力這麼主張，
% cue 994 2882.570 2885.270
\hypertarget{w02:cue:994}{}你得服從我。

\tc{0:48:05}
% cue 995 2885.270 2886.710
\hypertarget{w02:cue:995}{}你問我說，
% cue 996 2886.710 2887.710
\hypertarget{w02:cue:996}{}你跟我講，
% cue 997 2887.710 2888.710
\hypertarget{w02:cue:997}{}你要服從我，
% cue 998 2888.710 2889.800
\hypertarget{w02:cue:998}{}那我問你，
% cue 999 2889.800 2890.800
\hypertarget{w02:cue:999}{}為什麼我要服從你？
% cue 1000 2890.800 2892.370
\hypertarget{w02:cue:1000}{}你跟我講說，
% cue 1001 2892.370 2893.370
\hypertarget{w02:cue:1001}{}我具有道德正當性。

\tc{0:48:13}
% cue 1002 2893.370 2895.140
\hypertarget{w02:cue:1002}{}好。
% cue 1003 2895.140 2896.140
\hypertarget{w02:cue:1003}{}那什麼叫道德正當性？
% cue 1004 2896.140 2897.300
\hypertarget{w02:cue:1004}{}那就是，
% cue 1005 2897.300 2899.300
\hypertarget{w02:cue:1005}{}使得你必須服從我的那些，
% cue 1006 2899.300 2901.300
\hypertarget{w02:cue:1006}{}道德上的考慮。
% cue 1007 2901.300 2904.030
\hypertarget{w02:cue:1007}{}我這樣講，
% cue 1008 2904.030 2905.380
\hypertarget{w02:cue:1008}{}句子有錯亂嗎？
% cue 1009 2905.380 2907.380
\hypertarget{w02:cue:1009}{}你去換，
% cue 1010 2907.380 2908.350
\hypertarget{w02:cue:1010}{}代換你就會帶。

\tc{0:48:28}
% cue 1011 2908.350 2909.540
\hypertarget{w02:cue:1011}{}這裡開始的時候，
% cue 1012 2909.540 2910.540
\hypertarget{w02:cue:1012}{}是講 legitimacy，
% cue 1013 2910.540 2911.690
\hypertarget{w02:cue:1013}{}legitimacy，
% cue 1014 2911.690 2914.650
\hypertarget{w02:cue:1014}{}其實就是某些道德考慮，
% cue 1015 2914.650 2916.610
\hypertarget{w02:cue:1015}{}那這些道德考慮，
% cue 1016 2916.610 2918.610
\hypertarget{w02:cue:1016}{}可以使得你叫別人，
% cue 1017 2918.610 2920.570
\hypertarget{w02:cue:1017}{}要服從你的權威，
% cue 1018 2920.570 2922.720
\hypertarget{w02:cue:1018}{}就對了。

\tc{0:48:42}
% cue 1019 2922.720 2923.720
\hypertarget{w02:cue:1019}{}那我就記住，
% cue 1020 2923.720 2924.680
\hypertarget{w02:cue:1020}{}legitimacy，
% cue 1021 2924.680 2926.770
\hypertarget{w02:cue:1021}{}就是某些道德考慮。

\tc{0:48:46}
% cue 1022 2926.770 2928.210
\hypertarget{w02:cue:1022}{}然後，
% cue 1023 2928.210 2932.380
\hypertarget{w02:cue:1023}{}this idea must minimally，
% cue 1024 2932.380 2934.410
\hypertarget{w02:cue:1024}{}最起碼的 imply，
% cue 1025 2934.410 2936.460
\hypertarget{w02:cue:1025}{}it is in principle morally wrong
% cue 1026 2936.460 2938.580
\hypertarget{w02:cue:1026}{}for someone to interfere with the legal
% cue 1027 2938.580 2940.580
\hypertarget{w02:cue:1027}{}the legitimate political actions
% cue 1028 2940.580 2941.580
\hypertarget{w02:cue:1028}{}or decisions
% cue 1029 2941.580 2945.530
\hypertarget{w02:cue:1029}{}unless one can show that such interference
% cue 1030 2945.530 2947.560
\hypertarget{w02:cue:1030}{}is itself legitimate。

\tc{0:49:07}
% cue 1031 2947.560 2950.010
\hypertarget{w02:cue:1031}{}除非你能證明，
% cue 1032 2950.010 2951.970
\hypertarget{w02:cue:1032}{}干涉這件事情是正當的，
% cue 1033 2951.970 2953.020
\hypertarget{w02:cue:1033}{}否則，
% cue 1034 2953.020 2954.310
\hypertarget{w02:cue:1034}{}原則上，
% cue 1035 2954.310 2955.370
\hypertarget{w02:cue:1035}{}我們認為，
% cue 1036 2955.370 2956.370
\hypertarget{w02:cue:1036}{}你去干涉，
% cue 1037 2956.370 2959.610
\hypertarget{w02:cue:1037}{}和正當的政治行動，
% cue 1038 2959.610 2960.570
\hypertarget{w02:cue:1038}{}或者政治決定，
% cue 1039 2960.570 2962.570
\hypertarget{w02:cue:1039}{}這件事情本身，
% cue 1040 2962.570 2963.570
\hypertarget{w02:cue:1040}{}是道德上的錯誤。
% cue 1041 2963.570 2966.490
\hypertarget{w02:cue:1041}{}除非你能證明，
% cue 1042 2966.490 2967.490
\hypertarget{w02:cue:1042}{}我干涉是正當的，
% cue 1043 2967.490 2968.490
\hypertarget{w02:cue:1043}{}你證明不出來，
% cue 1044 2968.490 2970.460
\hypertarget{w02:cue:1044}{}干涉是正當的，
% cue 1045 2970.460 2972.460
\hypertarget{w02:cue:1045}{}那你就得必須承認，
% cue 1046 2972.460 2973.460
\hypertarget{w02:cue:1046}{}你的干涉是錯誤的，
% cue 1047 2973.460 2974.460
\hypertarget{w02:cue:1047}{}道德上是錯誤的。

\tc{0:49:34}
% cue 1048 2974.460 2975.610
\hypertarget{w02:cue:1048}{}所以你看到這句，
% cue 1049 2975.610 2977.610
\hypertarget{w02:cue:1049}{}that is morally permissible，
% cue 1050 2977.610 2978.610
\hypertarget{w02:cue:1050}{}legitimate，
% cue 1051 2978.610 2979.670
\hypertarget{w02:cue:1051}{}就是 morally permissible，
% cue 1052 2979.670 2981.670
\hypertarget{w02:cue:1052}{}而可不可以，
% cue 1053 2981.670 2985.210
\hypertarget{w02:cue:1053}{}是要看後面的 moral considerations，
% cue 1054 2985.210 2987.240
\hypertarget{w02:cue:1054}{}來決定。
% cue 1055 2987.240 2991.190
\hypertarget{w02:cue:1055}{}這個就算是明詞寢誓吧。
% cue 1056 2991.190 2993.280
\hypertarget{w02:cue:1056}{}你這樣子玩了，
% cue 1057 2993.280 2995.280
\hypertarget{w02:cue:1057}{}你就可以明白他要講什麼。

\tc{0:49:55}
% cue 1058 2995.280 2996.280
\hypertarget{w02:cue:1058}{}好。
% cue 1059 2996.280 2998.280
\hypertarget{w02:cue:1059}{}那麼他現在講呢，
% cue 1060 2998.280 2999.280
\hypertarget{w02:cue:1060}{}看下一句，
% cue 1061 2999.280 3000.410
\hypertarget{w02:cue:1061}{}下一段，
% cue 1062 3000.410 3001.540
\hypertarget{w02:cue:1062}{}in general
% cue 1063 3001.540 3004.540
\hypertarget{w02:cue:1063}{}the answer to the question has taken a form，
% cue 1064 3004.540 3007.970
\hypertarget{w02:cue:1064}{}用什麼樣的形式，
% cue 1065 3007.970 3009.030
\hypertarget{w02:cue:1065}{}moral principle
% cue 1066 3009.030 3012.700
\hypertarget{w02:cue:1066}{}indicating the positive fact
% cue 1067 3012.700 3014.920
\hypertarget{w02:cue:1067}{}by virtue of which
% cue 1068 3014.920 3017.950
\hypertarget{w02:cue:1068}{}political actions are specifically authorized，
% cue 1069 3017.950 3021.450
\hypertarget{w02:cue:1069}{}你採取了某樣的政治行動，
% cue 1070 3021.450 3024.450
\hypertarget{w02:cue:1070}{}那個政治行動是特定得到了授權，
% cue 1071 3024.450 3029.430
\hypertarget{w02:cue:1071}{}你有一個這樣子的具體的積極事實，
% cue 1072 3029.430 3031.400
\hypertarget{w02:cue:1072}{}有這個具體的積極事實，
% cue 1073 3031.400 3032.400
\hypertarget{w02:cue:1073}{}去授權，
% cue 1074 3032.400 3034.500
\hypertarget{w02:cue:1074}{}來採取某樣的政治行動，
% cue 1075 3034.500 3036.660
\hypertarget{w02:cue:1075}{}那這樣子我們才能處理是這張。

\tc{0:50:36}
% cue 1076 3036.660 3038.600
\hypertarget{w02:cue:1076}{}那這有很多種說法，
% cue 1077 3038.600 3041.430
\hypertarget{w02:cue:1077}{}很多個不同的 versions。
% cue 1078 3041.430 3043.430
\hypertarget{w02:cue:1078}{}那第一個就叫做 consent theory，
% cue 1079 3043.430 3044.940
\hypertarget{w02:cue:1079}{}叫做同意說。
% cue 1080 3044.940 3046.510
\hypertarget{w02:cue:1080}{}什麼意思？
% cue 1081 3046.510 3047.510
\hypertarget{w02:cue:1081}{}for example，
% cue 1082 3047.510 3050.020
\hypertarget{w02:cue:1082}{}It has stated that fact
% cue 1083 3050.020 3052.240
\hypertarget{w02:cue:1083}{}individual citizen have
% cue 1084 3052.240 3055.810
\hypertarget{w02:cue:1084}{}would have voluntarily consented to be governed
% cue 1085 3055.810 3057.890
\hypertarget{w02:cue:1085}{}by a given form of government
% cue 1086 3057.890 3062.140
\hypertarget{w02:cue:1086}{}constitute sufficient reason to regard the action of the government
% cue 1087 3062.140 3064.150
\hypertarget{w02:cue:1087}{}as specifically authorized。

\tc{0:51:04}
% cue 1088 3064.150 3065.430
\hypertarget{w02:cue:1088}{}當人民，
% cue 1089 3065.430 3067.060
\hypertarget{w02:cue:1089}{}當人民，
% cue 1090 3067.060 3068.440
\hypertarget{w02:cue:1090}{}當人民，
% cue 1091 3068.440 3069.440
\hypertarget{w02:cue:1091}{}在某些情況下說，
% cue 1092 3069.440 3070.790
\hypertarget{w02:cue:1092}{}我自願的，
% cue 1093 3070.790 3073.140
\hypertarget{w02:cue:1093}{}我是真心誠意的，
% cue 1094 3073.140 3076.610
\hypertarget{w02:cue:1094}{}同意了政府的統治行為，
% cue 1095 3076.610 3077.900
\hypertarget{w02:cue:1095}{}那麼政府，
% cue 1096 3077.900 3080.020
\hypertarget{w02:cue:1096}{}這個統治就具有正當性。
% cue 1097 3080.020 3081.530
\hypertarget{w02:cue:1097}{}政府這個決定，
% cue 1098 3081.530 3082.520
\hypertarget{w02:cue:1098}{}這個決定，
% cue 1099 3082.520 3083.480
\hypertarget{w02:cue:1099}{}就會具體的授權，
% cue 1100 3083.480 3084.980
\hypertarget{w02:cue:1100}{}特定的授權。

\tc{0:51:24}
% cue 1101 3084.980 3086.050
\hypertarget{w02:cue:1101}{}我有同意，
% cue 1102 3086.050 3089.440
\hypertarget{w02:cue:1102}{}我同意你這麼做，
% cue 1103 3089.440 3090.570
\hypertarget{w02:cue:1103}{}那我這個同意，
% cue 1104 3090.570 3091.570
\hypertarget{w02:cue:1104}{}我這個授權，
% cue 1105 3091.570 3093.730
\hypertarget{w02:cue:1105}{}就是使得你變得正當。
% cue 1106 3093.730 3096.650
\hypertarget{w02:cue:1106}{}我可能沒有泛泛的同意，
% cue 1107 3096.650 3099.470
\hypertarget{w02:cue:1107}{}但是我有具體的。

\tc{0:51:39}
% cue 1108 3099.470 3102.300
\hypertarget{w02:cue:1108}{}without such an actual consent，
% cue 1109 3102.300 3104.720
\hypertarget{w02:cue:1109}{}如果沒有這樣的同意，
% cue 1110 3104.720 3106.720
\hypertarget{w02:cue:1110}{}laws enacted by such government，
% cue 1111 3106.720 3107.910
\hypertarget{w02:cue:1111}{}這個政府所制定的法，
% cue 1112 3107.910 3109.390
\hypertarget{w02:cue:1112}{}它在道德上，
% cue 1113 3109.390 3110.390
\hypertarget{w02:cue:1113}{}不會被允許，
% cue 1114 3110.390 3112.390
\hypertarget{w02:cue:1114}{}也就是不會具有正當性。

\tc{0:51:52}
% cue 1115 3112.390 3113.390
\hypertarget{w02:cue:1115}{}道德上的允許，
% cue 1116 3113.390 3115.350
\hypertarget{w02:cue:1116}{}道德上的 permissible，
% cue 1117 3115.350 3116.350
\hypertarget{w02:cue:1117}{}就是 legent of 法。

\tc{0:51:56}
% cue 1118 3116.350 3118.090
\hypertarget{w02:cue:1118}{}我剛剛就畫了，
% cue 1119 3118.090 3119.210
\hypertarget{w02:cue:1119}{}我一個圖形，
% cue 1120 3119.210 3119.310
\hypertarget{w02:cue:1120}{}你畫得出來。
% cue 1121 3119.210 3120.690
\hypertarget{w02:cue:1121}{}你一面講，
% cue 1122 3120.690 3121.690
\hypertarget{w02:cue:1122}{}一面畫，
% cue 1123 3121.690 3121.790
\hypertarget{w02:cue:1123}{}一面寫，
% cue 1124 3121.690 3122.690
\hypertarget{w02:cue:1124}{}就不知道他講什麼了。

\tc{0:52:02}
% cue 1125 3122.690 3124.480
\hypertarget{w02:cue:1125}{}你當你讀到，
% cue 1126 3124.480 3125.440
\hypertarget{w02:cue:1126}{}覺得比較困難的時候，
% cue 1127 3125.440 3127.440
\hypertarget{w02:cue:1127}{}就把那個代曼語句拿出來，
% cue 1128 3127.440 3128.440
\hypertarget{w02:cue:1128}{}知道他在講什麼。
% cue 1129 3128.440 3129.590
\hypertarget{w02:cue:1129}{}他不要一直講這個具體，
% cue 1130 3129.590 3131.900
\hypertarget{w02:cue:1130}{}我覺得你這樣很反應。

\tc{0:52:11}
% cue 1131 3131.900 3134.180
\hypertarget{w02:cue:1131}{}以上是第一說，
% cue 1132 3134.180 3135.280
\hypertarget{w02:cue:1132}{}同意說，
% cue 1133 3135.280 3137.750
\hypertarget{w02:cue:1133}{}人民對某個具體行為有同意，
% cue 1134 3137.750 3139.780
\hypertarget{w02:cue:1134}{}那他所實施的這個，
% cue 1135 3139.780 3140.780
\hypertarget{w02:cue:1135}{}你就要同意，
% cue 1136 3140.780 3142.260
\hypertarget{w02:cue:1136}{}承認，
% cue 1137 3142.260 3143.260
\hypertarget{w02:cue:1137}{}而具有正當性。

\tc{0:52:23}
% cue 1138 3143.260 3145.680
\hypertarget{w02:cue:1138}{}第二個叫 democratic theory，
% cue 1139 3145.680 3148.280
\hypertarget{w02:cue:1139}{}supplies another principle，
% cue 1140 3148.280 3149.280
\hypertarget{w02:cue:1140}{}It states that the fact
% cue 1141 3149.280 3151.400
\hypertarget{w02:cue:1141}{}a given piece of legislation
% cue 1142 3151.400 3152.360
\hypertarget{w02:cue:1142}{}某個例子，
% cue 1143 3152.360 3154.400
\hypertarget{w02:cue:1143}{}has a democratic pedigree，
% cue 1144 3154.400 3156.400
\hypertarget{w02:cue:1144}{}具有某一種民主的血統，
% cue 1145 3156.400 3158.530
\hypertarget{w02:cue:1145}{}make it morally permissible，
% cue 1146 3158.530 3160.850
\hypertarget{w02:cue:1146}{}這個法律具有民主的血統，
% cue 1147 3160.850 3163.910
\hypertarget{w02:cue:1147}{}那他就具有道德上的正當性。
% cue 1148 3163.910 3167.520
\hypertarget{w02:cue:1148}{}這是第二個理論。

\tc{0:52:47}
% cue 1149 3167.520 3168.970
\hypertarget{w02:cue:1149}{}第三個理論，
% cue 1150 3168.970 3172.440
\hypertarget{w02:cue:1150}{}other principles may point to other facts
% cue 1151 3172.440 3175.520
\hypertarget{w02:cue:1151}{}such as immemorial custom\footnote{英文，immemorial 意為久遠到超出記憶所及，immemorial custom 即自古相沿、起源已不可考的習慣。本文第 528 頁以之與神法、自然法、憲法並列，為特定正當性所援引的積極事實。\url{https://en.wiktionary.org/wiki/immemorial}。}，
% cue 1152 3175.520 3178.410
\hypertarget{w02:cue:1152}{}有一種不朽的習慣，
% cue 1153 3178.410 3179.820
\hypertarget{w02:cue:1153}{}divine law，
% cue 1154 3179.820 3180.790
\hypertarget{w02:cue:1154}{}神聖的法律，
% cue 1155 3180.790 3182.980
\hypertarget{w02:cue:1155}{}the law of nature，
% cue 1156 3182.980 3184.110
\hypertarget{w02:cue:1156}{}自然法，
% cue 1157 3184.110 3186.300
\hypertarget{w02:cue:1157}{}a constitution and so on，
% cue 1158 3186.300 3189.460
\hypertarget{w02:cue:1158}{}好多好多種不同的來源，
% cue 1159 3189.460 3191.620
\hypertarget{w02:cue:1159}{}你說我得到他的授權，
% cue 1160 3191.620 3192.650
\hypertarget{w02:cue:1160}{}他都是正當的。

\tc{0:53:12}
% cue 1161 3192.650 3196.930
\hypertarget{w02:cue:1161}{}this answer seeks to establish what I should call
% cue 1162 3196.930 3200.310
\hypertarget{w02:cue:1162}{}a specific legitimacy of political authority，
% cue 1163 3200.310 3202.500
\hypertarget{w02:cue:1163}{}某個政治權威的特定授權。
% cue 1164 3202.500 3205.780
\hypertarget{w02:cue:1164}{}this they point to positive legitimating
% cue 1165 3205.780 3206.740
\hypertarget{w02:cue:1165}{}facts，
% cue 1166 3206.740 3209.510
\hypertarget{w02:cue:1166}{}這個是有一個具體的，
% cue 1167 3209.510 3211.610
\hypertarget{w02:cue:1167}{}使它成為正當的事實，
% cue 1168 3211.610 3214.340
\hypertarget{w02:cue:1168}{}沒有這樣的具體的事實，
% cue 1169 3214.340 3216.370
\hypertarget{w02:cue:1169}{}使它成為正當的事實，
% cue 1170 3216.370 3218.560
\hypertarget{w02:cue:1170}{}也就是說沒有這種具體的授權，
% cue 1171 3218.560 3219.910
\hypertarget{w02:cue:1171}{}這不可能阻障。
% cue 1172 3219.910 3220.910
\hypertarget{w02:cue:1172}{}是有一個具體的，
% cue 1173 3220.910 3230.020
\hypertarget{w02:cue:1173}{}使它成為合法的事實。

\tc{0:53:50}
% cue 1174 3230.020 3231.020
\hypertarget{w02:cue:1174}{}泛泛的講什麼叫授權，
% cue 1175 3231.020 3233.040
\hypertarget{w02:cue:1175}{}你說什麼叫授權，
% cue 1176 3233.040 3234.450
\hypertarget{w02:cue:1176}{}這就是授權。
% cue 1177 3234.450 3235.860
\hypertarget{w02:cue:1177}{}說法，
% cue 1178 3235.860 3236.860
\hypertarget{w02:cue:1178}{}有的是用同意來表現，
% cue 1179 3236.860 3238.880
\hypertarget{w02:cue:1179}{}有的是用民主來表現，
% cue 1180 3238.880 3240.260
\hypertarget{w02:cue:1180}{}有的是用其他的法院來表現，
% cue 1181 3240.260 3241.290
\hypertarget{w02:cue:1181}{}有的是用其他的法院來表現。

\tc{0:54:01}
% cue 1182 3241.290 3243.660
\hypertarget{w02:cue:1182}{}好。
% cue 1183 3243.660 3247.670
\hypertarget{w02:cue:1183}{}the judiciary exercises effective coercive power
% cue 1184 3247.670 3248.670
\hypertarget{w02:cue:1184}{}and has political authority
% cue 1185 3248.670 3249.920
\hypertarget{w02:cue:1185}{}and has political authority
% cue 1186 3249.920 3250.920
\hypertarget{w02:cue:1186}{}within society，
% cue 1187 3250.920 3256.390
\hypertarget{w02:cue:1187}{}司法部門具有強制力，
% cue 1188 3256.390 3257.380
\hypertarget{w02:cue:1188}{}有效的強制力量，
% cue 1189 3257.380 3259.540
\hypertarget{w02:cue:1189}{}而且它是一個政治的權威。

\tc{0:54:19}
% cue 1190 3259.540 3260.570
\hypertarget{w02:cue:1190}{}這件事情可能要稍微說一下，
% cue 1191 3260.570 3261.570
\hypertarget{w02:cue:1191}{}因為在臺灣，
% cue 1192 3261.570 3264.740
\hypertarget{w02:cue:1192}{}大概你講說司法是一個政治的權威。
% cue 1193 3264.740 3266.740
\hypertarget{w02:cue:1193}{}這件事情可能要稍微說一下，
% cue 1194 3266.740 3267.860
\hypertarget{w02:cue:1194}{}因為在臺灣，
% cue 1195 3267.860 3270.360
\hypertarget{w02:cue:1195}{}大概你講說司法是一個政治的權威。
% cue 1196 3270.360 3271.550
\hypertarget{w02:cue:1196}{}這件事情可能要稍微說一下，
% cue 1197 3271.550 3272.510
\hypertarget{w02:cue:1197}{}因為在臺灣，
% cue 1198 3272.510 3275.090
\hypertarget{w02:cue:1198}{}大概你講說司法是一個政治的權威。
% cue 1199 3275.090 3277.180
\hypertarget{w02:cue:1199}{}這件事情可能要稍微說一下，
% cue 1200 3277.180 3278.210
\hypertarget{w02:cue:1200}{}因為在臺灣，
% cue 1201 3278.210 3280.220
\hypertarget{w02:cue:1201}{}大概你講說司法是一個政治的權威。
% cue 1202 3280.220 3282.620
\hypertarget{w02:cue:1202}{}這件事情可能要稍微說一下，
% cue 1203 3282.620 3284.830
\hypertarget{w02:cue:1203}{}因為在臺灣，
% cue 1204 3284.830 3288.460
\hypertarget{w02:cue:1204}{}大概你講說司法是一個政治權。

\tc{0:54:48}
% cue 1205 3288.460 3289.810
\hypertarget{w02:cue:1205}{}哇，
% cue 1206 3289.810 3291.160
\hypertarget{w02:cue:1206}{}很多人就覺得你在那裡，
% cue 1207 3291.160 3292.540
\hypertarget{w02:cue:1207}{}在污衊司法。
% cue 1208 3292.540 3294.080
\hypertarget{w02:cue:1208}{}我們是獨立的，
% cue 1209 3294.080 3296.140
\hypertarget{w02:cue:1209}{}怎麼說我們是政治權呢？

\tc{0:54:56}
% cue 1210 3296.140 3298.740
\hypertarget{w02:cue:1210}{}但是這個作者他不是這樣想的，
% cue 1211 3298.740 3301.120
\hypertarget{w02:cue:1211}{}就是純粹從憲法的觀點來講，
% cue 1212 3301.120 3303.540
\hypertarget{w02:cue:1212}{}當然是一個權力，
% cue 1213 3303.540 3304.780
\hypertarget{w02:cue:1213}{}只要它是權力，
% cue 1214 3304.780 3305.780
\hypertarget{w02:cue:1214}{}它就是 political authority。
% cue 1215 3305.780 3307.800
\hypertarget{w02:cue:1215}{}聽得懂我講話的意思嗎？

\tc{0:55:07}
% cue 1216 3307.800 3309.990
\hypertarget{w02:cue:1216}{}用德國人的話講是，
% cue 1217 3309.990 3310.990
\hypertarget{w02:cue:1217}{}這是我們憲法機關\footnote{德文 Verfassungsorgan。德國聯邦憲法法院法第 1 條第 1 項規定：「Das Bundesverfassungsgericht ist ein allen übrigen Verfassungsorganen gegenüber selbständiger und unabhängiger Gerichtshof des Bundes.」（聯邦憲法法院為對其他一切憲法機關均獨立自主的聯邦法院）。\url{https://www.gesetze-im-internet.de/bverfgg/__1.html}。}。
% cue 1218 3310.990 3313.880
\hypertarget{w02:cue:1218}{}請問你憲法機關何以取得權力？
% cue 1219 3313.880 3314.910
\hypertarget{w02:cue:1219}{}依據憲法。
% cue 1220 3314.910 3317.030
\hypertarget{w02:cue:1220}{}那憲法是什麼東西？
% cue 1221 3317.030 3319.120
\hypertarget{w02:cue:1221}{}憲法是一套政治的權威，
% cue 1222 3319.120 3321.310
\hypertarget{w02:cue:1222}{}所以它一定是一個 political authority。
% cue 1223 3321.310 3325.620
\hypertarget{w02:cue:1223}{}只不過你這個 political authority 的權力來源跟行程運動方式，
% cue 1224 3325.620 3328.620
\hypertarget{w02:cue:1224}{}可能跟其他的政治部門也許不同，
% cue 1225 3328.620 3330.840
\hypertarget{w02:cue:1225}{}但是你絕對是一個 political authority，
% cue 1226 3330.840 3332.190
\hypertarget{w02:cue:1226}{}不然我為什麼要服從你，
% cue 1227 3332.190 3334.830
\hypertarget{w02:cue:1227}{}我幹嘛聊你法意要幹去。
% cue 1228 3334.830 3336.160
\hypertarget{w02:cue:1228}{}ok？

\tc{0:55:36}
% cue 1229 3336.160 3339.270
\hypertarget{w02:cue:1229}{}那，by virtue of what consideration 又來了，
% cue 1230 3339.270 3340.550
\hypertarget{w02:cue:1230}{}what consideration，
% cue 1231 3340.550 3342.700
\hypertarget{w02:cue:1231}{}是不是正當性嗎？
% cue 1232 3342.700 3344.470
\hypertarget{w02:cue:1232}{}哪些，
% cue 1233 3344.470 3347.470
\hypertarget{w02:cue:1233}{}can we say that the power of this authority is legitimate?
% cue 1234 3347.470 3350.170
\hypertarget{w02:cue:1234}{}是根據什麼樣的考慮，
% cue 1235 3350.170 3351.170
\hypertarget{w02:cue:1235}{}我們說它是正當？
% cue 1236 3351.170 3352.070
\hypertarget{w02:cue:1236}{}然後，
% cue 1237 3352.070 3353.140
\hypertarget{w02:cue:1237}{}從而，
% cue 1238 3353.140 3355.140
\hypertarget{w02:cue:1238}{}its action and decision are morally permissible，
% cue 1239 3355.140 3357.170
\hypertarget{w02:cue:1239}{}從而它的決定是道德上允許的，
% cue 1240 3357.170 3359.290
\hypertarget{w02:cue:1240}{}道德上允許就是正當。

\tc{0:55:59}
% cue 1241 3359.290 3364.050
\hypertarget{w02:cue:1241}{}the modern answer has taken the form of principles
% cue 1242 3364.050 3368.010
\hypertarget{w02:cue:1242}{}establishing the specific legitimacy，
% cue 1243 3368.010 3371.010
\hypertarget{w02:cue:1243}{}I call it the specific legitimacy thesis，
% cue 1244 3371.010 3374.580
\hypertarget{w02:cue:1244}{}那我就把它們通稱叫做
% cue 1245 3374.580 3376.740
\hypertarget{w02:cue:1245}{}"特定正當性命題"。

\tc{0:56:16}
% cue 1246 3376.740 3377.660
\hypertarget{w02:cue:1246}{}好。
% cue 1247 3377.660 3381.250
\hypertarget{w02:cue:1247}{}it can be summarized into four propositions，
% cue 1248 3381.250 3383.590
\hypertarget{w02:cue:1248}{}看看你寫，
% cue 1249 3383.590 3384.810
\hypertarget{w02:cue:1249}{}看你有沒有辦法寫，
% cue 1250 3384.810 3386.810
\hypertarget{w02:cue:1250}{}正當性會包含以下四命題，
% cue 1251 3386.810 3389.470
\hypertarget{w02:cue:1251}{}就是它做成同一共產部。

\tc{0:56:29}
% cue 1252 3389.470 3391.870
\hypertarget{w02:cue:1252}{}第一，
% cue 1253 3391.870 3394.620
\hypertarget{w02:cue:1253}{}"特定正當性命題"，
% cue 1254 3394.620 3395.620
\hypertarget{w02:cue:1254}{}一，
% cue 1255 3395.620 3397.060
\hypertarget{w02:cue:1255}{}SLT1，
% cue 1256 3397.060 3401.340
\hypertarget{w02:cue:1256}{}"The judiciary has no inherent legitimacy"，
% cue 1257 3401.340 3405.130
\hypertarget{w02:cue:1257}{}司法是沒有內在正當性的，
% cue 1258 3405.130 3407.130
\hypertarget{w02:cue:1258}{}或者沒有本質正當性。
% cue 1259 3407.130 3409.370
\hypertarget{w02:cue:1259}{}這個可以理解嗎？

\tc{0:56:49}
% cue 1260 3409.370 3411.010
\hypertarget{w02:cue:1260}{}我們是個民主政府，
% cue 1261 3411.010 3413.810
\hypertarget{w02:cue:1261}{}司法也不是我們選的，
% cue 1262 3413.810 3415.640
\hypertarget{w02:cue:1262}{}它也不對我們負責，
% cue 1263 3415.640 3418.150
\hypertarget{w02:cue:1263}{}從哪裡來的正當性？
% cue 1264 3418.150 3420.370
\hypertarget{w02:cue:1264}{}它當然本質上沒有正當性，
% cue 1265 3420.370 3423.230
\hypertarget{w02:cue:1265}{}它一定要從別的地方取得正當性，
% cue 1266 3423.230 3426.000
\hypertarget{w02:cue:1266}{}它不可能有本質正當性。

\tc{0:57:06}
% cue 1267 3426.000 3428.000
\hypertarget{w02:cue:1267}{}這裡可能換個句子來講，
% cue 1268 3428.000 3430.000
\hypertarget{w02:cue:1268}{}總統或者民意代表，
% cue 1269 3430.000 3431.000
\hypertarget{w02:cue:1269}{}他可能可以主張，
% cue 1270 3431.000 3432.960
\hypertarget{w02:cue:1270}{}他具有某種本質上的正當性，
% cue 1271 3432.960 3434.250
\hypertarget{w02:cue:1271}{}因為你們選我，
% cue 1272 3434.250 3435.410
\hypertarget{w02:cue:1272}{}你們授權給我，
% cue 1273 3435.410 3436.410
\hypertarget{w02:cue:1273}{}你們同意我，
% cue 1274 3436.410 3438.140
\hypertarget{w02:cue:1274}{}不是嗎？
% cue 1275 3438.140 3439.140
\hypertarget{w02:cue:1275}{}不然我怎麼會行使權力？

\tc{0:57:19}
% cue 1276 3439.140 3440.480
\hypertarget{w02:cue:1276}{}好，
% cue 1277 3440.480 3441.440
\hypertarget{w02:cue:1277}{}以上一。
% cue 1278 3441.440 3442.470
\hypertarget{w02:cue:1278}{}這一點要記住前提，
% cue 1279 3442.470 3444.470
\hypertarget{w02:cue:1279}{}是司法沒有本質正當性，
% cue 1280 3444.470 3446.980
\hypertarget{w02:cue:1280}{}所以你得去證立它的正當性。

\tc{0:57:26}
% cue 1281 3446.980 3448.070
\hypertarget{w02:cue:1281}{}第二，
% cue 1282 3448.070 3449.380
\hypertarget{w02:cue:1282}{}SLT2，
% cue 1283 3449.380 3454.660
\hypertarget{w02:cue:1283}{}Specific legitimacy of judicial power is conditioned，
% cue 1284 3454.660 3455.790
\hypertarget{w02:cue:1284}{}這個字圈起來，
% cue 1285 3455.790 3456.730
\hypertarget{w02:cue:1285}{}條件，
% cue 1286 3456.730 3460.190
\hypertarget{w02:cue:1286}{}by the rule of law，
% cue 1287 3460.190 3461.150
\hypertarget{w02:cue:1287}{}所謂的，
% cue 1288 3461.150 3464.170
\hypertarget{w02:cue:1288}{}違憲審查特定正當性，
% cue 1289 3464.170 3469.400
\hypertarget{w02:cue:1289}{}是以法治作為條件，
% cue 1290 3469.400 3470.810
\hypertarget{w02:cue:1290}{}這是副條件。
% cue 1291 3470.810 3471.960
\hypertarget{w02:cue:1291}{}換句話說，
% cue 1292 3471.960 3474.630
\hypertarget{w02:cue:1292}{}如果你不符合法治這個條件，
% cue 1293 3474.630 3478.050
\hypertarget{w02:cue:1293}{}你就不可能取得特定的正當性。

\tc{0:57:58}
% cue 1294 3478.050 3480.520
\hypertarget{w02:cue:1294}{}這就是三個例子。
% cue 1295 3480.520 3482.830
\hypertarget{w02:cue:1295}{}這個你會不會想，
% cue 1296 3482.830 3484.150
\hypertarget{w02:cue:1296}{}什麼叫法治？
% cue 1297 3484.150 3485.310
\hypertarget{w02:cue:1297}{}自止三件事。
% cue 1298 3485.310 3488.780
\hypertarget{w02:cue:1298}{}這個地方對於法治的看法，
% cue 1299 3488.780 3489.960
\hypertarget{w02:cue:1299}{}大家各有不同，
% cue 1300 3489.960 3492.470
\hypertarget{w02:cue:1300}{}這個是一種說法，
% cue 1301 3492.470 3494.470
\hypertarget{w02:cue:1301}{}我就這樣講就好了，
% cue 1302 3494.470 3496.470
\hypertarget{w02:cue:1302}{}我的見解也許不完全大。

\tc{0:58:16}
% cue 1303 3496.470 3497.470
\hypertarget{w02:cue:1303}{}第一，
% cue 1304 3497.470 3502.810
\hypertarget{w02:cue:1304}{}the constitution of the judiciary must be according to law，
% cue 1305 3502.810 3506.020
\hypertarget{w02:cue:1305}{}有關於司法的組成，
% cue 1306 3506.020 3509.330
\hypertarget{w02:cue:1306}{}看 the constitution 這個字，
% cue 1307 3509.330 3510.330
\hypertarget{w02:cue:1307}{}如果你還懷疑，
% cue 1308 3510.330 3511.550
\hypertarget{w02:cue:1308}{}那你再看它後面，
% cue 1309 3511.550 3514.700
\hypertarget{w02:cue:1309}{}its structure and composition，
% cue 1310 3514.700 3516.240
\hypertarget{w02:cue:1310}{}法院的組成，
% cue 1311 3516.240 3518.210
\hypertarget{w02:cue:1311}{}還有法官的組成，
% cue 1312 3518.210 3523.460
\hypertarget{w02:cue:1312}{}structure and composition must be according to law，
% cue 1313 3523.460 3524.430
\hypertarget{w02:cue:1313}{}一定要依法。

\tc{0:58:44}
% cue 1314 3524.430 3526.420
\hypertarget{w02:cue:1314}{}可以想想看，
% cue 1315 3526.420 3528.840
\hypertarget{w02:cue:1315}{}今天 114 憲判一，
% cue 1316 3528.840 3529.960
\hypertarget{w02:cue:1316}{}所以我要復活了。
% cue 1317 3529.960 3532.660
\hypertarget{w02:cue:1317}{}那它行為它的爭議，
% cue 1318 3532.660 3533.820
\hypertarget{w02:cue:1318}{}是跟這個有關的。

\tc{0:58:53}
% cue 1319 3533.820 3535.970
\hypertarget{w02:cue:1319}{}但是這個爭議有關的是，
% cue 1320 3535.970 3537.030
\hypertarget{w02:cue:1320}{}第一，
% cue 1321 3537.030 3538.090
\hypertarget{w02:cue:1321}{}它的人數，
% cue 1322 3538.090 3539.120
\hypertarget{w02:cue:1322}{}它的 composition，
% cue 1323 3539.120 3540.440
\hypertarget{w02:cue:1323}{}它的結構，
% cue 1324 3540.440 3541.440
\hypertarget{w02:cue:1324}{}它有 according to law 嗎？
% cue 1325 3541.440 3544.510
\hypertarget{w02:cue:1325}{}這樣講有聽懂了嗎？

\tc{0:59:04}
% cue 1326 3544.510 3546.160
\hypertarget{w02:cue:1326}{}你要講依法，
% cue 1327 3546.160 3547.130
\hypertarget{w02:cue:1327}{}我們現在在這邊講，
% cue 1328 3547.130 3549.130
\hypertarget{w02:cue:1328}{}抽象的話叫做依法審案，
% cue 1329 3549.130 3550.130
\hypertarget{w02:cue:1329}{}依法行政，
% cue 1330 3550.130 3553.340
\hypertarget{w02:cue:1330}{}那你看看這是不是都是屬於法治？
% cue 1331 3553.340 3554.560
\hypertarget{w02:cue:1331}{}rule of law，
% cue 1332 3554.560 3555.620
\hypertarget{w02:cue:1332}{}rule of law，
% cue 1333 3555.620 3556.620
\hypertarget{w02:cue:1333}{}對它來講，
% cue 1334 3556.620 3557.620
\hypertarget{w02:cue:1334}{}它一層層講給你聽。

\tc{0:59:17}
% cue 1335 3557.620 3558.580
\hypertarget{w02:cue:1335}{}第一層，
% cue 1336 3558.580 3561.660
\hypertarget{w02:cue:1336}{}有關於法院 judicial，
% cue 1337 3561.660 3563.660
\hypertarget{w02:cue:1337}{}這種司法和法案具體的話講，
% cue 1338 3563.660 3564.660
\hypertarget{w02:cue:1338}{}就是法院，
% cue 1339 3564.660 3566.660
\hypertarget{w02:cue:1339}{}法院的組成跟法院的結構，
% cue 1340 3566.660 3567.660
\hypertarget{w02:cue:1340}{}它必須依法，
% cue 1341 3567.660 3568.950
\hypertarget{w02:cue:1341}{}我沒有設立這個法院，
% cue 1342 3568.950 3569.950
\hypertarget{w02:cue:1342}{}法院可以自己說，
% cue 1343 3569.950 3570.950
\hypertarget{w02:cue:1343}{}我要成立嗎？
% cue 1344 3570.950 3572.290
\hypertarget{w02:cue:1344}{}現在有這種法院嗎？
% cue 1345 3572.290 3574.380
\hypertarget{w02:cue:1345}{}這種法院可以有正當性嗎？
% cue 1346 3574.490 3576.010
\hypertarget{w02:cue:1346}{}這是第一點。

\tc{0:59:36}
% cue 1347 3576.010 3577.020
\hypertarget{w02:cue:1347}{}第二點。
% cue 1348 3577.020 3577.820
\hypertarget{w02:cue:1348}{}第二點。
% cue 1349 3577.820 3580.900
\hypertarget{w02:cue:1349}{}the judicial process, the manner and procedure。
% cue 1350 3580.900 3583.360
\hypertarget{w02:cue:1350}{}它的運轉形式，運作形式。
% cue 1351 3583.290 3584.290
\hypertarget{w02:cue:1351}{}這是，
% cue 1352 3584.290 3586.290
\hypertarget{w02:cue:1352}{}by virtue of which judicial decisions are reached，
% cue 1353 3586.290 3587.320
\hypertarget{w02:cue:1353}{}經由這個程序，
% cue 1354 3587.320 3588.320
\hypertarget{w02:cue:1354}{}做這個司法決定，
% cue 1355 3588.320 3590.320
\hypertarget{w02:cue:1355}{}must also be according to law，
% cue 1356 3590.320 3591.320
\hypertarget{w02:cue:1356}{}很好，
% cue 1357 3591.320 3594.000
\hypertarget{w02:cue:1357}{}那，
% cue 1358 3594.000 3594.100
\hypertarget{w02:cue:1358}{}首先，
% cue 1359 3594.000 3595.170
\hypertarget{w02:cue:1359}{}你要十個\footnote{2025 年 1 月 23 日修正公布的憲法訴訟法第 30 條第 2 項規定：「前項參與評議之大法官人數不得低於十人。作成違憲之宣告時，同意違憲宣告之大法官人數不得低於九人。」中華民國憲法法庭 114 年憲判字第 1 號判決（2025 年 12 月 19 日）宣告該項及同條第 3 項至第 6 項等規定牴觸憲法，自判決公告之日起失其效力。\cite{ccj114-1}。}，
% cue 1360 3595.170 3596.170
\hypertarget{w02:cue:1360}{}你現在只有八個\footnote{憲法增修條文第 5 條第 1 項規定司法院設大法官十五人。2024 年 10 月 31 日許宗力、蔡烱燉、黃昭元、許志雄、張瓊文、黃瑞明、詹森林七位大法官任期屆滿後，在任者為謝銘洋、呂太郎、楊惠欽、蔡宗珍、蔡彩貞、朱富美、陳忠五、尤伯祥八人；114 年憲判字第 1 號判決理由第 41 段載明該庭當時大法官人數為 8 人。\cite[41]{ccj114-1}。}，
% cue 1361 3596.170 3597.160
\hypertarget{w02:cue:1361}{}第二，
% cue 1362 3597.160 3598.160
\hypertarget{w02:cue:1362}{}我叫你要合議\footnote{憲法訴訟法第 30 條第 1 項規定，判決除本法別有規定外，應經大法官現有總額三分之二以上參與評議，大法官現有總額過半數同意；114 年憲判字第 1 號判決理由第 42 至 43 段即依此項定其評議及評決門檻。該判決所附大法官立場表列名者為謝銘洋、呂太郎、蔡彩貞、陳忠五、尤伯祥五人，楊惠欽、蔡宗珍、朱富美三位大法官未列名。\cite[42–43]{ccj114-1}。}，
% cue 1363 3598.160 3599.220
\hypertarget{w02:cue:1363}{}你有合議嗎？
% cue 1364 3599.220 3603.040
\hypertarget{w02:cue:1364}{}那你不知道，
% cue 1365 3603.040 3604.580
\hypertarget{w02:cue:1365}{}依照我的理解，
% cue 1366 3604.580 3607.000
\hypertarget{w02:cue:1366}{}現在司法院大的問題，
% cue 1367 3607.000 3608.810
\hypertarget{w02:cue:1367}{}沒有合議。

\tc{1:00:08}
% cue 1368 3608.810 3610.380
\hypertarget{w02:cue:1368}{}這個事情以後，
% cue 1369 3610.380 3611.440
\hypertarget{w02:cue:1369}{}如果我們在學校，
% cue 1370 3611.440 3613.500
\hypertarget{w02:cue:1370}{}你不是哪一天要辦，
% cue 1371 3613.500 3614.470
\hypertarget{w02:cue:1371}{}哪一天要辦呢？
% cue 1372 3614.470 3614.570
\hypertarget{w02:cue:1372}{}證明，
% cue 1373 3614.470 3616.640
\hypertarget{w02:cue:1373}{}11 月 18 號，
% cue 1374 3616.640 3618.860
\hypertarget{w02:cue:1374}{}我希望要辦一場，
% cue 1375 3618.860 3620.300
\hypertarget{w02:cue:1375}{}導師的新書發表會，
% cue 1376 3620.300 3621.300
\hypertarget{w02:cue:1376}{}200 的新書發表會。

\tc{1:00:21}
% cue 1377 3621.300 3625.880
\hypertarget{w02:cue:1377}{}現在司法院，
% cue 1378 3625.880 3628.920
\hypertarget{w02:cue:1378}{}這個是沒有依法，
% cue 1379 3628.920 3631.300
\hypertarget{w02:cue:1379}{}你的運作方式沒有依法，
% cue 1380 3631.300 3632.620
\hypertarget{w02:cue:1380}{}第三點，
% cue 1381 3632.620 3633.750
\hypertarget{w02:cue:1381}{}最重要的，
% cue 1382 3633.750 3636.040
\hypertarget{w02:cue:1382}{}就本文而言，
% cue 1383 3636.040 3638.100
\hypertarget{w02:cue:1383}{}the content and substance，
% cue 1384 3638.100 3639.060
\hypertarget{w02:cue:1384}{}它的內容，
% cue 1385 3639.060 3640.020
\hypertarget{w02:cue:1385}{}它的內涵，
% cue 1386 3640.020 3641.150
\hypertarget{w02:cue:1386}{}of judicial decision，
% cue 1387 3641.150 3644.150
\hypertarget{w02:cue:1387}{}must be in accordance with law，
% cue 1388 3644.150 3646.590
\hypertarget{w02:cue:1388}{}你做成的判決內容，
% cue 1389 3646.590 3649.750
\hypertarget{w02:cue:1389}{}也必須要符合法律的規定，
% cue 1390 3649.750 3651.350
\hypertarget{w02:cue:1390}{}although contested，
% cue 1391 3651.350 3652.570
\hypertarget{w02:cue:1391}{}儘管有爭議，
% cue 1392 3652.570 3654.630
\hypertarget{w02:cue:1392}{}the concept of the rule of law，
% cue 1393 3654.630 3656.610
\hypertarget{w02:cue:1393}{}它曾認為有爭議嗎？
% cue 1394 3656.610 3657.720
\hypertarget{w02:cue:1394}{}所謂法治的這個概念，
% cue 1395 3657.720 3660.850
\hypertarget{w02:cue:1395}{}entails two basic ideas，
% cue 1396 3660.850 3663.940
\hypertarget{w02:cue:1396}{}它會導出兩個基本觀念。

\tc{1:01:03}
% cue 1397 3663.940 3665.090
\hypertarget{w02:cue:1397}{}first，
% cue 1398 3665.090 3668.150
\hypertarget{w02:cue:1398}{}government actions and decisions must be rational，
% cue 1399 3668.150 3669.720
\hypertarget{w02:cue:1399}{}政府的決定，
% cue 1400 3669.720 3671.740
\hypertarget{w02:cue:1400}{}必須是合理的，
% cue 1401 3671.740 3674.120
\hypertarget{w02:cue:1401}{}必須是理性的，
% cue 1402 3674.120 3676.310
\hypertarget{w02:cue:1402}{}那也就是說 based on reason，
% cue 1403 3676.310 3679.460
\hypertarget{w02:cue:1403}{}是基於理智而做的決定，
% cue 1404 3679.460 3681.030
\hypertarget{w02:cue:1404}{}基於理智，
% cue 1405 3681.030 3682.030
\hypertarget{w02:cue:1405}{}它就是 rational，
% cue 1406 3682.030 3686.140
\hypertarget{w02:cue:1406}{}你覺得法治是不是要這樣子，
% cue 1407 3686.140 3687.400
\hypertarget{w02:cue:1407}{}人治，
% cue 1408 3687.400 3688.400
\hypertarget{w02:cue:1408}{}我今天情緒不好，
% cue 1409 3688.400 3689.620
\hypertarget{w02:cue:1409}{}殺了他，
% cue 1410 3689.620 3691.000
\hypertarget{w02:cue:1410}{}法治，
% cue 1411 3691.000 3692.000
\hypertarget{w02:cue:1411}{}法要有個道理，
% cue 1412 3692.000 3694.460
\hypertarget{w02:cue:1412}{}要 based on reason，
% cue 1413 3694.460 3696.050
\hypertarget{w02:cue:1413}{}你講法治，
% cue 1414 3696.050 3698.050
\hypertarget{w02:cue:1414}{}你卻說政府的決定是不 rational，
% cue 1415 3698.050 3699.040
\hypertarget{w02:cue:1415}{}這我們聽不懂，
% cue 1416 3699.040 3701.130
\hypertarget{w02:cue:1416}{}這是第一。

\tc{1:01:41}
% cue 1417 3701.130 3702.260
\hypertarget{w02:cue:1417}{}第二，
% cue 1418 3702.260 3703.610
\hypertarget{w02:cue:1418}{}第二，
% cue 1419 3703.610 3704.610
\hypertarget{w02:cue:1419}{}這是更難的，
% cue 1420 3704.610 3705.700
\hypertarget{w02:cue:1420}{}請注意這裡，
% cue 1421 3705.700 3706.800
\hypertarget{w02:cue:1421}{}the reasons，
% cue 1422 3706.800 3707.800
\hypertarget{w02:cue:1422}{}你說你有道理，
% cue 1423 3707.800 3708.730
\hypertarget{w02:cue:1423}{}你有理由，
% cue 1424 3708.730 3709.960
\hypertarget{w02:cue:1424}{}好，
% cue 1425 3709.960 3712.310
\hypertarget{w02:cue:1425}{}你所主張的理由 must be legal，
% cue 1426 3712.310 3713.750
\hypertarget{w02:cue:1426}{}必須是合理的，
% cue 1427 3713.750 3714.750
\hypertarget{w02:cue:1427}{}合法的理由，
% cue 1428 3714.750 3718.060
\hypertarget{w02:cue:1428}{}等一下會舉例子。

\tc{1:01:58}
% cue 1429 3718.060 3719.500
\hypertarget{w02:cue:1429}{}it follows that，
% cue 1430 3719.500 3721.880
\hypertarget{w02:cue:1430}{}special legitimacy，
% cue 1431 3721.880 3723.390
\hypertarget{w02:cue:1431}{}其實這個地方，
% cue 1432 3723.390 3725.480
\hypertarget{w02:cue:1432}{}如果我會把它改的話，
% cue 1433 3725.480 3728.670
\hypertarget{w02:cue:1433}{}這個 special 這個詞要把它改成 specific，
% cue 1434 3728.670 3730.050
\hypertarget{w02:cue:1434}{}這樣它的用語才會統一，
% cue 1435 3730.050 3733.530
\hypertarget{w02:cue:1435}{}specific legitimacy is considered，
% cue 1436 3733.530 3734.810
\hypertarget{w02:cue:1436}{}conditioned by rational，
% cue 1437 3734.810 3736.810
\hypertarget{w02:cue:1437}{}judicial decisions，
% cue 1438 3736.810 3740.380
\hypertarget{w02:cue:1438}{}是有一個合理的司法決定，
% cue 1439 3740.380 3742.600
\hypertarget{w02:cue:1439}{}that are based upon reason，
% cue 1440 3742.600 3745.200
\hypertarget{w02:cue:1440}{}and certain in a sense，
% cue 1441 3745.200 3746.200
\hypertarget{w02:cue:1441}{}it's legal。

\tc{1:02:26}
% cue 1442 3746.200 3747.870
\hypertarget{w02:cue:1442}{}legal 有兩種意思，
% cue 1443 3747.870 3748.960
\hypertarget{w02:cue:1443}{}有一種是表象的，
% cue 1444 3748.960 3751.180
\hypertarget{w02:cue:1444}{}有一種是實質的，
% cue 1445 3751.180 3752.470
\hypertarget{w02:cue:1445}{}那後面這一點，
% cue 1446 3752.470 3753.470
\hypertarget{w02:cue:1446}{}他等一下會講，
% cue 1447 3753.470 3754.820
\hypertarget{w02:cue:1447}{}我跟你講說，
% cue 1448 3754.820 3756.820
\hypertarget{w02:cue:1448}{}你依據的理由不合法，
% cue 1449 3756.820 3758.270
\hypertarget{w02:cue:1449}{}這個是什麼意思，
% cue 1450 3758.270 3759.270
\hypertarget{w02:cue:1450}{}法律明明講起了，
% cue 1451 3759.270 3761.270
\hypertarget{w02:cue:1451}{}你怎麼說我依據的理由不合法，
% cue 1452 3761.270 3762.320
\hypertarget{w02:cue:1452}{}所以這個 legal，
% cue 1453 3762.320 3764.620
\hypertarget{w02:cue:1453}{}它有比較複雜的意思。

\tc{1:02:44}
% cue 1454 3764.620 3765.780
\hypertarget{w02:cue:1454}{}再下面一段，
% cue 1455 3765.780 3769.830
\hypertarget{w02:cue:1455}{}the rule of law condition is a formal，
% cue 1456 3769.830 3773.420
\hypertarget{w02:cue:1456}{}有關於依法，
% cue 1457 3773.420 3775.600
\hypertarget{w02:cue:1457}{}這個條件是形式化的，
% cue 1458 3775.600 3776.600
\hypertarget{w02:cue:1458}{}it holds that，
% cue 1459 3776.600 3777.530
\hypertarget{w02:cue:1459}{}它主張，
% cue 1460 3777.530 3779.560
\hypertarget{w02:cue:1460}{}the legitimacy of the judicial power，
% cue 1461 3779.560 3783.390
\hypertarget{w02:cue:1461}{}partially derives from the legal character。
% cue 1462 3783.390 3786.570
\hypertarget{w02:cue:1462}{}of the justification of the decision，
% cue 1463 3786.570 3787.630
\hypertarget{w02:cue:1463}{}made by the courts，
% cue 1464 3787.630 3789.690
\hypertarget{w02:cue:1464}{}independently of their content，
% cue 1465 3789.690 3791.740
\hypertarget{w02:cue:1465}{}跟他們的內容不關，
% cue 1466 3791.740 3793.740
\hypertarget{w02:cue:1466}{}跟他們的本質價值，
% cue 1467 3793.740 3796.080
\hypertarget{w02:cue:1467}{}跟結果不關，
% cue 1468 3796.080 3797.080
\hypertarget{w02:cue:1468}{}這些部分是實質的，
% cue 1469 3797.080 3799.170
\hypertarget{w02:cue:1469}{}我現在講是形式賬本。

\tc{1:03:19}
% cue 1470 3799.170 3801.500
\hypertarget{w02:cue:1470}{}within modern legal theory，
% cue 1471 3801.500 3805.770
\hypertarget{w02:cue:1471}{}the conditions have often been conceived as sufficient，
% cue 1472 3805.770 3807.800
\hypertarget{w02:cue:1472}{}mere legality，
% cue 1473 3807.800 3810.140
\hypertarget{w02:cue:1473}{}只要合法就好了，
% cue 1474 3810.140 3811.490
\hypertarget{w02:cue:1474}{}這個就是 legalism。

\tc{1:03:31}
% cue 1475 3811.490 3813.150
\hypertarget{w02:cue:1475}{}yet，
% cue 1476 3813.150 3814.150
\hypertarget{w02:cue:1476}{}然而，
% cue 1477 3814.150 3819.270
\hypertarget{w02:cue:1477}{}mere legality cannot be accepted as sufficient condition，
% cue 1478 3819.270 3821.520
\hypertarget{w02:cue:1478}{}不能被認為是充分條件。

\tc{1:03:41}
% cue 1479 3821.520 3824.770
\hypertarget{w02:cue:1479}{}他說 imagine 下面這種情況，
% cue 1480 3824.770 3829.300
\hypertarget{w02:cue:1480}{}如果，你看下面 529 頁，
% cue 1481 3829.300 3832.610
\hypertarget{w02:cue:1481}{}Imagine the case in which a dictator，
% cue 1482 3832.610 3833.600
\hypertarget{w02:cue:1482}{}一個獨裁者，
% cue 1483 3833.600 3837.160
\hypertarget{w02:cue:1483}{}他用他的法律上的命令，
% cue 1484 3837.160 3841.620
\hypertarget{w02:cue:1484}{}說把所有藍眼睛六歲以下的女孩，
% cue 1485 3841.620 3843.580
\hypertarget{w02:cue:1485}{}一律處死，
% cue 1486 3843.580 3845.060
\hypertarget{w02:cue:1486}{}為什麼是藍眼，
% cue 1487 3845.060 3846.060
\hypertarget{w02:cue:1487}{}為什麼是女孩，
% cue 1488 3846.060 3847.060
\hypertarget{w02:cue:1488}{}為什麼是六歲，
% cue 1489 3847.060 3849.760
\hypertarget{w02:cue:1489}{}Other things being equal，
% cue 1490 3849.760 3852.110
\hypertarget{w02:cue:1490}{}假如其他情況都不變的話，
% cue 1491 3852.110 3854.520
\hypertarget{w02:cue:1491}{}no one could seriously claim，
% cue 1492 3854.520 3857.900
\hypertarget{w02:cue:1492}{}沒有人可以很嚴肅地主張，
% cue 1493 3857.900 3859.900
\hypertarget{w02:cue:1493}{}that the judges would enforce this command，
% cue 1494 3859.900 3863.590
\hypertarget{w02:cue:1494}{}你可以執行這個命令嗎？
% cue 1495 3863.590 3866.100
\hypertarget{w02:cue:1495}{}would act in a legitimate way，
% cue 1496 3866.100 3868.640
\hypertarget{w02:cue:1496}{}你可以說你這樣做，
% cue 1497 3868.640 3871.690
\hypertarget{w02:cue:1497}{}是一種正當的方式來行使你的司法權嗎？

\tc{1:04:31}
% cue 1498 3871.690 3874.230
\hypertarget{w02:cue:1498}{}one might certainly claim，
% cue 1499 3874.230 3878.450
\hypertarget{w02:cue:1499}{}that such judges would act within the law，
% cue 1500 3878.450 3879.410
\hypertarget{w02:cue:1500}{}你是不是依法，
% cue 1501 3879.410 3881.570
\hypertarget{w02:cue:1501}{}有啊，是有法律，
% cue 1502 3881.570 3882.670
\hypertarget{w02:cue:1502}{}according，
% cue 1503 3882.670 3884.670
\hypertarget{w02:cue:1503}{}and their decisions would be legal，
% cue 1504 3884.670 3885.670
\hypertarget{w02:cue:1504}{}所以從形式的觀點，
% cue 1505 3885.670 3887.670
\hypertarget{w02:cue:1505}{}它是合法的判決。

\tc{1:04:47}
% cue 1506 3887.670 3892.950
\hypertarget{w02:cue:1506}{}but to infer from this form of legality，
% cue 1507 3892.950 3899.110
\hypertarget{w02:cue:1507}{}從這一種單純的合法性的形式來推論，
% cue 1508 3899.110 3902.650
\hypertarget{w02:cue:1508}{}the proposition that the decision is legitimate，
% cue 1509 3902.650 3904.940
\hypertarget{w02:cue:1509}{}你因為形式上合法，
% cue 1510 3904.940 3907.030
\hypertarget{w02:cue:1510}{}你就說它是正當的，
% cue 1511 3907.030 3909.150
\hypertarget{w02:cue:1511}{}那我不能不說，
% cue 1512 3909.150 3911.150
\hypertarget{w02:cue:1512}{}這是一個 non sequitur\footnote{拉丁文，原意為「不隨之而來」，指結論不能由前提推出的推論。Tremblay 刊本 529 頁以此稱由裁判在形式上合法推出裁判為正當的推論。}，
% cue 1513 3911.150 3914.930
\hypertarget{w02:cue:1513}{}意思是說前後對，
% cue 1514 3914.930 3917.970
\hypertarget{w02:cue:1514}{}你單純形式合法，
% cue 1515 3917.970 3919.970
\hypertarget{w02:cue:1515}{}不能推論，
% cue 1516 3919.970 3920.930
\hypertarget{w02:cue:1516}{}它就是實正當，
% cue 1517 3920.930 3923.590
\hypertarget{w02:cue:1517}{}話講這樣子。

\tc{1:05:23}
% cue 1518 3923.590 3931.530
\hypertarget{w02:cue:1518}{}the judiciary has no more legitimacy than the law that justifies its decisions，
% cue 1519 3931.530 3935.100
\hypertarget{w02:cue:1519}{}thus, unless the law，
% cue 1520 3935.100 3938.300
\hypertarget{w02:cue:1520}{}除非這個法律被你聲稱，
% cue 1521 3938.300 3940.670
\hypertarget{w02:cue:1521}{}你主張作為理由的，
% cue 1522 3940.670 3942.490
\hypertarget{w02:cue:1522}{}有理由的這個法律，
% cue 1523 3943.370 3944.570
\hypertarget{w02:cue:1523}{}made legitimate，
% cue 1524 3944.570 3947.690
\hypertarget{w02:cue:1524}{}in accordance with some material，
% cue 1525 3947.690 3949.880
\hypertarget{w02:cue:1525}{}要有一些實質上的，
% cue 1526 3949.880 3951.550
\hypertarget{w02:cue:1526}{}material 可以，
% cue 1527 3951.550 3955.030
\hypertarget{w02:cue:1527}{}你也給它換成 substantive 也可以，
% cue 1528 3955.030 3957.320
\hypertarget{w02:cue:1528}{}實質上的道德標準，
% cue 1529 3957.320 3960.760
\hypertarget{w02:cue:1529}{}所以要以實質的道德標準來判斷，
% cue 1530 3960.760 3962.760
\hypertarget{w02:cue:1530}{}它到底是不是正當，
% cue 1531 3962.760 3965.860
\hypertarget{w02:cue:1531}{}而不能單純以形式來看，
% cue 1532 3965.860 3966.860
\hypertarget{w02:cue:1532}{}它是否合法。

\tc{1:06:06}
% cue 1533 3966.860 3968.510
\hypertarget{w02:cue:1533}{}好，
% cue 1534 3968.510 3969.760
\hypertarget{w02:cue:1534}{}所以，
% cue 1535 3969.760 3970.760
\hypertarget{w02:cue:1535}{}我們剛才到二，
% cue 1536 3970.760 3972.170
\hypertarget{w02:cue:1536}{}它說，
% cue 1537 3972.170 3975.960
\hypertarget{w02:cue:1537}{}形式的特定正當，
% cue 1538 3975.960 3977.440
\hypertarget{w02:cue:1538}{}特定，
% cue 1539 3977.440 3982.190
\hypertarget{w02:cue:1539}{}違憲審查特定正當性，
% cue 1540 3982.190 3985.980
\hypertarget{w02:cue:1540}{}作為條件，
% cue 1541 3985.980 3987.160
\hypertarget{w02:cue:1541}{}必須符合法治的要求，
% cue 1542 3987.160 3988.700
\hypertarget{w02:cue:1542}{}它說講到這裡，
% cue 1543 3988.700 3990.850
\hypertarget{w02:cue:1543}{}這個法治一般的理解，
% cue 1544 3990.850 3992.170
\hypertarget{w02:cue:1544}{}我剛剛說它有爭議，
% cue 1545 3992.170 3993.200
\hypertarget{w02:cue:1545}{}但一般的理解，
% cue 1546 3993.200 3994.230
\hypertarget{w02:cue:1546}{}這個法治，
% cue 1547 3994.230 3995.610
\hypertarget{w02:cue:1547}{}形式要求，
% cue 1548 3995.610 3998.050
\hypertarget{w02:cue:1548}{}它說單純形式要求不行，
% cue 1549 3998.050 4000.480
\hypertarget{w02:cue:1549}{}所以還要再加兩個命題。

\tc{1:06:40}
% cue 1550 4000.480 4001.480
\hypertarget{w02:cue:1550}{}看 529 頁，
% cue 1551 4001.480 4003.640
\hypertarget{w02:cue:1551}{}SLT3，
% cue 1552 4003.640 4010.300
\hypertarget{w02:cue:1552}{}The law providing the reason for judicial action and decision must itself be legitimate，
% cue 1553 4010.300 4012.550
\hypertarget{w02:cue:1553}{}你剛剛不是講法律嗎？
% cue 1554 4012.550 4014.550
\hypertarget{w02:cue:1554}{}法律提供了它正當的理由嗎？
% cue 1555 4014.550 4015.550
\hypertarget{w02:cue:1555}{}正當性的基礎嗎？
% cue 1556 4015.550 4016.640
\hypertarget{w02:cue:1556}{}不對，
% cue 1557 4016.640 4017.890
\hypertarget{w02:cue:1557}{}不夠，
% cue 1558 4017.890 4021.670
\hypertarget{w02:cue:1558}{}這個法律還必須本身是正當的。

\tc{1:07:01}
% cue 1559 4021.670 4024.670
\hypertarget{w02:cue:1559}{}Since the judiciary has no inherent legitimacy，
% cue 1560 4024.670 4025.630
\hypertarget{w02:cue:1560}{}剛講過，
% cue 1561 4025.630 4028.080
\hypertarget{w02:cue:1561}{}沒有本質上，
% cue 1562 4028.080 4031.990
\hypertarget{w02:cue:1562}{}Since its legitimacy is conditioned by rule of law，
% cue 1563 4031.990 4036.120
\hypertarget{w02:cue:1563}{}而且必須要以合乎法治作為條件，
% cue 1564 4036.120 4041.270
\hypertarget{w02:cue:1564}{}It has no more legitimacy than the law that governs its constitution，
% cue 1565 4041.270 4043.620
\hypertarget{w02:cue:1565}{}its process and its action。

\tc{1:07:23}
% cue 1566 4043.620 4044.720
\hypertarget{w02:cue:1566}{}好，
% cue 1567 4044.720 4047.780
\hypertarget{w02:cue:1567}{}這個是司法決策所依據的法律，
% cue 1568 4047.780 4050.780
\hypertarget{w02:cue:1568}{}因為你是依據那個法律取得的正當性，
% cue 1569 4050.780 4051.780
\hypertarget{w02:cue:1569}{}那我們就說，
% cue 1570 4051.780 4056.290
\hypertarget{w02:cue:1570}{}那那個法律本身也必須要 legitimate，
% cue 1571 4056.290 4059.760
\hypertarget{w02:cue:1571}{}the law be legitimate，
% cue 1572 4059.760 4063.650
\hypertarget{w02:cue:1572}{}那個 law 是你作為裁判之準據，
% cue 1573 4063.650 4066.130
\hypertarget{w02:cue:1573}{}裁判之準據的法律，
% cue 1574 4066.130 4069.280
\hypertarget{w02:cue:1574}{}那個法律本身也必須是合法，
% cue 1575 4069.280 4070.340
\hypertarget{w02:cue:1575}{}必須是正當的。

\tc{1:07:50}
% cue 1576 4070.340 4072.750
\hypertarget{w02:cue:1576}{}第四點最難的，
% cue 1577 4072.750 4076.700
\hypertarget{w02:cue:1577}{}the law is legitimate if and only if，
% cue 1578 4076.700 4079.490
\hypertarget{w02:cue:1578}{}在這裡你們可以跟我討論一個問題，
% cue 1579 4079.490 4082.890
\hypertarget{w02:cue:1579}{}老師那三跟四有什麼差別？
% cue 1580 4082.890 4084.800
\hypertarget{w02:cue:1580}{}三跟四應該合起來啊，
% cue 1581 4084.800 4088.830
\hypertarget{w02:cue:1581}{}為什麼要一個三再加一個四到這裡，
% cue 1582 4088.830 4090.340
\hypertarget{w02:cue:1582}{}如果你沒有這種法律，
% cue 1583 4090.340 4094.450
\hypertarget{w02:cue:1583}{}表示我認為這個三跟四，
% cue 1584 4094.450 4095.450
\hypertarget{w02:cue:1584}{}如果是我寫，
% cue 1585 4095.450 4096.510
\hypertarget{w02:cue:1585}{}我是合在一起的，
% cue 1586 4096.510 4097.510
\hypertarget{w02:cue:1586}{}但我不懂，
% cue 1587 4097.510 4099.790
\hypertarget{w02:cue:1587}{}他說準據的法律必須要是正當的，
% cue 1588 4099.790 4103.200
\hypertarget{w02:cue:1588}{}然後是講準據的法律，
% cue 1589 4103.200 4104.450
\hypertarget{w02:cue:1589}{}若且唯若，
% cue 1590 4104.450 4106.870
\hypertarget{w02:cue:1590}{}if and only if，
% cue 1591 4106.870 4108.900
\hypertarget{w02:cue:1591}{}在這種情況之下才是正當的，
% cue 1592 4108.900 4110.630
\hypertarget{w02:cue:1592}{}那就把它合在一起，
% cue 1593 4110.630 4111.630
\hypertarget{w02:cue:1593}{}不就好，
% cue 1594 4111.630 4113.260
\hypertarget{w02:cue:1594}{}不不，
% cue 1595 4113.260 4114.260
\hypertarget{w02:cue:1595}{}如果是我，
% cue 1596 4114.260 4115.260
\hypertarget{w02:cue:1596}{}是這麼寫的。

\tc{1:08:35}
% cue 1597 4115.260 4116.510
\hypertarget{w02:cue:1597}{}然後，
% cue 1598 4116.510 4118.510
\hypertarget{w02:cue:1598}{}it is constituted by political decisions，
% cue 1599 4118.510 4121.080
\hypertarget{w02:cue:1599}{}那這個法律是正當的，
% cue 1600 4121.080 4122.050
\hypertarget{w02:cue:1600}{}若且唯若，
% cue 1601 4122.050 4124.140
\hypertarget{w02:cue:1601}{}如果這個法律呢，
% cue 1602 4124.140 4126.420
\hypertarget{w02:cue:1602}{}是一個政治決定所建構的，
% cue 1603 4126.420 4128.540
\hypertarget{w02:cue:1603}{}而那個政治的決定呢，
% cue 1604 4128.540 4130.540
\hypertarget{w02:cue:1604}{}have a certain democratic pedigree，
% cue 1605 4130.540 4133.540
\hypertarget{w02:cue:1605}{}它具有某一種民主的血統，
% cue 1606 4133.540 4136.550
\hypertarget{w02:cue:1606}{}你們聽不懂嗎？
% cue 1607 4136.550 4139.050
\hypertarget{w02:cue:1607}{}剛剛說這個法律必須是正當的，
% cue 1608 4139.050 4142.080
\hypertarget{w02:cue:1608}{}然後說這個法律只有在這種情況下才是正當的，
% cue 1609 4142.080 4145.080
\hypertarget{w02:cue:1609}{}你說這個法律必須具有某種民主血統。

\tc{1:09:05}
% cue 1610 4145.670 4149.130
\hypertarget{w02:cue:1610}{}Indeed, in the abstract，抽象化來說。
% cue 1611 4149.130 4152.010
\hypertarget{w02:cue:1611}{}legitimacy of law，法的正當性。
% cue 1612 4152.010 4156.140
\hypertarget{w02:cue:1612}{}may not be a matter of democratic entitlement\footnote{英文，指依民主而取得的正當資格。Tremblay 刊本 529 頁謂，抽象而言法律的正當性未必取決於 democratic entitlement，也可能取決於是否合於自然法。}。
% cue 1613 4156.140 4161.540
\hypertarget{w02:cue:1613}{}可能跟它是不是民主上可以主張的。
% cue 1614 4161.540 4162.950
\hypertarget{w02:cue:1614}{}這個沒有必然的關係。
% cue 1615 4162.950 4170.730
\hypertarget{w02:cue:1615}{}但是在現代來講，它們經常地判斷，
% cue 1616 4170.730 4178.170
\hypertarget{w02:cue:1616}{}那些特別的正義性的行動，是最終的一件事。
% cue 1617 4178.170 4186.110
\hypertarget{w02:cue:1617}{}理論上來講，這個法律是不是正當，跟這個法律是不是民主，沒有必然關係。
% cue 1618 4186.110 4188.330
\hypertarget{w02:cue:1618}{}可以是不民主而正當。
% cue 1619 4188.330 4192.210
\hypertarget{w02:cue:1619}{}可是在今天這個民主社會。
% cue 1620 4192.210 4193.690
\hypertarget{w02:cue:1620}{}我們都會認為。
% cue 1621 4193.690 4195.550
\hypertarget{w02:cue:1621}{}那個法律如果是不民主。
% cue 1622 4195.550 4196.990
\hypertarget{w02:cue:1622}{}它怎麼會是正當呢。
% cue 1623 4196.990 4199.580
\hypertarget{w02:cue:1623}{}所以你要說這個法律要是正當的。
% cue 1624 4199.580 4202.180
\hypertarget{w02:cue:1624}{}它必須要具備某種民主的血統。
% cue 1625 4202.180 4203.840
\hypertarget{w02:cue:1625}{}有了這個民主的血統。
% cue 1626 4203.840 4205.220
\hypertarget{w02:cue:1626}{}你才可以主張。

\tc{1:10:05}
% cue 1627 4205.220 4209.480
\hypertarget{w02:cue:1627}{}民主的抗論真的很難翻。
% cue 1628 4209.480 4211.310
\hypertarget{w02:cue:1628}{}這個不知道要怎麼翻。
% cue 1629 4211.310 4214.270
\hypertarget{w02:cue:1629}{}這個叫我把它翻成一個加持。
% cue 1630 4214.270 4218.670
\hypertarget{w02:cue:1630}{}就是 entitlement。
% cue 1631 4218.670 4224.370
\hypertarget{w02:cue:1631}{}This has derived from a commitment to the separation of powers doctrine。
% cue 1632 4224.370 4227.040
\hypertarget{w02:cue:1632}{}這個是來自於權力分立的原則。
% cue 1633 4227.040 4229.810
\hypertarget{w02:cue:1633}{}權力分立的原則要怎麼解讀呢。
% cue 1634 4229.810 4231.830
\hypertarget{w02:cue:1634}{}read in accordance with。
% cue 1635 4231.830 4237.790
\hypertarget{w02:cue:1635}{}怎麼解讀要根據 the idea of popular or national sovereignty。
% cue 1636 4237.790 4240.760
\hypertarget{w02:cue:1636}{}就是國民主權原則\footnote{中華民國憲法第 2 條規定：「中華民國之主權屬於國民全體。」Tremblay 刊本 529 頁所述權力分立原則，係依 the ideas of popular and national sovereignty 解讀。\url{https://law.moj.gov.tw/LawClass/LawSingle.aspx?pcode=A0000001&flno=2}。}。
% cue 1637 4240.760 4244.710
\hypertarget{w02:cue:1637}{}that underlie the modern idea of democracy。

\tc{1:10:44}
% cue 1638 4244.710 4246.170
\hypertarget{w02:cue:1638}{}你看這句話多複雜。
% cue 1639 4246.170 4248.740
\hypertarget{w02:cue:1639}{}我們有一個現代的民主觀念。
% cue 1640 4248.740 4252.820
\hypertarget{w02:cue:1640}{}現在民主的觀念的基礎是國民主權原則。
% cue 1641 4252.820 4257.950
\hypertarget{w02:cue:1641}{}而國民主權原則為基礎來導讀來解讀權力分立。
% cue 1642 4257.950 4260.430
\hypertarget{w02:cue:1642}{}這樣子得出來的結論就是。
% cue 1643 4260.430 4261.780
\hypertarget{w02:cue:1643}{}你要立法。
% cue 1644 4261.780 4263.190
\hypertarget{w02:cue:1644}{}你這個法要正當。
% cue 1645 4263.190 4265.730
\hypertarget{w02:cue:1645}{}這個法必須具有民主的血統。
% cue 1646 4265.730 4267.770
\hypertarget{w02:cue:1646}{}如果你違反了民主。
% cue 1647 4267.770 4269.630
\hypertarget{w02:cue:1647}{}那你這個法可能是正當的。
% cue 1648 4269.630 4273.070
\hypertarget{w02:cue:1648}{}這麼複雜。

\tc{1:11:13}
% cue 1649 4273.070 4274.950
\hypertarget{w02:cue:1649}{}好了到這裡。
% cue 1650 4274.950 4277.360
\hypertarget{w02:cue:1650}{}經過我這樣的逐步的拆解。
% cue 1651 4277.360 4279.340
\hypertarget{w02:cue:1651}{}只剩下你同意跟不同意。
% cue 1652 4279.340 4280.830
\hypertarget{w02:cue:1652}{}但是你們大概不會聽不懂。
% cue 1653 4280.830 4282.560
\hypertarget{w02:cue:1653}{}好再往下。
% cue 1654 4282.560 4284.320
\hypertarget{w02:cue:1654}{}These four propositions。
% cue 1655 4284.320 4285.350
\hypertarget{w02:cue:1655}{}這四個命題。
% cue 1656 4285.390 4292.090
\hypertarget{w02:cue:1656}{}Express the modern justification for the specific legitimacy of a judicial action。
% cue 1657 4292.090 4298.930
\hypertarget{w02:cue:1657}{}你要講違憲審查的具體或者是特定的正當性。
% cue 1658 4298.930 4302.900
\hypertarget{w02:cue:1658}{}那它就包含了這四個剛才講的這四個命題。
% cue 1659 4302.900 4307.240
\hypertarget{w02:cue:1659}{}Indeed, they clearly show，他們清楚地顯示。
% cue 1660 4307.240 4314.490
\hypertarget{w02:cue:1660}{}Why Judicial enforcement of the democratically enacted legislation against citizens。
% cue 1661 4314.490 4317.580
\hypertarget{w02:cue:1661}{}has generally been conceived as a legitimate。
% cue 1662 4317.580 4319.540
\hypertarget{w02:cue:1662}{}within democratic society。

\tc{1:11:59}
% cue 1663 4319.540 4320.600
\hypertarget{w02:cue:1663}{}好。
% cue 1664 4320.600 4321.660
\hypertarget{w02:cue:1664}{}那他說。
% cue 1665 4321.660 4323.650
\hypertarget{w02:cue:1665}{}這樣你就可以解釋。
% cue 1666 4323.650 4326.530
\hypertarget{w02:cue:1666}{}為什麼由司法來執行。
% cue 1667 4326.530 4330.580
\hypertarget{w02:cue:1667}{}由合乎民主程序所制定的法律。
% cue 1668 4330.580 4333.360
\hypertarget{w02:cue:1668}{}對付 against its citizens。
% cue 1669 4333.360 4336.860
\hypertarget{w02:cue:1669}{}has generally been conceived as a legitimate。
% cue 1670 4336.860 4340.080
\hypertarget{w02:cue:1670}{}為什麼在一個民主社會裡。
% cue 1671 4340.080 4343.180
\hypertarget{w02:cue:1671}{}由司法來執行這一種。
% cue 1672 4343.180 4346.140
\hypertarget{w02:cue:1672}{}民意代表通過的法律。
% cue 1673 4346.140 4347.640
\hypertarget{w02:cue:1673}{}對抗它的人民。
% cue 1674 4347.640 4351.150
\hypertarget{w02:cue:1674}{}我們會認為具有正當性。
% cue 1675 4351.150 4370.890
\hypertarget{w02:cue:1675}{}因為它既是合法的。
% cue 1676 4370.890 4373.020
\hypertarget{w02:cue:1676}{}而且也是正當的。
% cue 1677 4373.020 4374.840
\hypertarget{w02:cue:1677}{}所以它來執行。
% cue 1678 4374.840 4375.850
\hypertarget{w02:cue:1678}{}還有什麼好說。
% cue 1679 4375.850 4379.200
\hypertarget{w02:cue:1679}{}這個是講執行。
% cue 1680 4379.200 4382.550
\hypertarget{w02:cue:1680}{}以上到這裡都是講核心。
% cue 1681 4382.550 4383.650
\hypertarget{w02:cue:1681}{}所以我就執行了。
% cue 1682 4383.650 4385.410
\hypertarget{w02:cue:1682}{}那我執行的結果。

\tc{1:13:05}
% cue 1683 4385.410 4390.630
\hypertarget{w02:cue:1683}{}但是現在，吾人本文比較關心的。
% cue 1684 4390.630 4393.630
\hypertarget{w02:cue:1684}{}But we are not so much concerned with。
% cue 1685 4393.630 4395.170
\hypertarget{w02:cue:1685}{}我們不是那麼關心。
% cue 1686 4395.170 4398.220
\hypertarget{w02:cue:1686}{}judicial decisions that enforce。
% cue 1687 4398.220 4401.470
\hypertarget{w02:cue:1687}{}democratically enacted legislation。
% cue 1688 4401.470 4407.040
\hypertarget{w02:cue:1688}{}As，沒有那麼關心，沒有像下面這個是一品。
% cue 1689 4407.040 4408.450
\hypertarget{w02:cue:1689}{}我們吾人更關心。
% cue 1690 4408.450 4411.190
\hypertarget{w02:cue:1690}{}judicial decisions that do not enforce it。
% cue 1691 4411.190 4413.530
\hypertarget{w02:cue:1691}{}司法做了一個決定。
% cue 1692 4413.530 4415.900
\hypertarget{w02:cue:1692}{}說我不執行立法者通過的法律。
% cue 1693 4415.900 4418.240
\hypertarget{w02:cue:1693}{}吾人關心的是這個。

\tc{1:13:38}
% cue 1694 4418.240 4437.640
\hypertarget{w02:cue:1694}{}When judges decide not to enforce against citizens a relevant piece of legislation that has been enacted by a democratic body, they do not recognize as good reason for their decision a legitimate law that should otherwise apply to the case at hand.
% cue 1695 4437.640 4443.780
\hypertarget{w02:cue:1695}{}This is why judicial review is generally conceived as countermajoritarian force within the system。

\tc{1:14:04}
% cue 1696 4444.110 4452.950
\hypertarget{w02:cue:1696}{}Moreover, even if majoritarianism did not constitute the best normative theory of democracy for the purpose of political legitimacy，
% cue 1697 4452.950 4461.290
\hypertarget{w02:cue:1697}{}judges could not justify the decision, not to enforce legislation by appeal to their own democratic pedigree.
% cue 1698 4461.290 4472.290
\hypertarget{w02:cue:1698}{}As we know, judges do not personify the people, do not represent and do not elect a not formally accountable electorate。

\tc{1:14:33}
% cue 1699 4473.150 4478.650
\hypertarget{w02:cue:1699}{}The question then is to determine the considerations，
% cue 1700 4478.650 4488.830
\hypertarget{w02:cue:1700}{}by virtue of which judicial review of democratically enacted legislation can be conceived and accepted as legitimate。

\section{以成文憲法為據的三種說法及其反對}

\tc{1:14:48}
% cue 1701 4488.830 4494.600
\hypertarget{w02:cue:1701}{}好，那這裡呢，為什麼違憲審查會是正當的呢？
% cue 1702 4494.600 4498.840
\hypertarget{w02:cue:1702}{}就有提出不同的觀點。
% cue 1703 4498.840 4503.220
\hypertarget{w02:cue:1703}{}Various answers have been put forward.
% cue 1704 4503.220 4504.220
\hypertarget{w02:cue:1704}{}以後。
% cue 1705 4504.220 4510.320
\hypertarget{w02:cue:1705}{}" 特定授權"。
% cue 1706 4510.320 4514.110
\hypertarget{w02:cue:1706}{}" 當"。
% cue 1707 4514.110 4519.540
\hypertarget{w02:cue:1707}{}" 以法律的法律為限制"。
% cue 1708 4519.540 4525.860
\hypertarget{w02:cue:1708}{}" 以法律為限制"。
% cue 1709 4525.860 4529.750
\hypertarget{w02:cue:1709}{}" 以法律為限制"。
% cue 1710 4529.750 4532.980
\hypertarget{w02:cue:1710}{}這是一種結論。
% cue 1711 4532.980 4535.500
\hypertarget{w02:cue:1711}{}那就是說你要不執行。
% cue 1712 4535.500 4538.960
\hypertarget{w02:cue:1712}{}你必須要有一特定的授權。
% cue 1713 4538.960 4543.170
\hypertarget{w02:cue:1713}{}或者特定的使他得到正當。

\tc{1:15:43}
% cue 1714 4543.170 4546.790
\hypertarget{w02:cue:1714}{}那這裡有幾種可能的答案。
% cue 1715 4546.790 4548.640
\hypertarget{w02:cue:1715}{}有。
% cue 1716 4548.640 4552.990
\hypertarget{w02:cue:1716}{}我是依照一個更高的成文憲法的授權。
% cue 1717 4552.990 4553.980
\hypertarget{w02:cue:1717}{}我做的。
% cue 1718 4553.980 4555.480
\hypertarget{w02:cue:1718}{}所以我是正當的。
% cue 1719 4555.480 4558.290
\hypertarget{w02:cue:1719}{}I shall call this thesis。
% cue 1720 4558.290 4560.340
\hypertarget{w02:cue:1720}{}我把這種命題叫做。
% cue 1721 4560.340 4561.600
\hypertarget{w02:cue:1721}{}The traditional view。
% cue 1722 4561.600 4562.740
\hypertarget{w02:cue:1722}{}傳統觀點。
% cue 1723 4562.740 4568.470
\hypertarget{w02:cue:1723}{}有關於違憲審查特定正當性的傳統觀點。
% cue 1724 4568.470 4570.080
\hypertarget{w02:cue:1724}{}傳統觀點。
% cue 1725 4570.080 4574.640
\hypertarget{w02:cue:1725}{}就是我依照一個更高的超乎立法的成文憲法。
% cue 1726 4574.640 4576.480
\hypertarget{w02:cue:1726}{}我來宣告。
% cue 1727 4576.480 4577.370
\hypertarget{w02:cue:1727}{}我不執行。
% cue 1728 4577.370 4578.870
\hypertarget{w02:cue:1728}{}所以我是依據他。
% cue 1729 4578.870 4580.190
\hypertarget{w02:cue:1729}{}所以我有正當性。

\tc{1:16:20}
% cue 1730 4580.190 4584.000
\hypertarget{w02:cue:1730}{}那下面我們要看的是呢。
% cue 1731 4584.000 4585.260
\hypertarget{w02:cue:1731}{}這點並不陌生。
% cue 1732 4585.260 4587.160
\hypertarget{w02:cue:1732}{}我們一般尋常都是這樣。
% cue 1733 4587.160 4588.190
\hypertarget{w02:cue:1733}{}是不是。
% cue 1734 4588.190 4589.150
\hypertarget{w02:cue:1734}{}我是依據憲法。
% cue 1735 4589.150 4590.110
\hypertarget{w02:cue:1735}{}所以我有正當性。
% cue 1736 4590.110 4592.700
\hypertarget{w02:cue:1736}{}但是作者顯然不是這麼認為。
% cue 1737 4592.700 4594.390
\hypertarget{w02:cue:1737}{}不是滿意於此。
% cue 1738 4594.390 4595.530
\hypertarget{w02:cue:1738}{}不滿意這個答案。

\tc{1:16:35}
% cue 1739 4595.530 4597.660
\hypertarget{w02:cue:1739}{}那我們現在他怎麼不滿意。
% cue 1740 4597.660 4599.390
\hypertarget{w02:cue:1740}{}130 頁下一段。
% cue 1741 4599.390 4602.070
\hypertarget{w02:cue:1741}{}The problem with the traditional view。
% cue 1742 4602.070 4606.180
\hypertarget{w02:cue:1742}{}這種傳統說法的第一個問題。
% cue 1743 4606.180 4611.330
\hypertarget{w02:cue:1743}{}是，它不需要的特定的立法律的法律。
% cue 1744 4611.330 4613.630
\hypertarget{w02:cue:1744}{}剛才說了有四個命題。
% cue 1745 4620.770 4623.990
\hypertarget{w02:cue:1745}{}但是第四個命題是沒有的，
% cue 1746 4623.990 4626.970
\hypertarget{w02:cue:1746}{}for this reason，
% cue 1747 4626.970 4629.490
\hypertarget{w02:cue:1747}{}the traditional view does not resolve。
% cue 1748 4629.490 4631.840
\hypertarget{w02:cue:1748}{}the anti-democratic difficulties。

\tc{1:17:11}
% cue 1749 4631.840 4633.820
\hypertarget{w02:cue:1749}{}第四個命題不是說，
% cue 1750 4633.820 4637.040
\hypertarget{w02:cue:1750}{}那個法律必須要符合，
% cue 1751 4637.040 4640.350
\hypertarget{w02:cue:1751}{}必須要有某一種民主的血統，
% cue 1752 4640.350 4644.790
\hypertarget{w02:cue:1752}{}才能制成那個法律是正當，
% cue 1753 4644.790 4645.790
\hypertarget{w02:cue:1753}{}而那個法律就是，
% cue 1754 4645.790 4649.820
\hypertarget{w02:cue:1754}{}違憲審查據以推翻普通法律的基礎，
% cue 1755 4649.820 4651.010
\hypertarget{w02:cue:1755}{}那你現在說，
% cue 1756 4651.010 4653.010
\hypertarget{w02:cue:1756}{}我是依照一個更高的憲法，
% cue 1757 4653.010 4656.830
\hypertarget{w02:cue:1757}{}那麼憲法有沒有民主正當性，
% cue 1758 4656.830 4661.030
\hypertarget{w02:cue:1758}{}憲法有沒有某一種民主的在裡面啊，
% cue 1759 4661.030 4662.250
\hypertarget{w02:cue:1759}{}不然那個法，
% cue 1760 4662.250 4663.250
\hypertarget{w02:cue:1760}{}憲法也是個法，
% cue 1761 4663.250 4665.210
\hypertarget{w02:cue:1761}{}憲法那個法就不符了，
% cue 1762 4665.210 4666.140
\hypertarget{w02:cue:1762}{}第四個，
% cue 1763 4666.140 4667.230
\hypertarget{w02:cue:1763}{}剛才講的面具。

\tc{1:17:47}
% cue 1764 4667.230 4672.550
\hypertarget{w02:cue:1764}{}下面，
% cue 1765 4672.550 4673.550
\hypertarget{w02:cue:1765}{}the traditional，
% cue 1766 4673.550 4674.550
\hypertarget{w02:cue:1766}{}往下這段的，
% cue 1767 4674.550 4676.550
\hypertarget{w02:cue:1767}{}這個段落的最後倒數幾句，
% cue 1768 4676.550 4677.680
\hypertarget{w02:cue:1768}{}倒數兩句，
% cue 1769 4677.680 4682.020
\hypertarget{w02:cue:1769}{}the traditional view is no more satisfying。
% cue 1770 4682.020 4686.940
\hypertarget{w02:cue:1770}{}than the form of legalism criticized above，
% cue 1771 4686.940 4689.000
\hypertarget{w02:cue:1771}{}如果你說，
% cue 1772 4689.000 4690.060
\hypertarget{w02:cue:1772}{}我是依法，
% cue 1773 4690.060 4691.220
\hypertarget{w02:cue:1773}{}我是依照一個形式的法，
% cue 1774 4691.220 4694.210
\hypertarget{w02:cue:1774}{}那我現在是依照一個形式的憲法，
% cue 1775 4694.210 4696.630
\hypertarget{w02:cue:1775}{}而依法那個法沒有民主正當性，
% cue 1776 4696.630 4698.820
\hypertarget{w02:cue:1776}{}這個憲法也沒有民主正當性，
% cue 1777 4698.820 4700.850
\hypertarget{w02:cue:1777}{}那你這個憲法比高那個法，
% cue 1778 4700.850 4704.630
\hypertarget{w02:cue:1778}{}你好嗎？
% cue 1779 4704.630 4706.140
\hypertarget{w02:cue:1779}{}聽懂嗎？

\tc{1:18:26}
% cue 1780 4706.140 4707.140
\hypertarget{w02:cue:1780}{}所以，
% cue 1781 4707.140 4708.180
\hypertarget{w02:cue:1781}{}為什麼他要強調這是 political morality，
% cue 1782 4708.180 4709.390
\hypertarget{w02:cue:1782}{}我們講正當性，
% cue 1783 4709.390 4710.670
\hypertarget{w02:cue:1783}{}政治道德，
% cue 1784 4710.670 4711.670
\hypertarget{w02:cue:1784}{}不是單純合乎，
% cue 1785 4711.670 4712.670
\hypertarget{w02:cue:1785}{}形式上合乎那個法的要件。

\tc{1:18:32}
% cue 1786 4712.670 4713.670
\hypertarget{w02:cue:1786}{}好，
% cue 1787 4713.670 4715.800
\hypertarget{w02:cue:1787}{}因為有這個理由呢，
% cue 1788 4715.800 4717.040
\hypertarget{w02:cue:1788}{}for this reason，
% cue 1789 4717.040 4720.520
\hypertarget{w02:cue:1789}{}a second strategy has been adopted，
% cue 1790 4720.520 4721.520
\hypertarget{w02:cue:1790}{}單純說我是依據憲法，
% cue 1791 4721.520 4724.280
\hypertarget{w02:cue:1791}{}不管那個憲法有沒有民主正當性，
% cue 1792 4724.280 4725.950
\hypertarget{w02:cue:1792}{}我都不論，
% cue 1793 4725.950 4726.870
\hypertarget{w02:cue:1793}{}這樣是不夠的，
% cue 1794 4726.870 4727.870
\hypertarget{w02:cue:1794}{}所以又有其他的說法要產生了，
% cue 1795 4727.870 4728.890
\hypertarget{w02:cue:1795}{}531 頁，
% cue 1796 4728.890 4730.810
\hypertarget{w02:cue:1796}{}那麼有關於這一點呢，
% cue 1797 4730.810 4732.260
\hypertarget{w02:cue:1797}{}還有好幾種不同的看法，
% cue 1798 4732.260 4733.260 刪：與 cue 1788 同文（連續 4 個重複解碼）
% cue 1799 4733.260 4735.480 刪：與 cue 1789 同文（連續 4 個重複解碼）
% cue 1800 4735.480 4737.890 刪：與 cue 1790 同文（連續 4 個重複解碼）
% cue 1801 4737.890 4738.890 刪：與 cue 1791 同文（連續 4 個重複解碼）
% cue 1802 4738.890 4739.920
\hypertarget{w02:cue:1802}{}我都不論，
% cue 1803 4739.920 4740.920
\hypertarget{w02:cue:1803}{}這樣是不夠的，
% cue 1804 4740.920 4744.380
\hypertarget{w02:cue:1804}{}所以又有其他的說法要產生了，
% cue 1805 4744.380 4745.450
\hypertarget{w02:cue:1805}{}531 頁。

\tc{1:19:05}
% cue 1806 4745.450 4749.320
\hypertarget{w02:cue:1806}{}我要慢慢的講這裡的原因是因為，
% cue 1807 4749.320 4751.320
\hypertarget{w02:cue:1807}{}我們一般人到第一層已經結束了，
% cue 1808 4751.320 4753.440
\hypertarget{w02:cue:1808}{}我也走到下面了，
% cue 1809 4753.440 4755.940
\hypertarget{w02:cue:1809}{}那你現在看它下面還有好幾層，
% cue 1810 4755.940 4757.410
\hypertarget{w02:cue:1810}{}531 頁，
% cue 1811 4757.410 4759.720
\hypertarget{w02:cue:1811}{}那麼有關於這一點呢，
% cue 1812 4759.720 4763.160
\hypertarget{w02:cue:1812}{}還有好幾種不同的看法。

\tc{1:19:23}
% cue 1813 4763.160 4769.250
\hypertarget{w02:cue:1813}{}most versions claim，
% cue 1814 4769.250 4773.880
\hypertarget{w02:cue:1814}{}大多數人主張的 constitutional democratic legitimacy，
% cue 1815 4773.880 4778.250
\hypertarget{w02:cue:1815}{}憲法的民主正當性 must be source based，
% cue 1816 4778.250 4779.470
\hypertarget{w02:cue:1816}{}一定是法源做基礎，
% cue 1817 4779.470 4780.470
\hypertarget{w02:cue:1817}{}來源做基礎，
% cue 1818 4780.470 4782.730
\hypertarget{w02:cue:1818}{}以來源為基礎，
% cue 1819 4782.730 4789.060
\hypertarget{w02:cue:1819}{}and derive from an original act of popular consent，
% cue 1820 4789.060 4795.690
\hypertarget{w02:cue:1820}{}而且是來自於某一個全民同意的一個原始行為，
% cue 1821 4795.690 4796.610
\hypertarget{w02:cue:1821}{}簡單講就是，
% cue 1822 4796.610 4801.570
\hypertarget{w02:cue:1822}{}制憲的時候得到全民同意的這樣一個基礎，
% cue 1823 4801.570 4803.890
\hypertarget{w02:cue:1823}{}來取得它的民主正當性。

\tc{1:20:03}
% cue 1824 4803.890 4805.520
\hypertarget{w02:cue:1824}{}剛才我不要忘記，
% cue 1825 4805.520 4806.520
\hypertarget{w02:cue:1825}{}我們剛才批評的是說，
% cue 1826 4806.520 4807.520
\hypertarget{w02:cue:1826}{}你不能講，
% cue 1827 4807.520 4808.520
\hypertarget{w02:cue:1827}{}你不能不探討，
% cue 1828 4808.520 4811.580
\hypertarget{w02:cue:1828}{}憲法本身有沒有民主正當性，
% cue 1829 4811.580 4814.060
\hypertarget{w02:cue:1829}{}不然你就不能滿足第四個要件，
% cue 1830 4814.060 4815.120
\hypertarget{w02:cue:1830}{}第四個要件說，
% cue 1831 4815.120 4817.730
\hypertarget{w02:cue:1831}{}司法所談的那個法律，
% cue 1832 4817.730 4819.730
\hypertarget{w02:cue:1832}{}法律本身要具有某種民主血統，
% cue 1833 4819.730 4821.750
\hypertarget{w02:cue:1833}{}就是某種民主正當性，
% cue 1834 4821.750 4824.000
\hypertarget{w02:cue:1834}{}你不能不討論這一點，
% cue 1835 4824.000 4825.890
\hypertarget{w02:cue:1835}{}那要討論這一點，
% cue 1836 4825.890 4828.970
\hypertarget{w02:cue:1836}{}大多數的人的做法就是說，
% cue 1837 4828.970 4832.230
\hypertarget{w02:cue:1837}{}那是因為我們原始就有這個授權，
% cue 1838 4832.230 4834.320
\hypertarget{w02:cue:1838}{}所以我就取得正當性。

\tc{1:20:34}
% cue 1839 4834.320 4836.470
\hypertarget{w02:cue:1839}{}那看下面，
% cue 1840 4836.470 4837.470
\hypertarget{w02:cue:1840}{}I shall call this argument，
% cue 1841 4837.470 4838.470
\hypertarget{w02:cue:1841}{}這一段的最後，
% cue 1842 4838.470 4842.370
\hypertarget{w02:cue:1842}{}the dominant view on the 什麼，
% cue 1843 4842.370 4844.050
\hypertarget{w02:cue:1843}{}那我們現在講，
% cue 1844 4844.050 4845.080
\hypertarget{w02:cue:1844}{}違憲審查，
% cue 1845 4845.080 4848.170
\hypertarget{w02:cue:1845}{}特定正當性的主流觀點，
% cue 1846 4848.170 4850.060
\hypertarget{w02:cue:1846}{}這是不是主流說，
% cue 1847 4850.060 4853.530
\hypertarget{w02:cue:1847}{}主流說是不是長這樣，
% cue 1848 4853.530 4854.880
\hypertarget{w02:cue:1848}{}你們所理解，
% cue 1849 4854.880 4855.880
\hypertarget{w02:cue:1849}{}你們所聽到的，
% cue 1850 4855.880 4857.880
\hypertarget{w02:cue:1850}{}違憲審查正當性的主流說，
% cue 1851 4857.880 4860.700
\hypertarget{w02:cue:1851}{}我依據憲法，
% cue 1852 4860.700 4861.700
\hypertarget{w02:cue:1852}{}不夠，
% cue 1853 4861.700 4863.180
\hypertarget{w02:cue:1853}{}我依據的這個憲法，
% cue 1854 4863.180 4865.240
\hypertarget{w02:cue:1854}{}是得到原始授權的，
% cue 1855 4865.240 4866.360
\hypertarget{w02:cue:1855}{}所以它具有民主正當性，
% cue 1856 4866.360 4867.650
\hypertarget{w02:cue:1856}{}講沒話講了吧，
% cue 1857 4867.650 4869.780
\hypertarget{w02:cue:1857}{}四個要件都符合了吧，
% cue 1858 4869.780 4871.970
\hypertarget{w02:cue:1858}{}這是主流說，
% cue 1859 4871.970 4872.960
\hypertarget{w02:cue:1859}{}OK。

\tc{1:21:12}
% cue 1860 4872.960 4875.110
\hypertarget{w02:cue:1860}{}好，下面看，
% cue 1861 4875.110 4876.110
\hypertarget{w02:cue:1861}{}the dominant view is well known，
% cue 1862 4876.110 4878.520
\hypertarget{w02:cue:1862}{}within American constitutional law，
% cue 1863 4878.520 4880.870
\hypertarget{w02:cue:1863}{}這個說法在美國憲法裡面，
% cue 1864 4880.870 4881.850
\hypertarget{w02:cue:1864}{}是廣為人知的，
% cue 1865 4881.850 4884.880
\hypertarget{w02:cue:1865}{}at least since the famous case，
% cue 1866 4884.880 4886.850
\hypertarget{w02:cue:1866}{}Marbury vs Madison\footnote{Marbury v. Madison, 5 U.S. 137，美國聯邦最高法院 1803 年 2 月 24 日判決。首席大法官 John Marshall 執筆的法院意見謂「an act of the legislature, repugnant to the constitution, is void」，又謂「It is emphatically the province and duty of the judicial department to say what the law is.」美國國家檔案館將本判決列為確立違憲審查原則的判決。\url{https://www.archives.gov/milestone-documents/marbury-v-madison}。}。

\tc{1:21:26}
% cue 1867 4886.850 4888.390
\hypertarget{w02:cue:1867}{}那下面就有一段，
% cue 1868 4888.390 4889.390
\hypertarget{w02:cue:1868}{}比較詳細的論證了，
% cue 1869 4889.390 4891.740
\hypertarget{w02:cue:1869}{}簡單的講就是說，
% cue 1870 4891.740 4893.900
\hypertarget{w02:cue:1870}{}Marbury vs Madison 裡面的，
% cue 1871 4893.900 4895.990
\hypertarget{w02:cue:1871}{}Chief Justice Marshall\footnote{John Marshall（1755–1835），由 John Adams 總統提名，1801 年 2 月 4 日至 1835 年 7 月 6 日任美國聯邦最高法院首席大法官。\url{https://www.supremecourt.gov/about/members_text.aspx}。}，
% cue 1872 4895.990 4896.980
\hypertarget{w02:cue:1872}{}他說，
% cue 1873 4896.980 4898.050
\hypertarget{w02:cue:1873}{}法律抵觸憲法，
% cue 1874 4898.050 4899.460
\hypertarget{w02:cue:1874}{}應該無效，
% cue 1875 4899.460 4900.430
\hypertarget{w02:cue:1875}{}我們先不管說，
% cue 1876 4900.430 4902.430
\hypertarget{w02:cue:1876}{}是不是應該由法院來認定，
% cue 1877 4902.430 4904.400
\hypertarget{w02:cue:1877}{}這件事情是有爭議的，
% cue 1878 4904.400 4906.400
\hypertarget{w02:cue:1878}{}因為美國憲法也沒有明文，
% cue 1879 4906.400 4907.460
\hypertarget{w02:cue:1879}{}但是他的主要論證說，
% cue 1880 4907.460 4909.590
\hypertarget{w02:cue:1880}{}我們這個國家有法，
% cue 1881 4909.590 4910.590
\hypertarget{w02:cue:1881}{}那又有憲法，
% cue 1882 4910.590 4912.770
\hypertarget{w02:cue:1882}{}憲法說憲法是這個國家最高的法，
% cue 1883 4912.770 4914.850
\hypertarget{w02:cue:1883}{}這個憲法的效力高於法，
% cue 1884 4914.850 4917.270
\hypertarget{w02:cue:1884}{}所以但下面的法律，
% cue 1885 4917.270 4918.270
\hypertarget{w02:cue:1885}{}如果抵觸了憲法，
% cue 1886 4918.270 4920.230
\hypertarget{w02:cue:1886}{}那我們法院當然只能認定，
% cue 1887 4920.230 4921.230
\hypertarget{w02:cue:1887}{}效力優先，
% cue 1888 4921.230 4922.230
\hypertarget{w02:cue:1888}{}效力居高的，
% cue 1889 4922.230 4923.230
\hypertarget{w02:cue:1889}{}位階居高，
% cue 1890 4923.230 4924.230
\hypertarget{w02:cue:1890}{}效力優先的優先，
% cue 1891 4924.230 4925.900
\hypertarget{w02:cue:1891}{}三點。

\tc{1:22:05}
% cue 1892 4925.900 4927.060
\hypertarget{w02:cue:1892}{}第二點，
% cue 1893 4927.060 4928.060
\hypertarget{w02:cue:1893}{}他也會講說，
% cue 1894 4928.060 4929.060
\hypertarget{w02:cue:1894}{}不要忘記了，
% cue 1895 4929.060 4931.060
\hypertarget{w02:cue:1895}{}我們這個憲法從哪裡來的，
% cue 1896 4931.060 4933.020
\hypertarget{w02:cue:1896}{}我們是得到人民授權的，
% cue 1897 4933.020 4934.020
\hypertarget{w02:cue:1897}{}那是制憲的時候通過的，
% cue 1898 4934.020 4936.590
\hypertarget{w02:cue:1898}{}所以我當然有正當性，
% cue 1899 4936.590 4938.940
\hypertarget{w02:cue:1899}{}他用這樣的方法，
% cue 1900 4938.940 4940.030
\hypertarget{w02:cue:1900}{}他滿足了，
% cue 1901 4940.030 4941.290
\hypertarget{w02:cue:1901}{}他為什麼，
% cue 1902 4941.290 4942.830
\hypertarget{w02:cue:1902}{}他為什麼認為，
% cue 1903 4942.830 4943.890
\hypertarget{w02:cue:1903}{}他有正當性，
% cue 1904 4943.890 4945.890
\hypertarget{w02:cue:1904}{}可以去把違反憲法的法律，
% cue 1905 4945.890 4948.640
\hypertarget{w02:cue:1905}{}學校威脅不小。

\tc{1:22:28}
% cue 1906 4948.640 4950.640
\hypertarget{w02:cue:1906}{}然後 532 頁，
% cue 1907 4950.640 4951.890
\hypertarget{w02:cue:1907}{}這種主流觀點，
% cue 1908 4951.890 4954.950
\hypertarget{w02:cue:1908}{}也得到了加拿大最高法院\footnote{Supreme Court of Canada。Tremblay 刊本 532 頁引其判決，謂將權利憲章納入憲法的歷史性決定由加拿大人民選出的代表作成，並謂「Adjudication under the Charter must be approached free of any lingering doubts as to its legitimacy.」}的認識，
% cue 1909 4954.950 4956.390
\hypertarget{w02:cue:1909}{}加拿大最高法院，
% cue 1910 4956.390 4958.370
\hypertarget{w02:cue:1910}{}亦採此說，
% cue 1911 4958.370 4960.350
\hypertarget{w02:cue:1911}{}那這一段就比較複雜，
% cue 1912 4960.350 4961.350
\hypertarget{w02:cue:1912}{}牽涉到加拿大的憲法，
% cue 1913 4961.350 4964.730
\hypertarget{w02:cue:1913}{}原來是英國通過了一個法律，
% cue 1914 4964.730 4967.850
\hypertarget{w02:cue:1914}{}給北美洲民帝東，
% cue 1915 4967.850 4969.050
\hypertarget{w02:cue:1915}{}然後來呢，
% cue 1916 4969.050 4970.430
\hypertarget{w02:cue:1916}{}英國沒落了嗎，
% cue 1917 4970.430 4972.330
\hypertarget{w02:cue:1917}{}加拿大 1982\footnote{英國國會於 1982 年制定 Canada Act 1982，終止其對加拿大的立法權；該法附表 B 為 Constitution Act, 1982，內含加拿大權利與自由憲章及修憲程序。Constitution Act, 1982 第 52 條第 1 項規定加拿大憲法為加拿大的最高法，與之牴觸的法律於牴觸範圍內無效。1867 年由英國國會制定、據以成立加拿大的 British North America Act，現稱 Constitution Act, 1867。\url{https://laws-lois.justice.gc.ca/eng/const/page-13.html}。} 說，
% cue 1918 4972.330 4973.410
\hypertarget{w02:cue:1918}{}我可以不獨立，
% cue 1919 4973.410 4975.480
\hypertarget{w02:cue:1919}{}我可以不要脫離英國國協，
% cue 1920 4975.480 4977.540
\hypertarget{w02:cue:1920}{}但是我要了自己的憲法，
% cue 1921 4977.540 4979.700
\hypertarget{w02:cue:1921}{}所以他在立法技術上，
% cue 1922 4979.700 4981.600
\hypertarget{w02:cue:1922}{}他就做了一些轉折，
% cue 1923 4981.600 4983.760
\hypertarget{w02:cue:1923}{}這些轉折就牽涉到下面這一段話，
% cue 1924 4983.760 4984.760
\hypertarget{w02:cue:1924}{}這個事情比較複雜，
% cue 1925 4984.760 4985.760
\hypertarget{w02:cue:1925}{}我就不去討論。

\tc{1:23:05}
% cue 1926 4985.760 4986.760
\hypertarget{w02:cue:1926}{}好，
% cue 1927 4986.760 4987.760
\hypertarget{w02:cue:1927}{}現在看 533 頁，
% cue 1928 4987.760 4988.760
\hypertarget{w02:cue:1928}{}重點是，
% cue 1929 4988.760 4990.720
\hypertarget{w02:cue:1929}{}這樣的民主正當性，
% cue 1930 4990.720 4991.880
\hypertarget{w02:cue:1930}{}同意嗎，
% cue 1931 4991.880 4993.740
\hypertarget{w02:cue:1931}{}你同意嗎，
% cue 1932 4993.740 4994.740
\hypertarget{w02:cue:1932}{}533 頁第一段話，
% cue 1933 4994.740 4995.740
\hypertarget{w02:cue:1933}{}but there is，
% cue 1934 4995.740 4997.060
\hypertarget{w02:cue:1934}{}那個。

\tc{1:23:17}
% cue 1935 4997.060 4997.970
\hypertarget{w02:cue:1935}{}鄭偉現在幾點？

\tc{1:23:17}
\speaker{學生}
% cue 1936 4997.970 4999.060
\hypertarget{w02:cue:1936}{}呃，
% cue 1937 4999.060 5000.060
\hypertarget{w02:cue:1937}{}現在 11.42。

\tc{1:23:20}
\speaker{老師}
% cue 1938 5000.060 5002.240
\hypertarget{w02:cue:1938}{}好，
% cue 1939 5002.240 5004.320
\hypertarget{w02:cue:1939}{}但是，
% cue 1940 5004.320 5006.320
\hypertarget{w02:cue:1940}{}there is at least one important objection，
% cue 1941 5006.320 5007.570
\hypertarget{w02:cue:1941}{}對於主流說，
% cue 1942 5007.570 5009.570
\hypertarget{w02:cue:1942}{}至少有一點反對意見，
% cue 1943 5009.570 5030.310
\hypertarget{w02:cue:1943}{}The mere fact that the constitution was drafted, promulgated, ratified, and otherwise, explicitly or tacitly approved by the will of the majority of people, or their elected representatives, at the time of enactment，把它圈起來，在制憲的時候，確實得到了多數人民的同意，批准的。
% cue 1944 5030.310 5040.540
\hypertarget{w02:cue:1944}{}這個事實 does not necessarily entail，並不當然導出，it ought to prevail over subsequent legislation。

\tc{1:24:00}
% cue 1945 5040.540 5048.960
\hypertarget{w02:cue:1945}{}其他以後制定的法律，難道不是人民多數代表人民所同意的嗎？
% cue 1946 5048.960 5055.720
\hypertarget{w02:cue:1946}{}那為什麼你原先那個憲法就要效力高於以後的法律呢？
% cue 1947 5055.720 5062.700
\hypertarget{w02:cue:1947}{}這件事情在臺灣，我也很高興看到今天的這段，我們在臺灣幾乎沒有人討論。
% cue 1948 5062.700 5070.760
\hypertarget{w02:cue:1948}{}我們臺灣人的頭腦就是，憲法效力高於法律，句號，所以法律抵觸憲法者無效，句號。
% cue 1949 5070.760 5076.970
\hypertarget{w02:cue:1949}{}但是你忘記了一件事情，如果從民主正當性的觀點，要看誰能反映民主正當性，
% cue 1950 5076.970 5083.050
\hypertarget{w02:cue:1950}{}誰能反映民主，誰才具有正當性，誰才有民主正當性，是不是？
% cue 1951 5083.050 5088.290
\hypertarget{w02:cue:1951}{}那你現在的憲法是 1776 年之廳，我們這個法律是 2026 年之通過的，
% cue 1952 5088.290 5094.010
\hypertarget{w02:cue:1952}{}那哪一個法律比較代表民意，哪個法律比較具有正當性？
% cue 1953 5094.010 5101.130
\hypertarget{w02:cue:1953}{}為什麼是 1776 年\footnote{美國獨立宣言於 1776 年 7 月 4 日通過；美國聯邦憲法於 1787 年 9 月 17 日由制憲會議簽署。\url{https://www.archives.gov/founding-docs/declaration-history}。}這些作古的人，他們憑什麼拘束 30 後，250 後的我們憑什麼？
% cue 1954 5101.130 5103.470
\hypertarget{w02:cue:1954}{}你都等我講到個意思嗎？

\tc{1:25:03}
% cue 1955 5103.470 5110.220
\hypertarget{w02:cue:1955}{}你說說民主，那我才是民主啊，我自己要統治自己啊，不是以前人統治我啊，
% cue 1956 5110.220 5114.560
\hypertarget{w02:cue:1956}{}以前人當時他們是民主，我沒有話講，但是他們當時是民主，
% cue 1957 5114.560 5117.560
\hypertarget{w02:cue:1957}{}並不表示他們當時的民主要拘束我們今天的民主，
% cue 1958 5117.560 5122.410
\hypertarget{w02:cue:1958}{}那如果是叫到我們自己人的白活了，怎麼會叫呢？

\tc{1:25:22}
% cue 1959 5122.410 5127.290
\hypertarget{w02:cue:1959}{}下面再往下看。
% cue 1960 5127.290 5129.610
\hypertarget{w02:cue:1960}{}這種方法是說，我們不能夠找到我們的法律，
% cue 1961 5129.610 5130.610
\hypertarget{w02:cue:1961}{}所以說，我們要找到我們的法律，
% cue 1962 5130.610 5131.610
\hypertarget{w02:cue:1962}{}所以說，我們要找到我們的法律，
% cue 1963 5131.610 5133.800
\hypertarget{w02:cue:1963}{}我們要找到我們的法律，
% cue 1964 5133.800 5135.890
\hypertarget{w02:cue:1964}{}所以說，我們要找到我們的法律，
% cue 1965 5145.270 5145.830
\hypertarget{w02:cue:1965}{}請的推理。

\tc{1:25:45}
% cue 1966 5145.830 5161.200
\hypertarget{w02:cue:1966}{}當法院做成決定的時候。
% cue 1967 5158.290 5160.470
\hypertarget{w02:cue:1967}{}當法院做成決定的時候，
% cue 1968 5160.470 5162.620
\hypertarget{w02:cue:1968}{}法院法官應該問，
% cue 1969 5162.620 5165.830
\hypertarget{w02:cue:1969}{}是通過在後的法律，
% cue 1970 5165.830 5169.040
\hypertarget{w02:cue:1970}{}還是 250 年前通過的憲法，
% cue 1971 5169.040 5172.610
\hypertarget{w02:cue:1971}{}誰才最能代表，
% cue 1972 5172.610 5174.750
\hypertarget{w02:cue:1972}{}the best represent，
% cue 1973 5174.750 5177.040
\hypertarget{w02:cue:1973}{}最能代表人民的同意，
% cue 1974 5177.040 5179.890
\hypertarget{w02:cue:1974}{}誰能代表人民，
% cue 1975 5179.890 5182.500
\hypertarget{w02:cue:1975}{}誰最能代表人民的同意，
% cue 1976 5182.500 5184.750
\hypertarget{w02:cue:1976}{}那我要以這個為準。

\tc{1:26:24}
% cue 1977 5184.750 5196.340
\hypertarget{w02:cue:1977}{}the mere fact that the democratically enacted constitution has been recognized as a supra-legislative within a formal hierarchy of legal sources.
% cue 1978 5196.340 5197.400
\hypertarget{w02:cue:1978}{}它有一個 hierarchy，
% cue 1979 5197.400 5199.460
\hypertarget{w02:cue:1979}{}它有一個位階，
% cue 1980 5199.460 5200.520
\hypertarget{w02:cue:1980}{}它有一個位階，
% cue 1981 5200.520 5203.520
\hypertarget{w02:cue:1981}{}還像我們這個 Hans Kelsen\footnote{Hans Kelsen（1881–1973），奧地利法學家，純粹法學（Reine Rechtslehre）的創立者，為奧地利 1920 年聯邦憲法的主要起草人，該憲法所設的憲法法院（Verfassungsgerichtshof）出自其設計。課綱第 5 週（10 月 12 日）指定其 1929 年論文 Wesen und Entwicklung der Staatsgerichtsbarkeit 的英譯。\cite{kelsen2015}。} 的 Stufen theory\footnote{指法秩序階層構造說（Stufenbaulehre）：法秩序由上下位階的規範構成，下位規範的效力來自上位規範的授權，最上層為基礎規範（Grundnorm）。此說由 Kelsen 的學生 Adolf Julius Merkl 自 1918 年起提出，Kelsen 納入純粹法學。\url{https://resolve.cambridge.org/core/journals/canadian-journal-of-law-and-jurisprudence/article/problems-of-the-stufenbaulehre-kelsens-failure-to-derive-the-validity-of-a-norm-from-another-norm/B3F1B751D0BE4DB2CE7A271FD3CA16C5}。}，
% cue 1982 5203.520 5206.610
\hypertarget{w02:cue:1982}{}這個法規範階層構造說，
% cue 1983 5206.610 5207.610
\hypertarget{w02:cue:1983}{}你雖然站在最上面，
% cue 1984 5207.610 5209.890
\hypertarget{w02:cue:1984}{}這個是形式，
% cue 1985 5209.890 5210.890
\hypertarget{w02:cue:1985}{}cannot be sufficient reason，
% cue 1986 5210.890 5212.890
\hypertarget{w02:cue:1986}{}不足以作為一個充分的理由，
% cue 1987 5212.890 5215.930
\hypertarget{w02:cue:1987}{}to set aside the law that represent the better people，
% cue 1988 5215.930 5219.000
\hypertarget{w02:cue:1988}{}最能代表人民民意的法律，
% cue 1989 5219.000 5221.000
\hypertarget{w02:cue:1989}{}不因為它的位階在你之下，
% cue 1990 5221.000 5223.220
\hypertarget{w02:cue:1990}{}就被你幹掉，
% cue 1991 5223.220 5224.730
\hypertarget{w02:cue:1991}{}就是這樣的。

\tc{1:27:04}
% cue 1992 5224.730 5226.240
\hypertarget{w02:cue:1992}{}好，
% cue 1993 5226.240 5227.590
\hypertarget{w02:cue:1993}{}這個是第一個質疑，
% cue 1994 5227.590 5232.040
\hypertarget{w02:cue:1994}{}然後再往下一段，
% cue 1995 5232.040 5234.330
\hypertarget{w02:cue:1995}{}this objection to judicial review，
% cue 1996 5234.330 5238.440
\hypertarget{w02:cue:1996}{}is typical within sovereignty theory，
% cue 1997 5238.440 5244.920
\hypertarget{w02:cue:1997}{}這樣的一個對於違憲審查的質疑，
% cue 1998 5244.920 5249.730
\hypertarget{w02:cue:1998}{}也可以說在所謂的國民主權理論裡面，
% cue 1999 5249.730 5252.230
\hypertarget{w02:cue:1999}{}也是它的一部分，
% cue 2000 5252.230 5253.820
\hypertarget{w02:cue:2000}{}也是很常見的，
% cue 2001 5253.820 5256.670
\hypertarget{w02:cue:2001}{}the power of a sovereign，
% cue 2002 5256.670 5257.660
\hypertarget{w02:cue:2002}{}一個主權，
% cue 2003 5257.660 5259.730
\hypertarget{w02:cue:2003}{}主權的力量，
% cue 2004 5259.730 5262.500
\hypertarget{w02:cue:2004}{}has been conceived as continuous，
% cue 2005 5262.500 5266.010
\hypertarget{w02:cue:2005}{}主權被認為是持續的，
% cue 2006 5266.010 5268.090
\hypertarget{w02:cue:2006}{}okay，
% cue 2007 5268.090 5272.470
\hypertarget{w02:cue:2007}{}a sovereign people should never be bound by its previous political decision。

\tc{1:27:52}
% cue 2008 5272.470 5275.040
\hypertarget{w02:cue:2008}{}如果我是主權人民，
% cue 2009 5275.040 5280.760
\hypertarget{w02:cue:2009}{}我就不應該受到先前的政治決定的拘束，
% cue 2010 5280.760 5283.020
\hypertarget{w02:cue:2010}{}否則我還叫有主權嗎？
% cue 2011 5283.020 5284.050
\hypertarget{w02:cue:2011}{}你是主權者，
% cue 2012 5284.050 5285.050
\hypertarget{w02:cue:2012}{}你不能自己做決定，
% cue 2013 5285.050 5286.170
\hypertarget{w02:cue:2013}{}你要聽你爸爸，
% cue 2014 5286.170 5288.170
\hypertarget{w02:cue:2014}{}那你爸爸才是主權者，
% cue 2015 5288.170 5289.760
\hypertarget{w02:cue:2015}{}聽爺爺呢，
% cue 2016 5289.760 5290.740
\hypertarget{w02:cue:2016}{}聽三十代以前的主線，
% cue 2017 5290.740 5291.830
\hypertarget{w02:cue:2017}{}那你叫什麼主權？
% cue 2018 5291.830 5294.840
\hypertarget{w02:cue:2018}{}我是主權者，
% cue 2019 5294.840 5296.090
\hypertarget{w02:cue:2019}{}所以我說了算，
% cue 2020 5296.090 5297.110
\hypertarget{w02:cue:2020}{}我可以自己決定，
% cue 2021 5297.110 5298.360
\hypertarget{w02:cue:2021}{}我可以自為統治，
% cue 2022 5298.360 5300.100
\hypertarget{w02:cue:2022}{}不然我就不是主權者。

\tc{1:28:20}
% cue 2023 5300.100 5301.480
\hypertarget{w02:cue:2023}{}okay，
% cue 2024 5301.480 5306.840
\hypertarget{w02:cue:2024}{}the sovereign people should never be bound by its previous political decisions，
% cue 2025 5306.840 5309.060
\hypertarget{w02:cue:2025}{}任何先前的決定不能拘束，
% cue 2026 5309.060 5310.060
\hypertarget{w02:cue:2026}{}indeed，
% cue 2027 5310.060 5316.140
\hypertarget{w02:cue:2027}{}this has been a constitutive proposition of the doctrine of parliamentary sovereignty\footnote{英文，國會主權。Tremblay 刊本 533 頁註 18 引 A. V. Dicey, Introduction to the Study of the Law of the Constitution；依刊本 548 頁註 46 所引，Dicey 以國會主權指國會依英國憲法「the right to make or unmake any law whatever」，且英國法不承認任何人或機構有權推翻或擱置國會的立法。} in the United Kingdom，
% cue 2028 5316.140 5320.560
\hypertarget{w02:cue:2028}{}這個也是英國所謂的國會主權的理解。

\tc{1:28:40}
% cue 2029 5320.560 5322.120
\hypertarget{w02:cue:2029}{}好了，
% cue 2030 5322.120 5323.120
\hypertarget{w02:cue:2030}{}那這裡我們都跳過，
% cue 2031 5323.120 5324.120
\hypertarget{w02:cue:2031}{}在這裡面呢，
% cue 2032 5324.120 5326.120
\hypertarget{w02:cue:2032}{}它下面要講的就是，
% cue 2033 5326.120 5327.120
\hypertarget{w02:cue:2033}{}如果照這樣的觀點來講，
% cue 2034 5327.120 5329.180
\hypertarget{w02:cue:2034}{}我們可以理解後法優於前法，
% cue 2035 5329.180 5330.080
\hypertarget{w02:cue:2035}{}是不是？
% cue 2036 5330.080 5331.850
\hypertarget{w02:cue:2036}{}我們可以優於，
% cue 2037 5331.850 5333.850
\hypertarget{w02:cue:2037}{}我們可以理解特別法，
% cue 2038 5333.850 5334.850
\hypertarget{w02:cue:2038}{}優於普通法，
% cue 2039 5334.850 5338.130
\hypertarget{w02:cue:2039}{}基本上是要看它的民意正當性，
% cue 2040 5338.130 5340.320
\hypertarget{w02:cue:2040}{}它的民意基礎來決定，
% cue 2041 5340.320 5342.820
\hypertarget{w02:cue:2041}{}所以你不能跟我講說，
% cue 2042 5342.820 5343.820
\hypertarget{w02:cue:2042}{}我是憲法，
% cue 2043 5343.820 5346.870
\hypertarget{w02:cue:2043}{}一定當然在民意上要優先於法律，
% cue 2044 5346.870 5349.950
\hypertarget{w02:cue:2044}{}我們要就民意論民意，
% cue 2045 5349.950 5351.950
\hypertarget{w02:cue:2045}{}不是就抽象位階來論，
% cue 2046 5351.950 5355.190
\hypertarget{w02:cue:2046}{}這是他的意見。

\tc{1:29:15}
% cue 2047 5355.190 5356.730
\hypertarget{w02:cue:2047}{}那你再往下看，
% cue 2048 5356.730 5358.890
\hypertarget{w02:cue:2048}{}23 頁中間那一段，
% cue 2049 5358.890 5360.990
\hypertarget{w02:cue:2049}{}Bruce Ackerman 上場了，
% cue 2050 5360.990 5362.850
\hypertarget{w02:cue:2050}{}好，
% cue 2051 5362.850 5364.130
\hypertarget{w02:cue:2051}{}比如說，
% cue 2052 5364.130 5365.130
\hypertarget{w02:cue:2052}{}Bruce Ackerman has explained，
% cue 2053 5365.130 5367.250
\hypertarget{w02:cue:2053}{}the democratic superiority，
% cue 2054 5367.250 5369.500
\hypertarget{w02:cue:2054}{}民主的更高性，
% cue 2055 5369.500 5372.360
\hypertarget{w02:cue:2055}{}民主的優越性，
% cue 2056 5372.360 5374.640
\hypertarget{w02:cue:2056}{}of American Constitution，
% cue 2057 5374.640 5379.880
\hypertarget{w02:cue:2057}{}by reference to the quality of the deliberative processes，
% cue 2058 5379.880 5383.150
\hypertarget{w02:cue:2058}{}that have led to its popular ratification，
% cue 2059 5383.150 5387.560
\hypertarget{w02:cue:2059}{}這個牽涉到你們對 Bruce Ackerman 的理論的理解。

\tc{1:29:47}
% cue 2060 5387.560 5388.560
\hypertarget{w02:cue:2060}{}Bruce Ackerman 說，
% cue 2061 5388.560 5391.150
\hypertarget{w02:cue:2061}{}我們有兩個軌道立法，
% cue 2062 5391.150 5393.530
\hypertarget{w02:cue:2062}{}有一個高軌道，
% cue 2063 5393.530 5394.530
\hypertarget{w02:cue:2063}{}有一個低軌道，
% cue 2064 5394.530 5397.780
\hypertarget{w02:cue:2064}{}高軌道跟低軌道的差異，
% cue 2065 5397.780 5400.840
\hypertarget{w02:cue:2065}{}是在於他的民意的表現，
% cue 2066 5400.840 5403.870
\hypertarget{w02:cue:2066}{}他的民意的優越性不同，
% cue 2067 5403.870 5409.720
\hypertarget{w02:cue:2067}{}高軌道是由人民自己在所謂的憲法時刻\footnote{英文 constitutional moment，Bruce Ackerman 二元民主論（dualist democracy）的用語，指人民經動員並共同審議而行使高層次立法（higher lawmaking）的少數時期，與政治人物經代議機關的日常立法相對。Tremblay 刊本 534 頁引 Ackerman 1984 年的 Storrs Lectures（93 Yale Law Journal 1013）及 We the People: Foundations（Harvard University Press, 1991）。\url{https://law.yale.edu/bruce-ackerman}。}，
% cue 2068 5409.720 5411.720
\hypertarget{w02:cue:2068}{}憲法時刻是一個不好的譯語，
% cue 2069 5411.720 5413.720
\hypertarget{w02:cue:2069}{}因為憲法時刻讓我們馬上想到一個 point，
% cue 2070 5413.720 5414.720
\hypertarget{w02:cue:2070}{}一個點，
% cue 2071 5414.720 5415.900
\hypertarget{w02:cue:2071}{}三點十五分，
% cue 2072 5415.900 5417.090
\hypertarget{w02:cue:2072}{}不是，
% cue 2073 5417.090 5419.050
\hypertarget{w02:cue:2073}{}憲法時刻是一段很長的時期，
% cue 2074 5419.050 5422.360
\hypertarget{w02:cue:2074}{}所以我覺得他第一個可以批評的，
% cue 2075 5422.360 5423.270
\hypertarget{w02:cue:2075}{}用中文來說，
% cue 2076 5423.270 5425.320
\hypertarget{w02:cue:2076}{}就是你名不副實，
% cue 2077 5425.320 5427.350
\hypertarget{w02:cue:2077}{}你應該講憲法時期，
% cue 2078 5427.350 5429.630
\hypertarget{w02:cue:2078}{}或者憲法創造時期，
% cue 2079 5429.630 5431.630
\hypertarget{w02:cue:2079}{}憲法榮光時期，
% cue 2080 5431.630 5432.660
\hypertarget{w02:cue:2080}{}你不應該講憲法時刻，
% cue 2081 5432.660 5433.660
\hypertarget{w02:cue:2081}{}它不是一個點，
% cue 2082 5433.660 5437.100
\hypertarget{w02:cue:2082}{}那憲法的榮光時期是一個很長的時期，
% cue 2083 5437.100 5438.060
\hypertarget{w02:cue:2083}{}我們就不講了。

\tc{1:30:38}
% cue 2084 5438.060 5441.350
\hypertarget{w02:cue:2084}{}那就他說在這個很長的時期裡面，
% cue 2085 5441.350 5443.380
\hypertarget{w02:cue:2085}{}人民經過動員，
% cue 2086 5443.380 5445.600
\hypertarget{w02:cue:2086}{}高效的動員，
% cue 2087 5445.600 5448.760
\hypertarget{w02:cue:2087}{}而且經過非常仔細的思辨，
% cue 2088 5448.760 5451.490
\hypertarget{w02:cue:2088}{}他叫做 deliberation，
% cue 2089 5451.490 5452.490
\hypertarget{w02:cue:2089}{}mobilize，
% cue 2090 5452.490 5453.490
\hypertarget{w02:cue:2090}{}高度動員後的，
% cue 2091 5453.490 5454.490
\hypertarget{w02:cue:2091}{}大家都參與哦，
% cue 2092 5454.490 5455.490
\hypertarget{w02:cue:2092}{}不是只有一兩個，
% cue 2093 5455.490 5458.030
\hypertarget{w02:cue:2093}{}他們經過思辨，
% cue 2094 5458.030 5461.030
\hypertarget{w02:cue:2094}{}所以他們的民主的品質，
% cue 2095 5461.030 5462.060
\hypertarget{w02:cue:2095}{}要高於，
% cue 2096 5462.060 5463.380
\hypertarget{w02:cue:2096}{}就代表，
% cue 2097 5463.380 5465.310
\hypertarget{w02:cue:2097}{}在立法院裡面隨便混，
% cue 2098 5465.310 5467.400
\hypertarget{w02:cue:2098}{}然後各種不同的派系之間妥協，
% cue 2099 5467.400 5468.400
\hypertarget{w02:cue:2099}{}那有什麼思辨呢？

\tc{1:31:08}
% cue 2100 5468.400 5471.070
\hypertarget{w02:cue:2100}{}他意思是這樣子，
% cue 2101 5471.070 5472.420
\hypertarget{w02:cue:2101}{}他不是，
% cue 2102 5472.420 5474.770
\hypertarget{w02:cue:2102}{}不單純是因為人，
% cue 2103 5474.770 5475.790
\hypertarget{w02:cue:2103}{}不是因為我，
% cue 2104 5475.790 5476.770
\hypertarget{w02:cue:2104}{}跟我的代表，
% cue 2105 5476.770 5479.060
\hypertarget{w02:cue:2105}{}
% cue 2106 5479.060 5481.230
\hypertarget{w02:cue:2106}{}我是有高度思辨，
% cue 2107 5481.230 5483.290
\hypertarget{w02:cue:2107}{}我的代表是在哪裡混日子，
% cue 2108 5483.290 5485.320
\hypertarget{w02:cue:2108}{}他們在哪裡交易妥協的，
% cue 2109 5485.320 5487.400
\hypertarget{w02:cue:2109}{}所以我要優先於他，
% cue 2110 5487.400 5491.680
\hypertarget{w02:cue:2110}{}所以憲法要優先於他的解釋是這樣，
% cue 2111 5491.680 5495.840
\hypertarget{w02:cue:2111}{}這樣子有聽懂嗎？

\tc{1:31:35}
% cue 2112 5495.840 5498.790
\hypertarget{w02:cue:2112}{}所以，
% cue 2113 5498.790 5499.790
\hypertarget{w02:cue:2113}{}Bruce Ackerman，
% cue 2114 5499.790 5500.790
\hypertarget{w02:cue:2114}{}你要說他引起很多誤會，
% cue 2115 5500.790 5501.750
\hypertarget{w02:cue:2115}{}是活該，
% cue 2116 5501.750 5503.840
\hypertarget{w02:cue:2116}{}因為你自己講的話就是不清不楚，
% cue 2117 5503.840 5505.990
\hypertarget{w02:cue:2117}{}你是一個很像作家，
% cue 2118 5505.990 5506.990
\hypertarget{w02:cue:2118}{}高興寫，
% cue 2119 5506.990 5507.990
\hypertarget{w02:cue:2119}{}寫到哪裡算哪裡，
% cue 2120 5507.990 5509.920
\hypertarget{w02:cue:2120}{}你把憲法當小說寫，
% cue 2121 5509.920 5511.020
\hypertarget{w02:cue:2121}{}沒有錯，
% cue 2122 5511.020 5512.020
\hypertarget{w02:cue:2122}{}如果你講法制史，
% cue 2123 5512.020 5513.980
\hypertarget{w02:cue:2123}{}你正好變成一個很有趣的法制史，
% cue 2124 5513.980 5515.140
\hypertarget{w02:cue:2124}{}不然法制史，
% cue 2125 5515.140 5516.480
\hypertarget{w02:cue:2125}{}枯燥沒有人要讀，
% cue 2126 5516.480 5518.680
\hypertarget{w02:cue:2126}{}但是你這樣子建構你的法制史，
% cue 2127 5518.680 5519.680
\hypertarget{w02:cue:2127}{}你的問題就是，
% cue 2128 5519.680 5520.680
\hypertarget{w02:cue:2128}{}你的用語，
% cue 2129 5520.680 5521.770
\hypertarget{w02:cue:2129}{}你的概念，
% cue 2130 5521.770 5523.730
\hypertarget{w02:cue:2130}{}你的邏輯都不嚴謹，
% cue 2131 5523.730 5525.860
\hypertarget{w02:cue:2131}{}所以你會造成很多讀者，
% cue 2132 5525.860 5527.080
\hypertarget{w02:cue:2132}{}有不同的解讀，
% cue 2133 5527.080 5529.170
\hypertarget{w02:cue:2133}{}光光是 constitutional moment，
% cue 2134 5529.170 5530.170
\hypertarget{w02:cue:2134}{}批評的人一大堆，
% cue 2135 5530.170 5531.200
\hypertarget{w02:cue:2135}{}每個理解都不一樣，
% cue 2136 5531.200 5532.200
\hypertarget{w02:cue:2136}{}是你創造出來的，
% cue 2137 5532.200 5533.200
\hypertarget{w02:cue:2137}{}不能怪大家，
% cue 2138 5533.200 5534.450
\hypertarget{w02:cue:2138}{}而且他講的有道理，
% cue 2139 5534.450 5535.450
\hypertarget{w02:cue:2139}{}他也不隨便罵人。

\tc{1:32:15}
% cue 2140 5535.450 5536.740
\hypertarget{w02:cue:2140}{}好來，
% cue 2141 5536.740 5539.470
\hypertarget{w02:cue:2141}{}再看 530 試驗，
% cue 2142 5539.470 5543.820
\hypertarget{w02:cue:2142}{}The various versions of the Sophisticated View，
% cue 2143 5543.820 5545.300
\hypertarget{w02:cue:2143}{}剛才有主流說，
% cue 2144 5545.300 5546.870
\hypertarget{w02:cue:2144}{}現在呢，
% cue 2145 5546.870 5547.870
\hypertarget{w02:cue:2145}{}這是複雜觀點，
% cue 2146 5547.870 5550.560
\hypertarget{w02:cue:2146}{}剛才是 dominant，
% cue 2147 5550.560 5552.210
\hypertarget{w02:cue:2147}{}有沒有 dominant view，
% cue 2148 5552.210 5555.320
\hypertarget{w02:cue:2148}{}現在呢，
% cue 2149 5555.320 5558.240
\hypertarget{w02:cue:2149}{}是把 dominant view 加工了，
% cue 2150 5558.240 5560.130
\hypertarget{w02:cue:2150}{}所以我只能說，
% cue 2151 5560.130 5561.320
\hypertarget{w02:cue:2151}{}是複雜觀點，
% cue 2152 5561.320 5566.100
\hypertarget{w02:cue:2152}{}are complex indeed，
% cue 2153 5566.100 5567.350
\hypertarget{w02:cue:2153}{}本身很複雜，
% cue 2154 5567.350 5568.380
\hypertarget{w02:cue:2154}{}and highly controversial，
% cue 2155 5568.380 5569.280
\hypertarget{w02:cue:2155}{}and highly controversial，
% cue 2156 5569.280 5570.530
\hypertarget{w02:cue:2156}{}on their merit，
% cue 2157 5570.530 5571.560
\hypertarget{w02:cue:2157}{}就由他們的內容而言，
% cue 2158 5571.560 5572.590
\hypertarget{w02:cue:2158}{}也是高度爭議的。

\tc{1:32:52}
% cue 2159 5572.590 5575.000
\hypertarget{w02:cue:2159}{}The question I want to address is the following:。
% cue 2160 5575.000 5576.350
\hypertarget{w02:cue:2160}{}assuming that，
% cue 2161 5576.350 5577.350
\hypertarget{w02:cue:2161}{}假定，
% cue 2162 5577.350 5582.780
\hypertarget{w02:cue:2162}{}one version had succeeded in establishing the superior democratic credential，
% cue 2163 5582.780 5584.920
\hypertarget{w02:cue:2163}{}看這幾個字，
% cue 2164 5584.920 5586.520
\hypertarget{w02:cue:2164}{}它證明了什麼？
% cue 2165 5586.520 5590.960
\hypertarget{w02:cue:2165}{}證明了更優越的民主的，
% cue 2166 5590.960 5600.630
\hypertarget{w02:cue:2166}{}就是，
% cue 2167 5600.630 5601.650
\hypertarget{w02:cue:2167}{}overlegization，
% cue 2168 5601.650 5603.570
\hypertarget{w02:cue:2168}{}它不是位階的問題，
% cue 2169 5603.570 5604.840
\hypertarget{w02:cue:2169}{}它是一個 credential，
% cue 2170 5604.840 5608.980
\hypertarget{w02:cue:2170}{}我們說你這個人的資歷與履歷很豐富，
% cue 2171 5608.980 5610.070
\hypertarget{w02:cue:2171}{}credential，
% cue 2172 5610.070 5611.610
\hypertarget{w02:cue:2172}{}我這裡是這樣的，
% cue 2173 5611.610 5612.860
\hypertarget{w02:cue:2173}{}OK。

\tc{1:33:32}
% cue 2174 5612.860 5613.860
\hypertarget{w02:cue:2174}{}would it follow，
% cue 2175 5613.860 5614.860
\hypertarget{w02:cue:2175}{}我證明了這個，
% cue 2176 5614.860 5617.330
\hypertarget{w02:cue:2176}{}我的 democratic credential，
% cue 2177 5617.330 5618.330
\hypertarget{w02:cue:2177}{}is superior，
% cue 2178 5618.330 5619.330
\hypertarget{w02:cue:2178}{}比你好，
% cue 2179 5619.330 5620.900
\hypertarget{w02:cue:2179}{}能不能就說，
% cue 2180 5620.900 5621.990
\hypertarget{w02:cue:2180}{}would it follow that，
% cue 2181 5621.990 5627.520
\hypertarget{w02:cue:2181}{}specific legitimacy of judicial review has been established，
% cue 2182 5627.520 5630.470
\hypertarget{w02:cue:2182}{}你是證明你的民主品質比我優越，
% cue 2183 5630.470 5632.530
\hypertarget{w02:cue:2183}{}能不能就證明說，
% cue 2184 5632.530 5634.530
\hypertarget{w02:cue:2184}{}所以你就具有正當性，
% cue 2185 5634.530 5637.480
\hypertarget{w02:cue:2185}{}它在這裡它是打上問號，
% cue 2186 5637.480 5640.080
\hypertarget{w02:cue:2186}{}one might doubt so，
% cue 2187 5640.080 5641.490
\hypertarget{w02:cue:2187}{}我不可以這麼懷疑，
% cue 2188 5641.490 5650.940
\hypertarget{w02:cue:2188}{}there are at least two compelling objections against an affirmative answer to this question。

\tc{1:34:10}
% cue 2189 5650.940 5652.940
\hypertarget{w02:cue:2189}{}如果你認為我的論證，
% cue 2190 5652.940 5654.880
\hypertarget{w02:cue:2190}{}我的民主品質比較高，
% cue 2191 5654.880 5657.060
\hypertarget{w02:cue:2191}{}那我就當然應該具有正當性，
% cue 2192 5657.060 5658.060
\hypertarget{w02:cue:2192}{}我就可以做的話，
% cue 2193 5658.060 5660.050
\hypertarget{w02:cue:2193}{}那我至少有兩點懷疑，
% cue 2194 5660.050 5661.390
\hypertarget{w02:cue:2194}{}兩點挑戰，
% cue 2195 5661.390 5665.050
\hypertarget{w02:cue:2195}{}這是對 sophisticated view 的挑戰。

\tc{1:34:25}
% cue 2196 5665.050 5666.760
\hypertarget{w02:cue:2196}{}好，第一點，
% cue 2197 5666.760 5669.140
\hypertarget{w02:cue:2197}{}這個 534 頁最底下，
% cue 2198 5669.140 5670.140
\hypertarget{w02:cue:2198}{}first，
% cue 2199 5670.140 5674.590
\hypertarget{w02:cue:2199}{}the sophisticated view must necessarily presuppose，
% cue 2200 5674.590 5676.700
\hypertarget{w02:cue:2200}{}必然要預設，
% cue 2201 5676.700 5681.040
\hypertarget{w02:cue:2201}{}the meaning of the constitutional norm has been fixed，
% cue 2202 5681.040 5682.710
\hypertarget{w02:cue:2202}{}你們同意不同意？
% cue 2203 5682.710 5686.210
\hypertarget{w02:cue:2203}{}就是說，所謂憲法的意涵，
% cue 2204 5686.210 5690.110
\hypertarget{w02:cue:2204}{}是被確定定死的，
% cue 2205 5690.110 5692.510
\hypertarget{w02:cue:2205}{}and determined by the language of the constitution，
% cue 2206 5692.510 5696.700
\hypertarget{w02:cue:2206}{}就是被憲法的文句給定死在那裡，
% cue 2207 5696.700 5700.300
\hypertarget{w02:cue:2207}{}那麼法官呢，can discover，
% cue 2208 5700.300 5704.200
\hypertarget{w02:cue:2208}{}你只能發現，and apply it to relevant cases。

\tc{1:35:04}
% cue 2209 5704.200 5706.740
\hypertarget{w02:cue:2209}{}這個牽涉到憲法解釋方法，
% cue 2210 5706.740 5708.830
\hypertarget{w02:cue:2210}{}沒有前提問題，
% cue 2211 5708.830 5710.990
\hypertarget{w02:cue:2211}{}我們法官解釋憲法，
% cue 2212 5710.990 5713.280
\hypertarget{w02:cue:2212}{}你認為憲法的意思，
% cue 2213 5713.280 5715.660
\hypertarget{w02:cue:2213}{}因為憲法的文句定住了嗎？
% cue 2214 5715.660 5716.690
\hypertarget{w02:cue:2214}{}固定在那裡嗎？
% cue 2215 5716.690 5718.790
\hypertarget{w02:cue:2215}{}如果是這樣，那我們法官呢，
% cue 2216 5718.790 5719.790
\hypertarget{w02:cue:2216}{}沒得多說，
% cue 2217 5719.790 5720.850
\hypertarget{w02:cue:2217}{}我只能發現，
% cue 2218 5720.850 5721.940
\hypertarget{w02:cue:2218}{}然後運用了個案，
% cue 2219 5721.940 5724.450
\hypertarget{w02:cue:2219}{}可是如果他不確定，
% cue 2220 5724.450 5726.610
\hypertarget{w02:cue:2220}{}他沒有定義，
% cue 2221 5726.610 5727.700
\hypertarget{w02:cue:2221}{}他沒有確定是這個意思，
% cue 2222 5727.700 5730.370
\hypertarget{w02:cue:2222}{}而且有解釋的空間，
% cue 2223 5730.370 5731.560
\hypertarget{w02:cue:2223}{}那我法官做解釋的人，
% cue 2224 5731.560 5733.560
\hypertarget{w02:cue:2224}{}我又沒有上下其手，
% cue 2225 5733.560 5739.170
\hypertarget{w02:cue:2225}{}現在很少人採用機械說，
% cue 2226 5739.170 5741.590
\hypertarget{w02:cue:2226}{}籤則。

\tc{1:35:41}
% cue 2227 5741.590 5742.710
\hypertarget{w02:cue:2227}{}OK，
% cue 2228 5742.710 5744.960
\hypertarget{w02:cue:2228}{}如果是這樣，
% cue 2229 5744.960 5745.960
\hypertarget{w02:cue:2229}{}除非你是這樣子，
% cue 2230 5745.960 5746.960
\hypertarget{w02:cue:2230}{}不然你只不會有正當性，
% cue 2231 5746.960 5749.330
\hypertarget{w02:cue:2231}{}那這個叫做。
% cue 2232 5749.330 5750.650
\hypertarget{w02:cue:2232}{}the formalist condition，
% cue 2233 5750.650 5751.650
\hypertarget{w02:cue:2233}{}就是形式論的條件。

\tc{1:35:51}
% cue 2234 5751.650 5753.030
\hypertarget{w02:cue:2234}{}on various occasions，
% cue 2235 5753.030 5754.030
\hypertarget{w02:cue:2235}{}scholars, judges，
% cue 2236 5754.030 5755.160
\hypertarget{w02:cue:2236}{}constitutional lawyers have believed that the formalist conditions could be satisfied，
% cue 2237 5755.160 5756.160
\hypertarget{w02:cue:2237}{}有很多情況是可以的。
% cue 2238 5756.160 5758.470
\hypertarget{w02:cue:2238}{}這個我們就不管它什麼叫做 strict，
% cue 2239 5758.470 5759.460
\hypertarget{w02:cue:2239}{}constructionism，
% cue 2240 5759.460 5760.910
\hypertarget{w02:cue:2240}{}這也是，
% cue 2241 5760.910 5761.930
\hypertarget{w02:cue:2241}{}這也是，
% cue 2242 5761.930 5763.540
\hypertarget{w02:cue:2242}{}你會不會有正當性，
% cue 2243 5763.540 5766.030
\hypertarget{w02:cue:2243}{}這就是形式論的條件。
% cue 2244 5766.030 5768.560 刪：與 cue 2234 同文（連續 4 個重複解碼）
% cue 2245 5768.560 5769.560 刪：與 cue 2235 同文（連續 4 個重複解碼）
% cue 2246 5769.560 5774.710 刪：與 cue 2236 同文（連續 4 個重複解碼）
% cue 2247 5774.710 5777.800 刪：與 cue 2237 同文（連續 4 個重複解碼）

\tc{1:36:17}
% cue 2248 5777.800 5784.170
\hypertarget{w02:cue:2248}{}Chief Justice Marshall's theory of legal meaning and interpretation came within formalism，
% cue 2249 5784.170 5788.450
\hypertarget{w02:cue:2249}{}although not with strict constructionism\footnote{英文，嚴格解釋主義，主張法律文字依其字面嚴格適用，與擴張解釋相對。Tremblay 刊本 535 頁謂 Marshall 的法律意義與解釋理論屬於形式主義（formalism），但不屬於嚴格解釋主義。\url{https://www.law.cornell.edu/wex/strict_construction}。}，
% cue 2250 5788.450 5791.440
\hypertarget{w02:cue:2250}{}這個我們就不管它什麼叫做 strict constructionism，
% cue 2251 5791.440 5794.840
\hypertarget{w02:cue:2251}{}這也是純粹的解釋主義了，
% cue 2252 5794.840 5796.050
\hypertarget{w02:cue:2252}{}先不管這個。

\tc{1:36:36}
% cue 2253 5796.050 5799.050
\hypertarget{w02:cue:2253}{}那我們再往下跳一句，
% cue 2254 5799.050 5801.050
\hypertarget{w02:cue:2254}{}it is generally admitted，
% cue 2255 5801.050 5802.200
\hypertarget{w02:cue:2255}{}大家都同意，
% cue 2256 5802.200 5804.320
\hypertarget{w02:cue:2256}{}constitutional language is vague，
% cue 2257 5804.320 5806.630
\hypertarget{w02:cue:2257}{}憲法的語句是模糊的，
% cue 2258 5806.630 5808.890
\hypertarget{w02:cue:2258}{}abstract，抽象的，
% cue 2259 5808.890 5810.110
\hypertarget{w02:cue:2259}{}and general，
% cue 2260 5810.110 5811.200
\hypertarget{w02:cue:2260}{}一般的，
% cue 2261 5811.200 5813.200
\hypertarget{w02:cue:2261}{}and makes at least certain constitutional provisions，
% cue 2262 5813.200 5814.200
\hypertarget{w02:cue:2262}{}indeterminate，
% cue 2263 5814.200 5817.160
\hypertarget{w02:cue:2263}{}因為它用了這種模糊的語句，
% cue 2264 5817.160 5819.320
\hypertarget{w02:cue:2264}{}至少某些憲法的規定，
% cue 2265 5819.320 5823.010
\hypertarget{w02:cue:2265}{}是不確定它的意思，
% cue 2266 5823.010 5824.520
\hypertarget{w02:cue:2266}{}indeterminate，
% cue 2267 5824.520 5825.740
\hypertarget{w02:cue:2267}{}沒有確定的，
% cue 2268 5825.740 5828.240
\hypertarget{w02:cue:2268}{}因此，
% cue 2269 5828.240 5830.400
\hypertarget{w02:cue:2269}{}the language must be interpreted or constructed，
% cue 2270 5830.400 5831.520
\hypertarget{w02:cue:2270}{}所以，
% cue 2271 5831.520 5833.600
\hypertarget{w02:cue:2271}{}憲法才需要解釋，
% cue 2272 5833.600 5835.860
\hypertarget{w02:cue:2272}{}on the grounds of value，
% cue 2273 5835.860 5839.870
\hypertarget{w02:cue:2273}{}依照價值作為基礎來加以解釋，
% cue 2274 5839.870 5844.960
\hypertarget{w02:cue:2274}{}而不是只是單純依照那種形式的憲法文句，
% cue 2275 5844.960 5846.470
\hypertarget{w02:cue:2275}{}就可以決定。

\tc{1:37:26}
% cue 2276 5846.470 5849.350
\hypertarget{w02:cue:2276}{}since constitutional interpretation，
% cue 2277 5849.350 5852.530
\hypertarget{w02:cue:2277}{}raises deep controversial issue of the social policy and political morality，
% cue 2278 5852.530 5856.710
\hypertarget{w02:cue:2278}{}憲法解釋會牽涉到很有爭議的社會政策，
% cue 2279 5856.710 5857.930
\hypertarget{w02:cue:2279}{}嗯？
% cue 2280 5857.930 5858.930
\hypertarget{w02:cue:2280}{}比如說移民，
% cue 2281 5858.930 5859.900
\hypertarget{w02:cue:2281}{}要不要趕出去嗎？
% cue 2282 5859.900 5861.850
\hypertarget{w02:cue:2282}{}宗教基礎可不可以趕走移民啊？
% cue 2283 5861.850 5863.820
\hypertarget{w02:cue:2283}{}還有 political morality，
% cue 2284 5863.820 5865.750
\hypertarget{w02:cue:2284}{}政治上這樣道德嗎？
% cue 2285 5865.750 5867.850
\hypertarget{w02:cue:2285}{}計劃的計劃的計劃，
% cue 2286 5867.850 5870.000
\hypertarget{w02:cue:2286}{}計劃的計劃。
% cue 2287 5870.000 5872.140
\hypertarget{w02:cue:2287}{}計劃的計劃。
% cue 2288 5872.140 5873.310
\hypertarget{w02:cue:2288}{}計劃的計劃。
% cue 2289 5873.310 5874.310
\hypertarget{w02:cue:2289}{}計劃的計劃。
% cue 2290 5874.310 5876.590
\hypertarget{w02:cue:2290}{}可以終於提供的法律的觀點，
% cue 2291 5876.590 5878.660
\hypertarget{w02:cue:2291}{}所有觀點，
% cue 2292 5878.660 5879.800
\hypertarget{w02:cue:2292}{}尤其是理想和志，
% cue 2293 5879.800 5881.790
\hypertarget{w02:cue:2293}{}不可避免的，
% cue 2294 5881.790 5884.910
\hypertarget{w02:cue:2294}{}它會把他自己的看法，
% cue 2295 5884.910 5885.820
\hypertarget{w02:cue:2295}{}法官的看法，
% cue 2296 5885.820 5888.060
\hypertarget{w02:cue:2296}{}會滲入到解釋裡去。

\tc{1:38:08}
% cue 2297 5888.060 5892.110
\hypertarget{w02:cue:2297}{}since judges must determine for themselves。
% cue 2298 5892.110 5894.070
\hypertarget{w02:cue:2298}{}the kind of consideration。
% cue 2299 5894.070 5896.190
\hypertarget{w02:cue:2299}{}that are sufficiently strong。
% cue 2300 5896.190 5899.150
\hypertarget{w02:cue:2300}{}to justify their decisions not enforce，
% cue 2301 5899.150 5902.530
\hypertarget{w02:cue:2301}{}這是 judicial super legislature。

\tc{1:38:22}
% cue 2302 5902.530 5904.880
\hypertarget{w02:cue:2302}{}我們前面看到的是，
% cue 2303 5904.880 5906.910
\hypertarget{w02:cue:2303}{}你認為是憲法，
% cue 2304 5906.910 5908.930
\hypertarget{w02:cue:2304}{}它是一個超級的立法，
% cue 2305 5908.930 5911.530
\hypertarget{w02:cue:2305}{}但是經過了大法官解釋，
% cue 2306 5911.530 5913.530
\hypertarget{w02:cue:2306}{}體現出來的議詞，
% cue 2307 5913.530 5916.690
\hypertarget{w02:cue:2307}{}你會發現司法是一個超級立法者，
% cue 2308 5916.690 5920.030
\hypertarget{w02:cue:2308}{}這個你只要看過去這幾年，
% cue 2309 5920.030 5923.440
\hypertarget{w02:cue:2309}{}許宗力\footnote{許宗力於 2003 年 10 月 1 日至 2011 年 9 月 30 日任中華民國司法院大法官，2016 年 11 月 1 日至 2024 年 10 月 31 日再任大法官並任司法院院長；釋字第 791 號解釋由其擔任主席，113 年憲判字第 8 號判決由其擔任審判長。}法庭用司法積極主義，
% cue 2310 5923.440 5926.870
\hypertarget{w02:cue:2310}{}把過去的大法官解釋一個個翻掉了，
% cue 2311 5926.870 5929.110
\hypertarget{w02:cue:2311}{}那你覺得它是不是 super legislature，
% cue 2312 5929.110 5931.060
\hypertarget{w02:cue:2312}{}我講有錯嗎？

\tc{1:38:51}
% cue 2313 5931.060 5933.020
\hypertarget{w02:cue:2313}{}寫在法律裡呢，
% cue 2314 5933.020 5934.190
\hypertarget{w02:cue:2314}{}你說通姦\footnote{中華民國司法院大法官釋字第 791 號解釋（2020 年 5 月 29 日）解釋文：「刑法第239條規定：『有配偶而與人通姦者，處1年以下有期徒刑。其相姦者亦同。』對憲法第22條所保障性自主權之限制，與憲法第23條比例原則不符，應自本解釋公布之日起失其效力；於此範圍內，本院釋字第554號解釋應予變更。」\url{https://cons.judicial.gov.tw/docdata.aspx?fid=100&id=310972}。}，
% cue 2315 5934.190 5935.380
\hypertarget{w02:cue:2315}{}現在要除罪，
% cue 2316 5935.380 5937.380
\hypertarget{w02:cue:2316}{}通姦罪立法者不修改，
% cue 2317 5937.380 5939.440
\hypertarget{w02:cue:2317}{}他沒有修改的權利嗎？
% cue 2318 5939.440 5940.750
\hypertarget{w02:cue:2318}{}可以啊，
% cue 2319 5940.750 5941.750
\hypertarget{w02:cue:2319}{}他有沒有修改，
% cue 2320 5941.750 5942.770
\hypertarget{w02:cue:2320}{}他沒有修改，
% cue 2321 5942.770 5943.770
\hypertarget{w02:cue:2321}{}結果你這樣告訴他，
% cue 2322 5943.770 5945.770
\hypertarget{w02:cue:2322}{}通姦不可以成為犯罪，
% cue 2323 5945.770 5948.080
\hypertarget{w02:cue:2323}{}這意思 super legislature，
% cue 2324 5948.080 5950.650
\hypertarget{w02:cue:2324}{}所以是司法成為超級立法者，
% cue 2325 5950.650 5952.650
\hypertarget{w02:cue:2325}{}所以是司法成為超級立法者。

\tc{1:39:12}
% cue 2326 5952.650 5955.990
\hypertarget{w02:cue:2326}{}這些反對特別有力，
% cue 2327 5955.990 5959.240
\hypertarget{w02:cue:2327}{}我們這樣的反對特別有力，
% cue 2328 5959.240 5961.180
\hypertarget{w02:cue:2328}{}我們能不能在憲法的成文憲法之外，
% cue 2329 5961.180 5962.180
\hypertarget{w02:cue:2329}{}去找尋價值，
% cue 2330 5962.180 5964.180
\hypertarget{w02:cue:2330}{}然後來證立來解釋這個成文憲法。
% cue 2331 5964.180 5966.170
\hypertarget{w02:cue:2331}{}。
% cue 2332 5966.170 5967.170
\hypertarget{w02:cue:2332}{}我們能不能在憲法的成文憲法之外，
% cue 2333 5967.170 5968.170
\hypertarget{w02:cue:2333}{}去找尋價值，
% cue 2334 5968.170 5970.140
\hypertarget{w02:cue:2334}{}然後來證立來解釋這個成文憲法。
% cue 2335 5970.140 5972.140
\hypertarget{w02:cue:2335}{}。

\tc{1:39:32}
% cue 2336 5972.140 5973.120
\hypertarget{w02:cue:2336}{}Canadian Supreme Court，
% cue 2337 5973.120 5974.680
\hypertarget{w02:cue:2337}{}for example，
% cue 2338 5974.680 5976.680
\hypertarget{w02:cue:2338}{}has admitted that certain constitutional clauses，
% cue 2339 5976.680 5978.570
\hypertarget{w02:cue:2339}{}they can open，
% cue 2340 5978.570 5980.760
\hypertarget{w02:cue:2340}{}and that single incontrovertible meaning is not apparent，
% cue 2341 5980.760 5982.760
\hypertarget{w02:cue:2341}{}that their historical language used，
% cue 2342 5982.760 5985.460
\hypertarget{w02:cue:2342}{}that their historical language used，
% cue 2343 5985.460 5986.460
\hypertarget{w02:cue:2343}{}等等，
% cue 2344 5986.460 5987.460
\hypertarget{w02:cue:2344}{}也就可以了。
% cue 2345 5987.460 5988.550
\hypertarget{w02:cue:2345}{}好，
% cue 2346 5988.550 5989.550
\hypertarget{w02:cue:2346}{}所以這一點變成是有意義的。
% cue 2347 5989.550 5990.550
\hypertarget{w02:cue:2347}{}好，
% cue 2348 5990.550 5991.710
\hypertarget{w02:cue:2348}{}以上剛才講的是第一個反對。

\tc{1:39:51}
% cue 2349 5991.710 5993.020
\hypertarget{w02:cue:2349}{}第二個反對 537 頁，
% cue 2350 5993.020 5995.110
\hypertarget{w02:cue:2350}{}the second objection is more fundamental，
% cue 2351 6001.230 6002.310
\hypertarget{w02:cue:2351}{}也就可以了。
% cue 2352 6002.310 6005.870
\hypertarget{w02:cue:2352}{}好，所以這一點變成是有意義的。
% cue 2353 6005.870 6008.920
\hypertarget{w02:cue:2353}{}好，以上剛才講的是第一個反對。
% cue 2354 6008.920 6011.780
\hypertarget{w02:cue:2354}{}第二個反對 537 頁。
% cue 2355 6011.780 6016.740
\hypertarget{w02:cue:2355}{}第二個反對是更根本的。

\tc{1:40:16}
% cue 2356 6016.740 6019.150
\hypertarget{w02:cue:2356}{}剛才我們只是看著形式的。
% cue 2357 6019.150 6021.390
\hypertarget{w02:cue:2357}{}現在是更根本的反對。
% cue 2358 6021.390 6030.800
\hypertarget{w02:cue:2358}{}即便你說憲法不只是形式上位階高。
% cue 2359 6030.800 6036.830
\hypertarget{w02:cue:2359}{}而且民主的內涵高於普通的立法。
% cue 2360 6036.830 6041.970
\hypertarget{w02:cue:2360}{}因此，當兩者發生衝突的時候。
% cue 2361 6041.970 6043.630
\hypertarget{w02:cue:2361}{}衝突的時候以憲法為準。
% cue 2362 6043.630 6048.260
\hypertarget{w02:cue:2362}{}你就算你可以這麼講。
% cue 2363 6048.260 6051.190
\hypertarget{w02:cue:2363}{}也不當然導出。
% cue 2364 6051.190 6060.780
\hypertarget{w02:cue:2364}{}即便是這樣子。
% cue 2365 6060.780 6062.990
\hypertarget{w02:cue:2365}{}你也不能當然導出。
% cue 2366 6062.990 6066.840
\hypertarget{w02:cue:2366}{}法院有權利可以拒絕執行。
% cue 2367 6066.840 6070.770
\hypertarget{w02:cue:2367}{}民主程序所制定的立法。

\tc{1:41:10}
% cue 2368 6070.770 6077.910
\hypertarget{w02:cue:2368}{}代議民主。
% cue 2369 6077.910 6079.230
\hypertarget{w02:cue:2369}{}代議的國會。
% cue 2370 6079.230 6086.000
\hypertarget{w02:cue:2370}{}它可能是一個更適當的場所。
% cue 2371 6086.000 6087.300
\hypertarget{w02:cue:2371}{}幹嘛。
% cue 2372 6087.300 6089.300
\hypertarget{w02:cue:2372}{}to make such decisions。
% cue 2373 6089.300 6093.230
\hypertarget{w02:cue:2373}{}for the very reason that they best represent the will。
% cue 2374 6093.230 6095.730
\hypertarget{w02:cue:2374}{}你假如說要代表民意。
% cue 2375 6095.730 6097.000
\hypertarget{w02:cue:2375}{}以這個做基礎。
% cue 2376 6097.000 6100.610
\hypertarget{w02:cue:2376}{}那可能民意代表還是比法官更能代表民意。

\tc{1:41:40}
% cue 2377 6100.610 6105.540
\hypertarget{w02:cue:2377}{}這個不是只是民主討論品質的問題而已。
% cue 2378 6105.540 6107.830
\hypertarget{w02:cue:2378}{}是機構適當性的問題。
% cue 2379 6107.830 6110.850
\hypertarget{w02:cue:2379}{}用我們大陸法系的話來講的。
% cue 2380 6110.850 6112.190
\hypertarget{w02:cue:2380}{}這個應該叫做什麼？
% cue 2381 6112.190 6114.170
\hypertarget{w02:cue:2381}{}功能最適\footnote{中華民國憲法法庭 111 年憲判字第 9 號判決（2022 年 6 月 24 日）理由第 20 段闡述權力分立原則，謂應「將上述不同性質之權力，分別歸屬於組織及程序等功能最適之不同機關行使，以發揮功效」。\url{https://cons.judicial.gov.tw/docdata.aspx?fid=38&id=339757}。}吧。
% cue 2382 6114.170 6117.100
\hypertarget{w02:cue:2382}{}我們無論如何。
% cue 2383 6117.100 6120.480
\hypertarget{w02:cue:2383}{}還是從功能最適的觀點來講。
% cue 2384 6120.480 6122.350
\hypertarget{w02:cue:2384}{}有比你更適合的。
% cue 2385 6122.350 6124.300
\hypertarget{w02:cue:2385}{}也不是當然說。
% cue 2386 6124.300 6125.940
\hypertarget{w02:cue:2386}{}哇你就可以當然不執行。
% cue 2387 6125.940 6126.560
\hypertarget{w02:cue:2387}{}不是你。
% cue 2388 6126.560 6127.280
\hypertarget{w02:cue:2388}{}別人。

\tc{1:42:07}
% cue 2389 6127.280 6130.170
\hypertarget{w02:cue:2389}{}所以你想想另外一個有趣的問題。
% cue 2390 6130.170 6133.740
\hypertarget{w02:cue:2390}{}英國現在還沒有成文憲法。
% cue 2391 6133.740 6136.820
\hypertarget{w02:cue:2391}{}但是英國自從 1998 年以後。
% cue 2392 6136.820 6141.740
\hypertarget{w02:cue:2392}{}人權法\footnote{英國 Human Rights Act 1998 第 4 條規定，法院認定國會立法之規定與歐洲人權公約所保障之權利不相容時，得作成不相容宣告（declaration of incompatibility）；同條第 6 項規定該宣告不影響該規定之效力、繼續施行與執行，亦不拘束訴訟當事人。聯合王國最高法院自 2009 年 10 月 1 日起列為同條所稱之法院。\url{https://www.legislation.gov.uk/ukpga/1998/42/section/4}。}裡面有規定。
% cue 2393 6141.740 6145.980
\hypertarget{w02:cue:2393}{}這個法律不能侵害。
% cue 2394 6145.980 6158.430
\hypertarget{w02:cue:2394}{}最高法院現在不是宣告 1998 人權法。
% cue 2395 6158.430 6165.510
\hypertarget{w02:cue:2395}{}但它可以宣告法律確實抵觸 1998 人權法。
% cue 2396 6165.510 6170.430
\hypertarget{w02:cue:2396}{}然後退回不回看你們要怎麼處理。
% cue 2397 6170.430 6174.710
\hypertarget{w02:cue:2397}{}這個就比較簡單。
% cue 2398 6174.710 6175.900
\hypertarget{w02:cue:2398}{}它的民主品質有問題。
% cue 2399 6175.900 6178.140
\hypertarget{w02:cue:2399}{}主論是不對。
% cue 2400 6178.140 6180.810
\hypertarget{w02:cue:2400}{}但是不是民主法來決定的。
% cue 2401 6180.810 6185.530
\hypertarget{w02:cue:2401}{}比較國會比較接近現在我講的這個法的意思。

\tc{1:43:05}
% cue 2402 6185.530 6193.960
\hypertarget{w02:cue:2402}{}Judicial review takes certain controversial political and moral issues out of the public arena。
% cue 2403 6193.960 6196.920
\hypertarget{w02:cue:2403}{}where, in a decent democratic society。
% cue 2404 6196.920 6200.470
\hypertarget{w02:cue:2404}{}They should be debated and decided in principle。
% cue 2405 6207.250 6210.880
\hypertarget{w02:cue:2405}{}第一審就是一個有教養的社會，
% cue 2406 6210.880 6214.930
\hypertarget{w02:cue:2406}{}他們應該由人民來辯論，
% cue 2407 6214.930 6216.410
\hypertarget{w02:cue:2407}{}而不是由法官來做決定，
% cue 2408 6216.410 6218.810
\hypertarget{w02:cue:2408}{}要不要廢死\footnote{中華民國憲法法庭 113 年憲判字第 8 號判決（2024 年 9 月 20 日）主文第一項宣告，刑法第 271 條第 1 項等規定以死刑為最重本刑部分，「僅得適用於個案犯罪情節屬最嚴重，且其刑事程序符合憲法最嚴密之正當法律程序要求之情形。於此範圍內，與憲法保障人民生命權之意旨尚屬無違。」\cite{ccj113-8}。}，
% cue 2409 6218.810 6221.470
\hypertarget{w02:cue:2409}{}為什麼不是人民決定，
% cue 2410 6221.470 6222.470
\hypertarget{w02:cue:2410}{}為什麼是法官決定，
% cue 2411 6222.470 6224.390
\hypertarget{w02:cue:2411}{}憲法沒寫啊，
% cue 2412 6224.390 6226.810
\hypertarget{w02:cue:2412}{}聽懂嗎？
% cue 2413 6226.810 6228.950
\hypertarget{w02:cue:2413}{}它是在憲法以外的價值，
% cue 2414 6228.950 6231.100
\hypertarget{w02:cue:2414}{}或者是憲法裡面找到的價值，
% cue 2415 6231.100 6233.300
\hypertarget{w02:cue:2415}{}但是憲法沒有任何規定。

\tc{1:43:53}
% cue 2416 6233.300 6234.440
\hypertarget{w02:cue:2416}{}好了，以上到這裡，
% cue 2417 6234.440 6237.440
\hypertarget{w02:cue:2417}{}我們已經把兩個東西都講完了，
% cue 2418 6237.440 6240.440
\hypertarget{w02:cue:2418}{}然後到 538 頁看一下最後的結論，
% cue 2419 6240.440 6256.920
\hypertarget{w02:cue:2419}{}因為現在的這些理論，
% cue 2420 6256.920 6260.450
\hypertarget{w02:cue:2420}{}沒有辦法成功地證立它的正當性，
% cue 2421 6260.450 6265.230
\hypertarget{w02:cue:2421}{}所以我認為不要這麼複雜，
% cue 2422 6265.230 6267.310
\hypertarget{w02:cue:2422}{}我要來提出我的一般性，
% cue 2423 6267.310 6271.870
\hypertarget{w02:cue:2423}{}我們不需要再去尋求特定的授權，
% cue 2424 6271.870 6273.090
\hypertarget{w02:cue:2424}{}或者是特定的具體，
% cue 2425 6273.090 6275.460
\hypertarget{w02:cue:2425}{}特定的正當化理論，
% cue 2426 6275.460 6278.460
\hypertarget{w02:cue:2426}{}我現在來講一講一般的正當化理論。

\section{一般正當性的六個命題}

\tc{1:44:38}
% cue 2427 6278.460 6280.780
\hypertarget{w02:cue:2427}{}好，現在看 538 頁 2，
% cue 2428 6280.780 6283.090
\hypertarget{w02:cue:2428}{}In this section，
% cue 2429 6283.090 6284.540
\hypertarget{w02:cue:2429}{}I argue，
% cue 2430 6284.540 6287.430
\hypertarget{w02:cue:2430}{}judicial review is morally legitimate，
% cue 2431 6287.430 6288.430
\hypertarget{w02:cue:2431}{}法律，
% cue 2432 6288.430 6289.490
\hypertarget{w02:cue:2432}{}道德上正當的，
% cue 2433 6289.490 6291.760
\hypertarget{w02:cue:2433}{}原則上是道德正當的，
% cue 2434 6291.760 6297.180
\hypertarget{w02:cue:2434}{}而且也是道德上有正當，
% cue 2435 6297.180 6299.180
\hypertarget{w02:cue:2435}{}A matter of political morality，
% cue 2436 6299.180 6300.460
\hypertarget{w02:cue:2436}{}for this purpose，
% cue 2437 6300.460 6304.430
\hypertarget{w02:cue:2437}{}it is necessary to recall the concept of legitimacy，
% cue 2438 6304.430 6306.680
\hypertarget{w02:cue:2438}{}can be understood in different way，
% cue 2439 6306.680 6309.030
\hypertarget{w02:cue:2439}{}又回到正當性，
% cue 2440 6309.030 6312.380
\hypertarget{w02:cue:2440}{}正當性這個概念可以有不同的理解，
% cue 2441 6312.380 6314.600
\hypertarget{w02:cue:2441}{}下面這裡你們要自己看，
% cue 2442 6314.600 6315.680
\hypertarget{w02:cue:2442}{}沒有時間了。

\tc{1:45:15}
% cue 2443 6315.680 6318.330
\hypertarget{w02:cue:2443}{}我們現在要來講的是，
% cue 2444 6318.330 6320.000
\hypertarget{w02:cue:2444}{}下面，
% cue 2445 6320.000 6320.100
\hypertarget{w02:cue:2445}{}他說，
% cue 2446 6320.000 6325.280
\hypertarget{w02:cue:2446}{}這個案子跟，
% cue 2447 6325.280 6326.470
\hypertarget{w02:cue:2447}{}別的，
% cue 2448 6326.470 6327.470
\hypertarget{w02:cue:2448}{}跟前面那個案子不同的，
% cue 2449 6327.470 6328.470
\hypertarget{w02:cue:2449}{}不一樣的是，
% cue 2450 6328.470 6332.970
\hypertarget{w02:cue:2450}{}在 538 頁的最底下，
% cue 2451 6332.970 6335.350
\hypertarget{w02:cue:2451}{}this form of legitimacy，
% cue 2452 6335.350 6337.670
\hypertarget{w02:cue:2452}{}這一種形式的正當性，
% cue 2453 6337.670 6339.790
\hypertarget{w02:cue:2453}{}logically pre-supposes，
% cue 2454 6339.790 6341.850
\hypertarget{w02:cue:2454}{}在邏輯上預先假設，
% cue 2455 6341.850 6348.290
\hypertarget{w02:cue:2455}{}the existence of at least one general moral principle，
% cue 2456 6348.290 6352.560
\hypertarget{w02:cue:2456}{}一個普遍適用的道德原則，
% cue 2457 6352.560 6356.190
\hypertarget{w02:cue:2457}{}conferring some general authorization to act，
% cue 2458 6356.190 6358.500
\hypertarget{w02:cue:2458}{}它可以行動，
% cue 2459 6358.500 6360.060
\hypertarget{w02:cue:2459}{}除非，
% cue 2460 6360.060 6363.050
\hypertarget{w02:cue:2460}{}所以它的預設是倒過來的，
% cue 2461 6363.050 6365.170
\hypertarget{w02:cue:2461}{}除非證明它這樣這樣這樣，
% cue 2462 6365.170 6367.390
\hypertarget{w02:cue:2462}{}否則它可以行動，
% cue 2463 6367.390 6369.390
\hypertarget{w02:cue:2463}{}可以認為它是有正當性，
% cue 2464 6369.390 6373.280
\hypertarget{w02:cue:2464}{}unless some specific moral principles，
% cue 2465 6373.280 6375.280
\hypertarget{w02:cue:2465}{}makes them impermissible，
% cue 2466 6375.280 6379.590
\hypertarget{w02:cue:2466}{}除非有一個特定的道德原則，
% cue 2467 6379.590 6381.620
\hypertarget{w02:cue:2467}{}使得它變成不正當，
% cue 2468 6381.620 6385.220
\hypertarget{w02:cue:2468}{}不然你要假設它是正當。

\tc{1:46:25}
% cue 2469 6385.220 6388.940
\hypertarget{w02:cue:2469}{}general principle of this sort ground，
% cue 2470 6388.940 6391.700
\hypertarget{w02:cue:2470}{}when I call it a general legitimacy，
% cue 2471 6391.700 6394.140
\hypertarget{w02:cue:2471}{}所以這個 general legitimacy，
% cue 2472 6394.140 6396.140
\hypertarget{w02:cue:2472}{}跟 specific legitimacy，
% cue 2473 6396.140 6397.050
\hypertarget{w02:cue:2473}{}剛才一開始，
% cue 2474 6397.050 6398.170
\hypertarget{w02:cue:2474}{}有講過話，
% cue 2475 6398.170 6400.200
\hypertarget{w02:cue:2475}{}它已經再重複，
% cue 2476 6400.200 6402.260
\hypertarget{w02:cue:2476}{}現在再簡單講最後一次，
% cue 2477 6402.260 6405.600
\hypertarget{w02:cue:2477}{}specific legitimacy，
% cue 2478 6405.600 6408.100
\hypertarget{w02:cue:2478}{}是你要獲得授權才會正當，
% cue 2479 6408.100 6410.960
\hypertarget{w02:cue:2479}{}除非你獲得授權，
% cue 2480 6410.960 6412.440
\hypertarget{w02:cue:2480}{}正當，general legitimacy，
% cue 2481 6412.440 6413.410
\hypertarget{w02:cue:2481}{}相反，
% cue 2482 6413.410 6415.620
\hypertarget{w02:cue:2482}{}原則上你就是正當的，
% cue 2483 6415.620 6417.750
\hypertarget{w02:cue:2483}{}除非你可以證明，
% cue 2484 6417.750 6421.150
\hypertarget{w02:cue:2484}{}有一個特別的道德原則，
% cue 2485 6421.150 6424.490
\hypertarget{w02:cue:2485}{}來證明它是不正當的，
% cue 2486 6424.490 6425.490
\hypertarget{w02:cue:2486}{}沒有證明為不正當以前，
% cue 2487 6425.490 6428.200
\hypertarget{w02:cue:2487}{}就推定它是正當，
% cue 2488 6428.200 6429.200
\hypertarget{w02:cue:2488}{}它的舉證責任少的非常多。

\tc{1:47:09}
% cue 2489 6429.200 6430.990
\hypertarget{w02:cue:2489}{}好，
% cue 2490 6430.990 6431.990
\hypertarget{w02:cue:2490}{}那我們現在來看呢，
% cue 2491 6431.990 6433.040
\hypertarget{w02:cue:2491}{}它說，
% cue 2492 6433.040 6435.110
\hypertarget{w02:cue:2492}{}the general legitimacy thesis，
% cue 2493 6435.110 6436.010
\hypertarget{w02:cue:2493}{}剛才有，
% cue 2494 6436.010 6437.920
\hypertarget{w02:cue:2494}{}specific legitimacy thesis，
% cue 2495 6438.100 6440.130
\hypertarget{w02:cue:2495}{}現在是 general legitimacy thesis，
% cue 2496 6440.130 6444.120
\hypertarget{w02:cue:2496}{}也可以被 summarize or stated in six propositions，
% cue 2497 6444.120 6445.150
\hypertarget{w02:cue:2497}{}有六個命題，
% cue 2498 6445.150 6446.970
\hypertarget{w02:cue:2498}{}明白，
% cue 2499 6446.970 6448.650
\hypertarget{w02:cue:2499}{}GLT1，
% cue 2500 6448.650 6451.980
\hypertarget{w02:cue:2500}{}all political authorities in the state，
% cue 2501 6451.980 6456.270
\hypertarget{w02:cue:2501}{}are bound by a general principle of political morality，
% cue 2502 6456.270 6457.930
\hypertarget{w02:cue:2502}{}所以，
% cue 2503 6457.930 6458.930
\hypertarget{w02:cue:2503}{}我們現在說，
% cue 2504 6458.930 6459.930
\hypertarget{w02:cue:2504}{}一個，
% cue 2505 6459.930 6460.930
\hypertarget{w02:cue:2505}{}政治道德一般原則，
% cue 2506 6460.930 6462.280
\hypertarget{w02:cue:2506}{}是什麼呢？
% cue 2507 6462.280 6463.190
\hypertarget{w02:cue:2507}{}應該儘量的，
% cue 2508 6463.190 6465.400
\hypertarget{w02:cue:2508}{}以合法的方式來行使權力，
% cue 2509 6465.400 6468.220
\hypertarget{w02:cue:2509}{}這裡它剛開始所有角度，
% cue 2510 6468.220 6471.240
\hypertarget{w02:cue:2510}{}這裡它正式變成第一命題，
% cue 2511 6471.240 6472.240
\hypertarget{w02:cue:2511}{}這跟剛才是不一樣的哦。

\tc{1:47:52}
% cue 2512 6472.240 6474.220
\hypertarget{w02:cue:2512}{}那這裡它就引述了一些西方，
% cue 2513 6474.220 6479.210
\hypertarget{w02:cue:2513}{}尤其是自由主義的一些界限，
% cue 2514 6479.210 6480.370
\hypertarget{w02:cue:2514}{}比如說，
% cue 2515 6480.370 6481.590
\hypertarget{w02:cue:2515}{}魯梭啦，
% cue 2516 6481.590 6482.590
\hypertarget{w02:cue:2516}{}康德\footnote{Immanuel Kant（1724–1804），德意志哲學家。Tremblay 刊本 539 頁謂 GLT1 所述的政治道德一般原則，至少自 Jean-Jacques Rousseau 與 Kant 的著作以來即為規範性政治哲學所預設。\url{https://plato.stanford.edu/entries/kant/}。}啦等等，
% cue 2517 6482.590 6483.550
\hypertarget{w02:cue:2517}{}這個我們就不去談論，
% cue 2518 6483.550 6485.570
\hypertarget{w02:cue:2518}{}證明這個是，
% cue 2519 6485.570 6486.850
\hypertarget{w02:cue:2519}{}一般性是正當的，
% cue 2520 6486.850 6489.710
\hypertarget{w02:cue:2520}{}所有的國家機關原則上，
% cue 2521 6489.710 6490.870
\hypertarget{w02:cue:2521}{}應儘可能的，
% cue 2522 6490.870 6493.960
\hypertarget{w02:cue:2522}{}以合法的方式來行使它的權力，
% cue 2523 6493.960 6495.050
\hypertarget{w02:cue:2523}{}這個是第一個，
% cue 2524 6495.050 6498.010
\hypertarget{w02:cue:2524}{}政治道德就是正當性的原則。

\tc{1:48:18}
% cue 2525 6498.010 6499.390
\hypertarget{w02:cue:2525}{}好，
% cue 2526 6499.390 6500.860
\hypertarget{w02:cue:2526}{}第二個，
% cue 2527 6500.860 6502.810
\hypertarget{w02:cue:2527}{}the judiciary as a matter of fact，
% cue 2528 6502.810 6504.890
\hypertarget{w02:cue:2528}{}exercise political authority in the state，
% cue 2529 6504.890 6510.250
\hypertarget{w02:cue:2529}{}司法部門也是行使政治權力在一個國家裡，
% cue 2530 6510.250 6512.430
\hypertarget{w02:cue:2530}{}這兩點，
% cue 2531 6512.430 6513.650
\hypertarget{w02:cue:2531}{}會跟剛才說，
% cue 2532 6513.650 6516.710
\hypertarget{w02:cue:2532}{}司法是沒有本質上正當性，
% cue 2533 6516.710 6518.930
\hypertarget{w02:cue:2533}{}它的出發點非常不一樣，
% cue 2534 6518.930 6520.430
\hypertarget{w02:cue:2534}{}密尼是機關哦，
% cue 2535 6520.430 6522.620
\hypertarget{w02:cue:2535}{}所有國家機關都有正當性，
% cue 2536 6522.620 6523.620
\hypertarget{w02:cue:2536}{}存念都有正當性的，
% cue 2537 6523.620 6526.030
\hypertarget{w02:cue:2537}{}你看這個命題改了，
% cue 2538 6526.030 6527.030
\hypertarget{w02:cue:2538}{}答案就不一樣了。

\tc{1:48:47}
% cue 2539 6527.030 6528.250
\hypertarget{w02:cue:2539}{}GLT 3，
% cue 2540 6528.250 6533.470
\hypertarget{w02:cue:2540}{}Therefore the judiciary ought to act as far as possible in the legitimate way，
% cue 2541 6533.470 6535.520
\hypertarget{w02:cue:2541}{}所以應該要儘可能的，
% cue 2542 6535.520 6539.280
\hypertarget{w02:cue:2542}{}和合法的方式來行使權力，
% cue 2543 6539.280 6540.860
\hypertarget{w02:cue:2543}{}所以到這裡，
% cue 2544 6540.860 6541.780
\hypertarget{w02:cue:2544}{}GLT 3，
% cue 2545 6541.780 6543.850
\hypertarget{w02:cue:2545}{}我也不覺得它有特殊的比較多，
% cue 2546 6543.850 6544.850
\hypertarget{w02:cue:2546}{}沒關係，
% cue 2547 6544.850 6545.850
\hypertarget{w02:cue:2547}{}它這樣子比較完整，
% cue 2548 6545.850 6546.850
\hypertarget{w02:cue:2548}{}這算一段。

\tc{1:49:06}
% cue 2549 6546.850 6547.840
\hypertarget{w02:cue:2549}{}好，GLT 4，
% cue 2550 6547.840 6553.470
\hypertarget{w02:cue:2550}{}The legitimacy of judicial action is conditioned by the rule of law，
% cue 2551 6553.470 6556.720
\hypertarget{w02:cue:2551}{}這個跟剛才的形式化的字，
% cue 2552 6556.720 6558.130
\hypertarget{w02:cue:2552}{}是一樣的，
% cue 2553 6558.130 6566.890
\hypertarget{w02:cue:2553}{}provided that the legal reason for the action and decisions are legitimate by virtue of some democratic pedigree，
% cue 2554 6566.890 6567.890
\hypertarget{w02:cue:2554}{}這個是剛才的，
% cue 2555 6567.890 6569.880
\hypertarget{w02:cue:2555}{}二、三、四通通放進來了，
% cue 2556 6569.880 6571.520
\hypertarget{w02:cue:2556}{}所以這個時候，
% cue 2557 6571.520 6573.520
\hypertarget{w02:cue:2557}{}要符合這三個條件，
% cue 2558 6573.520 6575.680
\hypertarget{w02:cue:2558}{}這樣它才具有正當性。

\tc{1:49:35}
% cue 2559 6575.680 6578.600
\hypertarget{w02:cue:2559}{}GLT 5，
% cue 2560 6578.600 6579.950
\hypertarget{w02:cue:2560}{}Accordingly，
% cue 2561 6579.950 6591.190
\hypertarget{w02:cue:2561}{}Judicial action and decision that are not justified by law or not based upon the legal reasons that are not legitimate within the ambit，
% cue 2562 6591.190 6592.950
\hypertarget{w02:cue:2562}{}this is not，
% cue 2563 6592.950 6594.250
\hypertarget{w02:cue:2563}{}this is not，
% cue 2564 6594.250 6596.710
\hypertarget{w02:cue:2564}{}ceases are not moral permissible，
% cue 2565 6596.710 6599.340
\hypertarget{w02:cue:2565}{}如果不符合的話，
% cue 2566 6599.340 6601.460
\hypertarget{w02:cue:2566}{}那它們就具有正當性，
% cue 2567 6601.460 6603.280
\hypertarget{w02:cue:2567}{}剛才說，
% cue 2568 6603.280 6605.280
\hypertarget{w02:cue:2568}{}要具備那幾個是具有正當性，
% cue 2569 6605.280 6607.470
\hypertarget{w02:cue:2569}{}如果不符合，
% cue 2570 6607.470 6608.470
\hypertarget{w02:cue:2570}{}那就沒有正當性。

\tc{1:50:08}
% cue 2571 6608.470 6609.720
\hypertarget{w02:cue:2571}{}GLT 6，
% cue 2572 6609.720 6621.740
\hypertarget{w02:cue:2572}{}it follows that judiciary should not recognize as good reasons for action or decision reasons that are not legal，
% cue 2573 6621.740 6624.920
\hypertarget{w02:cue:2573}{}nor legal reasons that are not legitimate，
% cue 2574 6624.920 6627.240
\hypertarget{w02:cue:2574}{}nor legal reasons that are not legitimate。

\tc{1:50:27}
% cue 2575 6627.240 6629.080
\hypertarget{w02:cue:2575}{}好，
% cue 2576 6629.080 6632.530
\hypertarget{w02:cue:2576}{}所以一到三算一段，
% cue 2577 6632.530 6633.910
\hypertarget{w02:cue:2577}{}然後下面呢，
% cue 2578 6633.910 6636.390
\hypertarget{w02:cue:2578}{}大概說到四算一段，
% cue 2579 6636.390 6637.390
\hypertarget{w02:cue:2579}{}五跟六是一段，
% cue 2580 6637.390 6638.390
\hypertarget{w02:cue:2580}{}五跟六是一段，
% cue 2581 6638.390 6639.390
\hypertarget{w02:cue:2581}{}五跟六是反面的命題，
% cue 2582 6639.390 6641.610
\hypertarget{w02:cue:2582}{}那這樣子呢，
% cue 2583 6641.610 6642.580
\hypertarget{w02:cue:2583}{}它就把它的理論呢，
% cue 2584 6642.580 6644.540
\hypertarget{w02:cue:2584}{}大體上都講完了，
% cue 2585 6644.540 6645.630
\hypertarget{w02:cue:2585}{}它要講的講完了。

\section{憲法命題的內容與兩個假設}

\tc{1:50:45}
% cue 2586 6645.630 6647.920
\hypertarget{w02:cue:2586}{}現在是真正比較難的，
% cue 2587 6647.920 6648.920
\hypertarget{w02:cue:2587}{}沒有辦法接應講，
% cue 2588 6648.920 6650.110
\hypertarget{w02:cue:2588}{}是 5444，
% cue 2589 6650.110 6654.970
\hypertarget{w02:cue:2589}{}就看那個 Constitutional Law Thesis，
% cue 2590 6654.970 6657.250
\hypertarget{w02:cue:2590}{}我們只講個頭吧，
% cue 2591 6657.250 6662.380
\hypertarget{w02:cue:2591}{}The Constitutional Law Thesis states，
% cue 2592 6662.380 6663.420
\hypertarget{w02:cue:2592}{}說什麼，
% cue 2593 6663.420 6667.810
\hypertarget{w02:cue:2593}{}The basic rules and principles of the Constitutional Law，
% cue 2594 6667.810 6675.270
\hypertarget{w02:cue:2594}{}providing the criteria of validity and of proper interpretation of legislation，
% cue 2595 6675.270 6680.440
\hypertarget{w02:cue:2595}{}derived from the substance of antecedent norms，
% cue 2596 6680.440 6683.960
\hypertarget{w02:cue:2596}{}respect to the legitimacy of legislation，
% cue 2597 6683.960 6687.630
\hypertarget{w02:cue:2597}{}as understood and enforced by the courts。

\tc{1:51:27}
% cue 2598 6687.630 6690.050
\hypertarget{w02:cue:2598}{}我就把這句話講完，
% cue 2599 6690.050 6692.330
\hypertarget{w02:cue:2599}{}這樣子的話就算是今天結束了，
% cue 2600 6692.330 6696.670
\hypertarget{w02:cue:2600}{}不講，只講到剛才那樣子太抽象了，
% cue 2601 6696.670 6698.400
\hypertarget{w02:cue:2601}{}你沒有辦法理解，
% cue 2602 6698.400 6700.480
\hypertarget{w02:cue:2602}{}一個法律具備，
% cue 2603 6700.480 6705.630
\hypertarget{w02:cue:2603}{}一個司法所準據的法律，
% cue 2604 6705.630 6710.930
\hypertarget{w02:cue:2604}{}它需要具備一定的民主血統，
% cue 2605 6710.930 6712.240
\hypertarget{w02:cue:2605}{}這句話是什麼？
% cue 2606 6712.240 6714.990
\hypertarget{w02:cue:2606}{}聽不懂，
% cue 2607 6714.990 6716.990
\hypertarget{w02:cue:2607}{}它現在這個 Constitutional Law Thesis，
% cue 2608 6716.990 6720.960
\hypertarget{w02:cue:2608}{}就是把這個所謂法律，
% cue 2609 6720.960 6724.300
\hypertarget{w02:cue:2609}{}應具備一定的民主血統給一下解決。

\tc{1:52:04}
% cue 2610 6724.300 6725.880
\hypertarget{w02:cue:2610}{}它說什麼呢，
% cue 2611 6725.880 6728.740
\hypertarget{w02:cue:2611}{}一個憲法裡面，
% cue 2612 6728.740 6731.180
\hypertarget{w02:cue:2612}{}它一定有一些基本原則，
% cue 2613 6731.180 6733.620
\hypertarget{w02:cue:2613}{}或者是基本法則，
% cue 2614 6733.620 6735.780
\hypertarget{w02:cue:2614}{}Rules or Principles，
% cue 2615 6735.780 6737.230
\hypertarget{w02:cue:2615}{}這是 basic，
% cue 2616 6737.230 6739.150
\hypertarget{w02:cue:2616}{}這個東西呢，
% cue 2617 6739.150 6743.010
\hypertarget{w02:cue:2617}{}它會提供一套有效性的標準，
% cue 2618 6743.010 6744.840
\hypertarget{w02:cue:2618}{}有效性的標準，
% cue 2619 6744.840 6751.450
\hypertarget{w02:cue:2619}{}也是提供一套所謂的立法解釋，
% cue 2620 6751.450 6753.450
\hypertarget{w02:cue:2620}{}立法的適當解釋方法，
% cue 2621 6753.450 6756.850
\hypertarget{w02:cue:2621}{}那麼這兩樣東西，
% cue 2622 6756.850 6762.690
\hypertarget{w02:cue:2622}{}是來自於，
% cue 2623 6762.690 6765.770
\hypertarget{w02:cue:2623}{}先前就已經存在，
% cue 2624 6765.770 6767.790
\hypertarget{w02:cue:2624}{}先存的，
% cue 2625 6767.790 6768.790
\hypertarget{w02:cue:2625}{}或者是預制的，
% cue 2626 6768.790 6772.870
\hypertarget{w02:cue:2626}{}這個一套規範。

\tc{1:52:52}
% cue 2627 6772.870 6775.400
\hypertarget{w02:cue:2627}{}那這就是剛才，
% cue 2628 6775.400 6777.340
\hypertarget{w02:cue:2628}{}在這個結論裡面，
% cue 2629 6777.340 6778.360
\hypertarget{w02:cue:2629}{}也有談到，
% cue 2630 6778.360 6782.440
\hypertarget{w02:cue:2630}{}我用一句比較經典的話，
% cue 2631 6782.440 6783.440
\hypertarget{w02:cue:2631}{}把它講一講，
% cue 2632 6783.440 6787.260
\hypertarget{w02:cue:2632}{}有關於立法正當性的，
% cue 2633 6787.260 6788.290
\hypertarget{w02:cue:2633}{}前置規範。

\tc{1:53:08}
% cue 2634 6788.290 6790.080
\hypertarget{w02:cue:2634}{}看一下這句，
% cue 2635 6790.080 6794.340
\hypertarget{w02:cue:2635}{}憲法解釋，
% cue 2636 6794.340 6796.970
\hypertarget{w02:cue:2636}{}因為有助於，
% cue 2637 6796.970 6800.290
\hypertarget{w02:cue:2637}{}使得憲法原則和法則，
% cue 2638 6800.290 6801.340
\hypertarget{w02:cue:2638}{}變得正當，
% cue 2639 6801.340 6805.200
\hypertarget{w02:cue:2639}{}的憲法解釋性論證，
% cue 2640 6805.200 6806.290
\hypertarget{w02:cue:2640}{}而變得正當，
% cue 2641 6806.290 6811.230
\hypertarget{w02:cue:2641}{}憲法解釋不會當然正當，
% cue 2642 6811.230 6812.870
\hypertarget{w02:cue:2642}{}憲法解釋是因為，
% cue 2643 6812.870 6814.830
\hypertarget{w02:cue:2643}{}有下面這個東西，
% cue 2644 6814.830 6815.980
\hypertarget{w02:cue:2644}{}才會變得正當，
% cue 2645 6815.980 6818.350
\hypertarget{w02:cue:2645}{}有一個東西，
% cue 2646 6818.350 6820.610
\hypertarget{w02:cue:2646}{}那個東西是一套解釋性的論證，
% cue 2647 6820.610 6822.190
\hypertarget{w02:cue:2647}{}憲法解釋，
% cue 2648 6822.190 6823.160
\hypertarget{w02:cue:2648}{}我們透過的是，
% cue 2649 6823.160 6825.160
\hypertarget{w02:cue:2649}{}Interpretive Argument，
% cue 2650 6825.160 6826.760
\hypertarget{w02:cue:2650}{}這個 Interpretive Argument，
% cue 2651 6826.760 6827.760
\hypertarget{w02:cue:2651}{}要幹什麼？
% cue 2652 6827.760 6829.140
\hypertarget{w02:cue:2652}{}要使得，
% cue 2653 6829.140 6832.740
\hypertarget{w02:cue:2653}{}憲法原則和法則，
% cue 2654 6832.740 6833.800
\hypertarget{w02:cue:2654}{}變得正當，
% cue 2655 6833.800 6837.650
\hypertarget{w02:cue:2655}{}那這樣子憲法解釋，
% cue 2656 6837.650 6838.680
\hypertarget{w02:cue:2656}{}才會變得正當。

\tc{1:53:58}
% cue 2657 6838.680 6841.670
\hypertarget{w02:cue:2657}{}那如果你要這樣子，
% cue 2658 6841.670 6842.700
\hypertarget{w02:cue:2658}{}那就是說，
% cue 2659 6842.700 6843.790
\hypertarget{w02:cue:2659}{}我們需要，
% cue 2660 6843.790 6845.780
\hypertarget{w02:cue:2660}{}找到一套，
% cue 2661 6845.780 6847.130
\hypertarget{w02:cue:2661}{}解釋，
% cue 2662 6847.130 6849.210
\hypertarget{w02:cue:2662}{}那個立法正當性的，
% cue 2663 6849.210 6850.530
\hypertarget{w02:cue:2663}{}前置性規範，
% cue 2664 6850.530 6852.110
\hypertarget{w02:cue:2664}{}我們有一套標準，
% cue 2665 6852.110 6853.110
\hypertarget{w02:cue:2665}{}你看它叫，
% cue 2666 6853.110 6854.110
\hypertarget{w02:cue:2666}{}Criteria of Validity，
% cue 2667 6854.110 6855.040
\hypertarget{w02:cue:2667}{}如果符合，
% cue 2668 6855.040 6858.640
\hypertarget{w02:cue:2668}{}那它就是 Valid，
% cue 2669 6858.640 6861.340
\hypertarget{w02:cue:2669}{}那它必須要先符合一套，
% cue 2670 6861.340 6863.560
\hypertarget{w02:cue:2670}{}憲法最基本的，
% cue 2671 6863.560 6864.610
\hypertarget{w02:cue:2671}{}價值，
% cue 2672 6864.610 6867.340
\hypertarget{w02:cue:2672}{}一個前置性規範，
% cue 2673 6867.340 6869.160
\hypertarget{w02:cue:2673}{}Values，
% cue 2674 6869.160 6870.950
\hypertarget{w02:cue:2674}{}那個前置性規範，
% cue 2675 6870.950 6871.950
\hypertarget{w02:cue:2675}{}如果符合了，
% cue 2676 6871.950 6874.240
\hypertarget{w02:cue:2676}{}那麼立法就具有正當性，
% cue 2677 6874.240 6876.030
\hypertarget{w02:cue:2677}{}立法具有正當性了，
% cue 2678 6876.030 6878.310
\hypertarget{w02:cue:2678}{}那你去運用這套，
% cue 2679 6878.310 6879.310
\hypertarget{w02:cue:2679}{}具有立法，
% cue 2680 6879.310 6880.300
\hypertarget{w02:cue:2680}{}正當性的立法，
% cue 2681 6880.300 6881.280
\hypertarget{w02:cue:2681}{}來做裁判，
% cue 2682 6881.280 6883.960
\hypertarget{w02:cue:2682}{}然後你在這裡面，
% cue 2683 6883.960 6884.960
\hypertarget{w02:cue:2683}{}論證，
% cue 2684 6884.960 6885.960
\hypertarget{w02:cue:2684}{}你要有很多憲法解釋，
% cue 2685 6885.960 6887.310
\hypertarget{w02:cue:2685}{}你也會因此，
% cue 2686 6887.310 6888.220
\hypertarget{w02:cue:2686}{}就變得正當，
% cue 2687 6888.220 6891.260
\hypertarget{w02:cue:2687}{}講的意思是這樣。

\tc{1:54:51}
% cue 2688 6891.260 6892.380
\hypertarget{w02:cue:2688}{}那它這個理論，
% cue 2689 6892.380 6893.410
\hypertarget{w02:cue:2689}{}一，二，三，
% cue 2690 6893.410 6894.410
\hypertarget{w02:cue:2690}{}是連在一起的，
% cue 2691 6894.410 6896.850
\hypertarget{w02:cue:2691}{}那下面它發展，
% cue 2692 6896.850 6897.860
\hypertarget{w02:cue:2692}{}長的文字呢，
% cue 2693 6897.860 6898.810
\hypertarget{w02:cue:2693}{}是在解釋，
% cue 2694 6898.810 6900.620
\hypertarget{w02:cue:2694}{}什麼情況，
% cue 2695 6900.620 6901.620
\hypertarget{w02:cue:2695}{}我現在講的這一套，
% cue 2696 6901.620 6902.620
\hypertarget{w02:cue:2696}{}所謂的，
% cue 2697 6902.620 6904.770
\hypertarget{w02:cue:2697}{}憲法的命題，
% cue 2698 6904.770 6907.660
\hypertarget{w02:cue:2698}{}不只是適用在，
% cue 2699 6907.660 6908.770
\hypertarget{w02:cue:2699}{}成文憲法，
% cue 2700 6908.770 6910.030
\hypertarget{w02:cue:2700}{}也適用在不成文憲法，
% cue 2701 6910.030 6912.590
\hypertarget{w02:cue:2701}{}那都是有效的命題，
% cue 2702 6912.590 6913.660
\hypertarget{w02:cue:2702}{}那不成文憲法，
% cue 2703 6913.660 6914.560
\hypertarget{w02:cue:2703}{}它花了很長的時間，
% cue 2704 6914.560 6915.650
\hypertarget{w02:cue:2704}{}這個部分我們不容易懂，
% cue 2705 6915.650 6916.680
\hypertarget{w02:cue:2705}{}但是這一段呢，
% cue 2706 6916.680 6918.650
\hypertarget{w02:cue:2706}{}如果你們對於，
% cue 2707 6918.650 6920.710
\hypertarget{w02:cue:2707}{}什麼叫不成文憲法有興趣的話，
% cue 2708 6920.710 6921.930
\hypertarget{w02:cue:2708}{}我覺得這一段，
% cue 2709 6921.930 6924.820
\hypertarget{w02:cue:2709}{}至少我個人不是寫不出來，
% cue 2710 6924.820 6927.170
\hypertarget{w02:cue:2710}{}因為我對不成文憲法的理解太，
% cue 2711 6927.170 6929.290
\hypertarget{w02:cue:2711}{}太，太，太，太，
% cue 2712 6929.290 6930.290
\hypertarget{w02:cue:2712}{}太膚淺，
% cue 2713 6930.290 6931.480
\hypertarget{w02:cue:2713}{}做不出來。

\tc{1:55:31}
% cue 2714 6931.480 6932.480
\hypertarget{w02:cue:2714}{}那至於成文憲法的這一段，
% cue 2715 6932.480 6933.480
\hypertarget{w02:cue:2715}{}大家誤一頁一下，
% cue 2716 6933.480 6934.480
\hypertarget{w02:cue:2716}{}是大家比較熟的，
% cue 2717 6934.480 6935.610
\hypertarget{w02:cue:2717}{}那你，
% cue 2718 6935.610 6936.640
\hypertarget{w02:cue:2718}{}因為它前面講的比較多了，
% cue 2719 6936.640 6937.640
\hypertarget{w02:cue:2719}{}這就，
% cue 2720 6937.640 6938.640
\hypertarget{w02:cue:2720}{}這一段呢，
% cue 2721 6938.640 6939.610
\hypertarget{w02:cue:2721}{}相對講的比較少，
% cue 2722 6939.610 6941.640
\hypertarget{w02:cue:2722}{}但是還是可以去檢視一下，
% cue 2723 6941.640 6943.050
\hypertarget{w02:cue:2723}{}並不是，
% cue 2724 6943.050 6944.050
\hypertarget{w02:cue:2724}{}就說有了成文憲法，
% cue 2725 6944.050 6946.110
\hypertarget{w02:cue:2725}{}然後我們就一定可以依照，
% cue 2726 6946.110 6947.110
\hypertarget{w02:cue:2726}{}成文憲法來宣告，
% cue 2727 6947.110 6948.140
\hypertarget{w02:cue:2727}{}違憲，
% cue 2728 6948.140 6949.140
\hypertarget{w02:cue:2728}{}然後說我這樣就是有效。

\tc{1:55:49}
% cue 2729 6949.140 6950.200
\hypertarget{w02:cue:2729}{}那這個裡面呢，
% cue 2730 6950.200 6952.260
\hypertarget{w02:cue:2730}{}會看到它有許多有趣的人，
% cue 2731 6952.260 6954.210
\hypertarget{w02:cue:2731}{}有會出現了，
% cue 2732 6954.210 6955.280
\hypertarget{w02:cue:2732}{}比方說，
% cue 2733 6955.280 6956.280
\hypertarget{w02:cue:2733}{}555 頁，
% cue 2734 6956.280 6957.400
\hypertarget{w02:cue:2734}{}有供應，
% cue 2735 6957.400 6959.130
\hypertarget{w02:cue:2735}{}他提出來的是，
% cue 2736 6959.130 6961.130
\hypertarget{w02:cue:2736}{}民主代議補強論，
% cue 2737 6961.130 6963.790
\hypertarget{w02:cue:2737}{}民主代議補強論，
% cue 2738 6963.790 6964.730
\hypertarget{w02:cue:2738}{}就是剛才講的，
% cue 2739 6964.730 6967.040
\hypertarget{w02:cue:2739}{}使得立法，
% cue 2740 6967.040 6968.320
\hypertarget{w02:cue:2740}{}解釋立法，
% cue 2741 6968.320 6970.280
\hypertarget{w02:cue:2741}{}是不是正當的，
% cue 2742 6970.280 6971.460
\hypertarget{w02:cue:2742}{}那一套預設回法，
% cue 2743 6971.460 6976.080
\hypertarget{w02:cue:2743}{}antecedent，
% cue 2744 6976.080 6980.150
\hypertarget{w02:cue:2744}{}這個東西會使得立法，
% cue 2745 6980.150 6982.300
\hypertarget{w02:cue:2745}{}決定那個立法是正當與否，
% cue 2746 6982.300 6984.800
\hypertarget{w02:cue:2746}{}然後你用幾例照這樣的，
% cue 2747 6984.800 6985.770
\hypertarget{w02:cue:2747}{}來判斷，
% cue 2748 6985.770 6986.860
\hypertarget{w02:cue:2748}{}你可以使得這些法解釋，
% cue 2749 6986.860 6987.860
\hypertarget{w02:cue:2749}{}變成正當，
% cue 2750 6987.860 6990.880
\hypertarget{w02:cue:2750}{}這好幾件事綁在一起。

\tc{1:56:30}
% cue 2751 6990.880 6991.880
\hypertarget{w02:cue:2751}{}所以 555 頁，
% cue 2752 6991.880 6992.880
\hypertarget{w02:cue:2752}{}有 John Ely\footnote{John Hart Ely（1938–2003），美國憲法學者，1982 年至 1987 年任 Stanford 法學院院長，著有 Democracy and Distrust: A Theory of Judicial Review（Harvard University Press, 1980），主張違憲審查以參與導向、代議補強（representation-reinforcement）為其任務。Tremblay 刊本 555 頁以此說為成文憲法假設下先行規範的一例。\url{https://www.thecrimson.com/article/2003/10/29/constitutional-scholar-former-law-school-professor/}。}，
% cue 2753 6992.880 6994.290
\hypertarget{w02:cue:2753}{}556 頁，
% cue 2754 6994.290 6995.290
\hypertarget{w02:cue:2754}{}有 Michelman，
% cue 2755 6995.290 6996.280
\hypertarget{w02:cue:2755}{}Frank Michelman\footnote{Frank I. Michelman，美國憲法學者。Tremblay 刊本 556 頁引其 Law's Republic（97 Yale Law Journal 1493, 1988）論審議的民主。課綱第 11 週（11 月 23 日）指定其 Brennan and Democracy（86 California Law Review 399, 1998）。}，
% cue 2756 6996.280 6997.630
\hypertarget{w02:cue:2756}{}這是一個老先生，
% cue 2757 6997.630 6998.820
\hypertarget{w02:cue:2757}{}這個人非常有趣，
% cue 2758 6998.820 6999.820
\hypertarget{w02:cue:2758}{}他是猶太人，
% cue 2759 6999.820 7003.100
\hypertarget{w02:cue:2759}{}他寫了很多有趣的文章，
% cue 2760 7003.100 7004.190
\hypertarget{w02:cue:2760}{}我說真的，
% cue 2761 7004.190 7005.640
\hypertarget{w02:cue:2761}{}他的文章很難懂，
% cue 2762 7005.640 7007.770
\hypertarget{w02:cue:2762}{}但是我覺得他的文章非常好。

\tc{1:56:47}
% cue 2763 7007.770 7009.660
\hypertarget{w02:cue:2763}{}再下面呢，
% cue 2764 7009.660 7010.820
\hypertarget{w02:cue:2764}{}556 頁的下一點，
% cue 2765 7010.820 7011.790
\hypertarget{w02:cue:2765}{}有 Sunstein\footnote{Cass R. Sunstein，美國憲法學者，現任 Harvard 大學 Robert Walmsley University Professor。Tremblay 刊本 556 頁引其 Designing Democracy: What Constitutions Do（Oxford University Press, 2001）與 The Partial Constitution（Harvard University Press, 1993），作為審議民主（deliberative democracy）的一說。\url{https://hls.harvard.edu/faculty/cass-r-sunstein/}。}，
% cue 2766 7011.790 7014.230
\hypertarget{w02:cue:2766}{}他講的 deliberative democracy，
% cue 2767 7014.230 7015.550
\hypertarget{w02:cue:2767}{}557 頁，
% cue 2768 7015.550 7016.550
\hypertarget{w02:cue:2768}{}有上學期我們學到的，
% cue 2769 7016.550 7017.500
\hypertarget{w02:cue:2769}{}Ronald Dworkin\footnote{Ronald Dworkin（1931–2013），美國法哲學家，曾任 Oxford 大學法理學講座教授與 New York University 法學院教授。Tremblay 刊本 557 頁引其 Freedom's Law: The Moral Reading of the American Constitution（Harvard University Press, 1996），謂其民主觀以政治制度「treat all members of the community, as individuals, with equal concern and respect」為政治正當性的條件。\url{https://www.law.ox.ac.uk/news/2013-02-15-ronald-dworkin-1931-2013}。}，
% cue 2770 7017.500 7020.830
\hypertarget{w02:cue:2770}{}Equal citizenship。

\tc{1:57:00}
% cue 2771 7020.830 7023.640
\hypertarget{w02:cue:2771}{}這些都是我剛才講的，
% cue 2772 7023.640 7025.050
\hypertarget{w02:cue:2772}{}那個 underlying，
% cue 2773 7025.050 7027.340
\hypertarget{w02:cue:2773}{}antecedent norm，
% cue 2774 7027.340 7029.560
\hypertarget{w02:cue:2774}{}來決定某個法律，
% cue 2775 7029.560 7030.650
\hypertarget{w02:cue:2775}{}是不是具有民主的，
% cue 2776 7030.650 7031.580
\hypertarget{w02:cue:2776}{}血統，
% cue 2777 7031.580 7034.960
\hypertarget{w02:cue:2777}{}從而為正當的法律，
% cue 2778 7034.960 7036.130
\hypertarget{w02:cue:2778}{}然後，
% cue 2779 7036.130 7037.650
\hypertarget{w02:cue:2779}{}司法要去適用那個，
% cue 2780 7037.650 7039.650
\hypertarget{w02:cue:2780}{}具有民主血統的法律，
% cue 2781 7039.650 7041.860
\hypertarget{w02:cue:2781}{}會簡直變得正當，
% cue 2782 7041.860 7044.110
\hypertarget{w02:cue:2782}{}然後另外他也講一點，
% cue 2783 7044.110 7045.170
\hypertarget{w02:cue:2783}{}這個東西也會決定，
% cue 2784 7045.170 7048.200
\hypertarget{w02:cue:2784}{}你可以用的解釋方法，
% cue 2785 7048.200 7049.840
\hypertarget{w02:cue:2785}{}你離開了這裡，
% cue 2786 7049.840 7050.840
\hypertarget{w02:cue:2786}{}你的解釋也會不正當，
% cue 2787 7050.840 7052.890
\hypertarget{w02:cue:2787}{}就是太複雜了。

\section{翻譯作業、停課與後續進度}

\tc{1:57:32}
% cue 2788 7052.890 7054.110
\hypertarget{w02:cue:2788}{}好了，
% cue 2789 7054.110 7055.110
\hypertarget{w02:cue:2789}{}講到這裡，
% cue 2790 7055.110 7057.020
\hypertarget{w02:cue:2790}{}我覺得今天時間已經到了，
% cue 2791 7057.020 7059.170
\hypertarget{w02:cue:2791}{}我並沒有把這個工作做完，
% cue 2792 7059.170 7061.260
\hypertarget{w02:cue:2792}{}那這個工作確實很難做，
% cue 2793 7061.260 7063.210
\hypertarget{w02:cue:2793}{}我想徵求一個，
% cue 2794 7063.210 7065.220
\hypertarget{w02:cue:2794}{}你們有人在座人不考試的局所，
% cue 2795 7065.220 7067.720
\hypertarget{w02:cue:2795}{}我看看，
% cue 2796 7067.720 7068.720
\hypertarget{w02:cue:2796}{}你們可能兩個人，
% cue 2797 7068.720 7070.690
\hypertarget{w02:cue:2797}{}你們可以幫一個，
% cue 2798 7070.690 7071.840
\hypertarget{w02:cue:2798}{}你們用 AI，
% cue 2799 7071.840 7073.320
\hypertarget{w02:cue:2799}{}去把這段文字呢，
% cue 2800 7073.320 7074.320
\hypertarget{w02:cue:2800}{}通通變成中文，
% cue 2801 7074.320 7077.720
\hypertarget{w02:cue:2801}{}做完了。

\tc{1:57:57}
% cue 2802 7077.720 7078.720
\hypertarget{w02:cue:2802}{}好好好，
% cue 2803 7078.720 7079.750
\hypertarget{w02:cue:2803}{}非常好，
% cue 2804 7079.750 7080.750
\hypertarget{w02:cue:2804}{}你做完了，
% cue 2805 7080.750 7081.740
\hypertarget{w02:cue:2805}{}然後，
% cue 2806 7081.740 7082.770
\hypertarget{w02:cue:2806}{}你是你要做，
% cue 2807 7082.770 7083.770
\hypertarget{w02:cue:2807}{}如果你能做更好，
% cue 2808 7083.770 7085.240
\hypertarget{w02:cue:2808}{}那你試試看，
% cue 2809 7085.240 7086.850
\hypertarget{w02:cue:2809}{}就是把它，
% cue 2810 7086.850 7090.380
\hypertarget{w02:cue:2810}{}像我剛才講的那樣，
% cue 2811 7090.380 7091.380
\hypertarget{w02:cue:2811}{}尤其是最後那個結論，
% cue 2812 7091.380 7092.890
\hypertarget{w02:cue:2812}{}那個結論非常難講，
% cue 2813 7092.890 7095.970
\hypertarget{w02:cue:2813}{}我不知道我剛才那樣子，
% cue 2814 7095.970 7097.800
\hypertarget{w02:cue:2814}{}因為你做完了，
% cue 2815 7097.800 7100.310
\hypertarget{w02:cue:2815}{}再聽我剛才那樣的解說，
% cue 2816 7100.310 7103.530
\hypertarget{w02:cue:2816}{}你認為對你所做出來的文字的翻譯，
% cue 2817 7103.530 7104.680
\hypertarget{w02:cue:2817}{}的理解，
% cue 2818 7104.680 7105.780
\hypertarget{w02:cue:2818}{}會不會有幫助，
% cue 2819 7105.780 7109.470
\hypertarget{w02:cue:2819}{}也會讓你更看不懂那些有字天書，
% cue 2820 7109.470 7112.010
\hypertarget{w02:cue:2820}{}那個做出來像個有字天書，
% cue 2821 7112.010 7113.940
\hypertarget{w02:cue:2821}{}尤其用中文的表述，
% cue 2822 7113.940 7114.940
\hypertarget{w02:cue:2822}{}你真的沒辦法，
% cue 2823 7114.940 7115.940
\hypertarget{w02:cue:2823}{}不知道，
% cue 2824 7115.940 7116.940
\hypertarget{w02:cue:2824}{}還講，
% cue 2825 7116.940 7118.380
\hypertarget{w02:cue:2825}{}英文都非常難懂。

\tc{1:58:38}
% cue 2826 7118.380 7119.380
\hypertarget{w02:cue:2826}{}所以我剛一直跟你講說，
% cue 2827 7119.380 7121.440
\hypertarget{w02:cue:2827}{}這個等於這個等於這個等於這個等於這個，
% cue 2828 7121.440 7122.440
\hypertarget{w02:cue:2828}{}那有些經典的語句，
% cue 2829 7122.440 7123.440
\hypertarget{w02:cue:2829}{}它一直在帶外，
% cue 2830 7123.440 7124.440
\hypertarget{w02:cue:2830}{}那個經典語句，
% cue 2831 7124.440 7125.440
\hypertarget{w02:cue:2831}{}你要把它套裝，
% cue 2832 7125.440 7126.440
\hypertarget{w02:cue:2832}{}套裝怎麼樣，
% cue 2833 7126.440 7127.700
\hypertarget{w02:cue:2833}{}套子套裝 A，
% cue 2834 7127.700 7128.700
\hypertarget{w02:cue:2834}{}套餐 A，
% cue 2835 7128.700 7129.700
\hypertarget{w02:cue:2835}{}套餐 B，
% cue 2836 7129.700 7130.690
\hypertarget{w02:cue:2836}{}套餐 C，
% cue 2837 7130.690 7131.790
\hypertarget{w02:cue:2837}{}然後看他們自己的關係，
% cue 2838 7131.790 7133.430
\hypertarget{w02:cue:2838}{}那你試試看，
% cue 2839 7133.430 7134.430
\hypertarget{w02:cue:2839}{}如果你能做，
% cue 2840 7134.430 7135.620
\hypertarget{w02:cue:2840}{}你願不願意？
% cue 2841 7135.620 7136.620
\hypertarget{w02:cue:2841}{}這實在是，
% cue 2842 7136.620 7137.620
\hypertarget{w02:cue:2842}{}你的，
% cue 2843 7137.620 7138.580
\hypertarget{w02:cue:2843}{}對大家的貢獻，
% cue 2844 7138.580 7139.700
\hypertarget{w02:cue:2844}{}你如果可以的話，
% cue 2845 7139.700 7140.700
\hypertarget{w02:cue:2845}{}你試試看。

\tc{1:59:00}
% cue 2846 7140.700 7142.700
\hypertarget{w02:cue:2846}{}那下個禮拜，
% cue 2847 7142.700 7144.020
\hypertarget{w02:cue:2847}{}等你做出來，
% cue 2848 7144.020 7145.020
\hypertarget{w02:cue:2848}{}我等，
% cue 2849 7145.020 7146.020
\hypertarget{w02:cue:2849}{}我給你一點時間，
% cue 2850 7146.020 7147.020
\hypertarget{w02:cue:2850}{}如果一個禮拜不夠，
% cue 2851 7147.020 7148.020
\hypertarget{w02:cue:2851}{}不再一個禮拜，
% cue 2852 7148.020 7149.020
\hypertarget{w02:cue:2852}{}那麼下禮拜不上課\footnote{2026 年 9 月 28 日為孔子誕辰紀念日（教師節）。紀念日及節日實施條例於 2025 年 5 月 28 日公布，教師節自此列為國定假日；課綱第 3 週（9 月 28 日）記為教師節放假。\url{https://www.mol.gov.tw/1607/28162/28166/28218/28226/81488/}。}，
% cue 2853 7149.020 7150.100
\hypertarget{w02:cue:2853}{}對不對？
% cue 2854 7150.100 7151.180
\hypertarget{w02:cue:2854}{}那我們課綱，
% cue 2855 7151.180 7152.180
\hypertarget{w02:cue:2855}{}我決定不改了，
% cue 2856 7152.180 7154.180
\hypertarget{w02:cue:2856}{}因為下個禮拜既然不上課，
% cue 2857 7154.180 7156.180
\hypertarget{w02:cue:2857}{}那再下一個禮拜，
% cue 2858 7156.180 7157.180
\hypertarget{w02:cue:2858}{}有 80 頁，
% cue 2859 7157.180 7158.310
\hypertarget{w02:cue:2859}{}那老師也認了，
% cue 2860 7158.310 7160.720
\hypertarget{w02:cue:2860}{}就是下個禮拜不上課的時候，
% cue 2861 7160.720 7161.720
\hypertarget{w02:cue:2861}{}我準備 40 頁，
% cue 2862 7161.720 7162.720
\hypertarget{w02:cue:2862}{}再下禮拜再要上課，
% cue 2863 7162.720 7163.680
\hypertarget{w02:cue:2863}{}再準備 40 頁，
% cue 2864 7163.680 7164.870
\hypertarget{w02:cue:2864}{}所以課綱就不改了，
% cue 2865 7164.870 7165.960
\hypertarget{w02:cue:2865}{}就 80 頁，
% cue 2866 7165.960 7166.960
\hypertarget{w02:cue:2866}{}我們就讀一讀，
% cue 2867 7166.960 7168.960
\hypertarget{w02:cue:2867}{}那個 Hans Kelsen 跟 Carl Schmitt\footnote{Carl Schmitt（1888–1985），德國公法學家與政治理論家，1931 年出版 Der Hüter der Verfassung（憲法的守護者），反對由法院審查法律的合憲性。課綱第 5 週（10 月 12 日）指定該書與 Kelsen 1929 年論文的 Lars Vinx 英譯（2015）。\cite{schmitt2015}。}，
% cue 2868 7172.530 7173.910
\hypertarget{w02:cue:2868}{}我沒辦法，
% cue 2869 7173.910 7174.910
\hypertarget{w02:cue:2869}{}因為課綱已經排了，
% cue 2870 7174.910 7176.100
\hypertarget{w02:cue:2870}{}我換一篇文章，
% cue 2871 7176.100 7177.100
\hypertarget{w02:cue:2871}{}當然比較短，
% cue 2872 7177.100 7178.100
\hypertarget{w02:cue:2872}{}比較容易，
% cue 2873 7178.100 7179.060
\hypertarget{w02:cue:2873}{}但是我就試試看，
% cue 2874 7179.060 7181.250
\hypertarget{w02:cue:2874}{}反正那個我也沒有認真讀，
% cue 2875 7181.250 7183.080
\hypertarget{w02:cue:2875}{}我也讀一讀，
% cue 2876 7183.080 7184.050
\hypertarget{w02:cue:2876}{}好不好？

\tc{1:59:44}
% cue 2877 7184.050 7185.050
\hypertarget{w02:cue:2877}{}那我們今天上到這邊，
% cue 2878 7185.050 7186.010
\hypertarget{w02:cue:2878}{}下課，
% cue 2879 7186.010 7187.010
\hypertarget{w02:cue:2879}{}那就看看呢，
% cue 2880 7187.010 7188.010
\hypertarget{w02:cue:2880}{}語方能幫忙，
% cue 2881 7188.010 7189.010
\hypertarget{w02:cue:2881}{}你如果能做一做，
% cue 2882 7189.010 7190.490
\hypertarget{w02:cue:2882}{}你就做，
% cue 2883 7190.490 7191.490
\hypertarget{w02:cue:2883}{}那如果不能做呢，
% cue 2884 7191.490 7192.490
\hypertarget{w02:cue:2884}{}你至少，
% cue 2885 7192.490 7194.490
\hypertarget{w02:cue:2885}{}你可以把它中文，
% cue 2886 7194.490 7196.740
\hypertarget{w02:cue:2886}{}那個 AI 通過翻譯的，
% cue 2887 7196.740 7197.660
\hypertarget{w02:cue:2887}{}至少翻出來嘛，
% cue 2888 7197.660 7198.900
\hypertarget{w02:cue:2888}{}大家先乒乓簡，
% cue 2889 7198.900 7200.000
\hypertarget{w02:cue:2889}{}看看，
% cue 2890 7200.000 7202.560
\hypertarget{w02:cue:2890}{}聞起來是香的還是臭的，
% cue 2891 7202.560 7203.560
\hypertarget{w02:cue:2891}{}好，
% cue 2892 7203.560 7204.560
\hypertarget{w02:cue:2892}{}謝謝。

\tc{2:00:04}
% cue 2893 7204.560 7205.970
\hypertarget{w02:cue:2893}{}好，
% cue 2894 7205.970 7206.970
\hypertarget{w02:cue:2894}{}那你該怎麼辦，
% cue 2895 7206.970 7207.970
\hypertarget{w02:cue:2895}{}你怎麼辦，
% cue 2896 7207.970 7208.970
\hypertarget{w02:cue:2896}{}你有事情要做。

\tc{2:00:08}
\speaker{學生}
% cue 2897 7208.970 7209.970
\hypertarget{w02:cue:2897}{}大家各位先等一下，
% cue 2898 7209.970 7210.970
\hypertarget{w02:cue:2898}{}我們要做一些事情，
% cue 2899 7210.970 7212.450
\hypertarget{w02:cue:2899}{}一件事情是，
% cue 2900 7212.450 7215.450
\hypertarget{w02:cue:2900}{}各位想要把我們這學期，
% cue 2901 7215.450 7216.450
\hypertarget{w02:cue:2901}{}的文本，
% cue 2902 7216.450 7217.450
\hypertarget{w02:cue:2902}{}列印出來，
% cue 2903 7217.450 7218.450
\hypertarget{w02:cue:2903}{}做成一策碼，
% cue 2904 7218.450 7220.700
\hypertarget{w02:cue:2904}{}A4 的時候，
% cue 2905 7220.700 7221.700
\hypertarget{w02:cue:2905}{}他像大大一策碼，
% cue 2906 7221.700 7222.700
\hypertarget{w02:cue:2906}{}被裝進程，
% cue 2907 7222.700 7223.650
\hypertarget{w02:cue:2907}{}他就那麼好，
% cue 2908 7223.650 7224.660
\hypertarget{w02:cue:2908}{}他就看不到，
% cue 2909 7224.660 7225.660
\hypertarget{w02:cue:2909}{}管理應該起的，
% cue 2910 7225.660 7226.660
\hypertarget{w02:cue:2910}{}我們上去的學期，
% cue 2911 7226.660 7228.780
\hypertarget{w02:cue:2911}{}你各位會有興趣嗎？
% cue 2912 7228.780 7229.690
\hypertarget{w02:cue:2912}{}想要這樣做，
% cue 2913 7229.690 7233.640
\hypertarget{w02:cue:2913}{}大家會想做一個評價，
% cue 2914 7233.640 7236.660
\hypertarget{w02:cue:2914}{}因為這個要交錢，
% cue 2915 7236.660 7237.660
\hypertarget{w02:cue:2915}{}對，
% cue 2916 7237.660 7239.660
\hypertarget{w02:cue:2916}{}你們就按照你們付的錢，
% cue 2917 7239.660 7240.660
\hypertarget{w02:cue:2917}{}就去，
% cue 2918 7240.660 7241.660
\hypertarget{w02:cue:2918}{}多少錢，
% cue 2919 7241.660 7248.190
\hypertarget{w02:cue:2919}{}我是比較傾向於，
% cue 2920 7248.190 7249.190
\hypertarget{w02:cue:2920}{}你們要做，
% cue 2921 7249.190 7250.250
\hypertarget{w02:cue:2921}{}因為你們做的，
% cue 2922 7250.250 7252.040
\hypertarget{w02:cue:2922}{}真的很難讀。

\tc{2:00:52}
\speaker{學生}
% cue 2923 7252.040 7256.630
\hypertarget{w02:cue:2923}{}像我們到時候去列印，
% cue 2924 7256.630 7257.630
\hypertarget{w02:cue:2924}{}然後到時候再給各位，
% cue 2925 7257.630 7258.630
\hypertarget{w02:cue:2925}{}然後我們到時候，
% cue 2926 7258.630 7259.630
\hypertarget{w02:cue:2926}{}費用的部分，
% cue 2927 7259.630 7260.630
\hypertarget{w02:cue:2927}{}我會再跟各位一起。

\tc{2:01:00}
\speaker{老師}
% cue 2928 7260.630 7261.630
\hypertarget{w02:cue:2928}{}你先預收一筆錢，
% cue 2929 7261.630 7262.720
\hypertarget{w02:cue:2929}{}不然到時候你收不到錢，
% cue 2930 7262.720 7263.710
\hypertarget{w02:cue:2930}{}你會哭。

\tc{2:01:03}
\speaker{學生}
% cue 2931 7263.710 7265.390
\hypertarget{w02:cue:2931}{}好，
% cue 2932 7265.390 7266.390
\hypertarget{w02:cue:2932}{}那我先，
% cue 2933 7266.390 7267.390
\hypertarget{w02:cue:2933}{}到時候我會先跟這個。

\tc{2:01:07}
\speaker{老師}
% cue 2934 7267.390 7268.390
\hypertarget{w02:cue:2934}{}這種事情就證明，
% cue 2935 7268.390 7269.380
\hypertarget{w02:cue:2935}{}你不是會做生意的。

\tc{2:01:09}
\speaker{學生}
% cue 2936 7269.380 7270.380
\hypertarget{w02:cue:2936}{}對啊，
% cue 2937 7270.380 7270.480
\hypertarget{w02:cue:2937}{}老師，
% cue 2938 7270.380 7271.380
\hypertarget{w02:cue:2938}{}我會做生意的，
% cue 2939 7271.380 7272.380
\hypertarget{w02:cue:2939}{}哈哈哈，
% cue 2940 7272.380 7273.380
\hypertarget{w02:cue:2940}{}我會做生意，
% cue 2941 7273.380 7274.380
\hypertarget{w02:cue:2941}{}是很正常的。

\tc{2:01:14}
\speaker{老師}
% cue 2942 7274.380 7275.380
\hypertarget{w02:cue:2942}{}然後，
% cue 2943 7275.380 7276.380
\hypertarget{w02:cue:2943}{}老師是不想害你，
% cue 2944 7276.380 7277.600
\hypertarget{w02:cue:2944}{}你做了班長，
% cue 2945 7277.600 7278.600
\hypertarget{w02:cue:2945}{}就說，
% cue 2946 7278.600 7278.700
\hypertarget{w02:cue:2946}{}哇，
% cue 2947 7278.600 7279.600
\hypertarget{w02:cue:2947}{}老師你害死我，
% cue 2948 7279.600 7280.600
\hypertarget{w02:cue:2948}{}先提醒我，
% cue 2949 7280.600 7281.600
\hypertarget{w02:cue:2949}{}我要先提醒你，
% cue 2950 7281.600 7282.600
\hypertarget{w02:cue:2950}{}你可以先收一筆錢，
% cue 2951 7282.600 7283.600
\hypertarget{w02:cue:2951}{}然後說，
% cue 2952 7283.600 7284.600
\hypertarget{w02:cue:2952}{}先收個五百塊，
% cue 2953 7284.600 7286.600
\hypertarget{w02:cue:2953}{}五百塊應該差不多了。

\tc{2:01:26}
\speaker{學生}
% cue 2954 7286.600 7287.890
\hypertarget{w02:cue:2954}{}對啊，
% cue 2955 7287.890 7289.100
\hypertarget{w02:cue:2955}{}再給我吧，
% cue 2956 7289.100 7290.100
\hypertarget{w02:cue:2956}{}正前一臉，
% cue 2957 7290.100 7291.100
\hypertarget{w02:cue:2957}{}還一臉。

\tc{2:01:31}
\speaker{老師}
% cue 2958 7291.100 7291.200
\hypertarget{w02:cue:2958}{}沒關係，
% cue 2959 7291.100 7292.100
\hypertarget{w02:cue:2959}{}我們再見吧，
% cue 2960 7292.100 7293.010
\hypertarget{w02:cue:2960}{}拜拜，
% cue 2961 7293.010 7294.160
\hypertarget{w02:cue:2961}{}拜拜。

\tc{2:01:34}
\speaker{學生}
% cue 2962 7294.160 7294.260
\hypertarget{w02:cue:2962}{}呃，
% cue 2963 7294.160 7295.160
\hypertarget{w02:cue:2963}{}然後還有我這邊，
% cue 2964 7295.160 7296.160
\hypertarget{w02:cue:2964}{}在座的同學，
% cue 2965 7296.160 7297.160
\hypertarget{w02:cue:2965}{}有人沒加入到群組的嗎？
% cue 2966 7297.160 7298.160
\hypertarget{w02:cue:2966}{}不知道我加入了，
% cue 2967 7298.160 7299.070
\hypertarget{w02:cue:2967}{}對吧？
% cue 2968 7299.070 7301.150
\hypertarget{w02:cue:2968}{}那我們接下來，
